% Generated mechanically from the frozen pione.tex target.
% Do not edit this generated file; edit the bound locale source instead.

% OK, start here.
%

\setcounter{chapter}{57}
\chapter{概形的基本群}



\phantomsection
\label{pione-section-phantom}


\section{引言}
\label{pione-section-introduction}

\noindent
本章讨论概形的 Grothendieck 基本群及其应用。基础性参考文献是 \cite{SGA1}。
一篇很好的导论是 \cite{Lenstra}。
其他参考文献包括 \cite{Murre-lectures} 和 \cite{Grothendieck-Murre}。










\section{点上的 \'etale 概形}
\label{pione-section-schemes-etale-point}

\noindent
本节描述域的谱上的 \'etale 概形。在陈述结果之前，我们先对拓扑群 $G$
引入 $G$-集范畴。

\begin{definition}
\label{pione-definition-G-set-continuous}
设 $G$ 为拓扑群。一个{\it $G$-集}，有时也称为{\it 离散 $G$-集}，
是一个集合 $X$ 连同一个左作用 $a : G \times X \to X$；要求 $a$ 连续，
其中给 $X$ 赋予离散拓扑，给 $G \times X$ 赋予乘积拓扑。
一个{\it $G$-集的态射} $f : X \to Y$ 就是任意的 $G$-等变映射，
其源为 $X$、靶为 $Y$。$G$-集范畴记作 {\it $G\textit{-Sets}$}。
\end{definition}

\noindent
条件 $a : G \times X \to X$ 连续，无非表示任意 $x \in X$ 的稳定子群在 $G$
中是开的。如果 $G$ 是抽象群 $G$（即它是群但不是拓扑群），只要给它赋予
那么上述定义与先前的定义一致（例如见 Sites, Example
\ref{sites-example-site-on-group}），前提是给 $G$ 赋予离散拓扑。

\medskip\noindent
回忆，如果 $L/K$ 是无限 Galois 扩张，则 Galois 群
$G = \text{Gal}(L/K)$ 带有典范拓扑；见 Fields, Section
\ref{fields-section-infinite-galois}。

\begin{lemma}
\label{pione-lemma-sheaves-point}
设 $K$ 为域。设 $K^{sep}$ 为 $K$ 的一个可分闭包。考虑 profinite 群
$G = \text{Gal}(K^{sep}/K)$。函子
$$
\begin{matrix}
\text{schemes \'etale over }K &
\longrightarrow &
G\textit{-Sets} \\
X/K & \longmapsto &
\Mor_{\Spec(K)}(\Spec(K^{sep}), X)
\end{matrix}
$$
是范畴等价。
\end{lemma}

\begin{proof}
概形 $X$ 定义在 $K$ 上；它在 $K$ 上为 \'etale，当且仅当
$X \cong \coprod_{i\in I} \Spec(K_i)$，其中每个 $K_i$ 都是 $K$ 的有限可分扩张
（Morphisms, Lemma \ref{morphisms-lemma-etale-over-field}）。该引理的函子将
$X$ 送到 $G$-集
$$
\coprod\nolimits_i \Hom_K(K_i, K^{sep})
$$
及其自然左 $G$-作用。由 $G$ 上拓扑的定义，每个元素都有开稳定子群。反之，
任意 $G$-集 $S$ 都是其轨道的不交并。设 $S = \coprod S_i$。选取
$s_i \in S_i$，并用 $G_i \subset G$ 表示它的开稳定子群。由 Galois 理论
（Fields, Theorem \ref{fields-theorem-inifinite-galois-theory}），域
$(K^{sep})^{G_i}$ 是 $K$ 的有限可分扩张，因而概形
$$
\coprod\nolimits_i \Spec((K^{sep})^{G_i})
$$
在 $K$ 上为 \'etale。这给出了该引理函子的逆。略去一些细节。
\end{proof}

\begin{remark}
\label{pione-remark-covering-surjective}
在引理 \ref{pione-lemma-sheaves-point} 的对应下，小 \'etale 位点
$\Spec(K)_\etale$，即 $K$ 的小 \'etale 位点，其中的覆盖对应于
$G\textit{-Sets}$ 中的满射映射族。
\end{remark}









\section{Galois 范畴}
\label{pione-section-galois}

\noindent
本节讨论读者可在 \cite[Expos\'e V, Sections 4, 5, and 6]{SGA1}
中找到的部分材料。

\medskip\noindent
设 $F : \mathcal{C} \to \textit{Sets}$ 为函子。回忆，依照我们的约定，
范畴的对象构成集合，并且任意一对对象之间的态射也构成集合。有典范单射
\begin{equation}
\label{pione-equation-embedding-product}
\text{Aut}(F)
\longrightarrow
\prod\nolimits_{X \in \Ob(\mathcal{C})} \text{Aut}(F(X))
\end{equation}
对集合 $E$，给 $\text{Aut}(E)$ 赋予紧开拓扑；见 Topology, Example
\ref{topology-example-automorphisms-of-a-set}。当然，当 $E$ 有限时，这就是
离散拓扑；本节所关心的正是这种情形\footnote{讨论 pro-\'etale 基本群时，
一般情形也将是我们所关心的。}。给 $\text{Aut}(F)$ 赋予由
(\ref{pione-equation-embedding-product}) 右端的乘积拓扑诱导的拓扑。特别地，作用映射
$$
\text{Aut}(F) \times F(X) \longrightarrow F(X)
$$
在给 $F(X)$ 赋予离散拓扑时是连续的，因为作用映射
$\text{Aut}(E) \times E \to E$ 对任意集合 $E$ 都具有这一性质。
$\text{Aut}(F)$ 上该拓扑的泛性质如下：设 $G$ 为拓扑群，而
$G \to \text{Aut}(F)$ 为群同态，并且诱导作用 $G \times F(X) \to F(X)$
对所有 $X \in \Ob(\mathcal{C})$ 都连续，其中给 $F(X)$ 赋予离散拓扑。
那么 $G \to \text{Aut}(F)$ 连续。

\medskip\noindent
下述引理说明，取值于有限集范畴的函子之自同构群自动成为 profinite 群。

\begin{lemma}
\label{pione-lemma-aut-inverse-limit}
设 $\mathcal{C}$ 为范畴，$F : \mathcal{C} \to \textit{Sets}$ 为函子。映射
(\ref{pione-equation-embedding-product}) 将 $\text{Aut}(F)$ 识别为
$\prod_{X \in \Ob(\mathcal{C})} \text{Aut}(F(X))$ 的闭子群。特别地，如果
$F(X)$ 对所有 $X$ 都有限，那么 $\text{Aut}(F)$ 是 profinite 群。
\end{lemma}

\begin{proof}
设 $\xi = (\gamma_X) \in \prod \text{Aut}(F(X))$ 是不属于
$\text{Aut}(F)$ 的元素。于是存在态射 $f : X \to X'$（属于 $\mathcal{C}$）
及元素 $x \in F(X)$，使得
$F(f)(\gamma_X(x)) \not = \gamma_{X'}(F(f)(x))$.
考虑开邻域
$U = \{\gamma \in \text{Aut}(F(X)) \mid \gamma(x) = \gamma_X(x)\}$，
它是 $\gamma_X$ 的邻域；以及开邻域
$U' = \{\gamma' \in \text{Aut}(F(X')) \mid \gamma'(F(f)(x)) =
\gamma_{X'}(F(f)(x))\}$.
于是
$U \times U' \times \prod_{X'' \not = X, X'} \text{Aut}(F(X''))$
是 $\xi$ 的开邻域，且不与 $\text{Aut}(F)$ 相交。最后一个断言来自如下事实：
$\prod \text{Aut}(F(X))$ 在每个 $F(X)$ 都有限时是 profinite 空间。
\end{proof}

\begin{example}
\label{pione-example-galois-category-G-sets}
设 $G$ 为拓扑群。一个重要例子是遗忘函子
\begin{equation}
\label{pione-equation-forgetful}
\textit{Finite-}G\textit{-Sets} \longrightarrow \textit{Sets}
\end{equation}
其中 $\textit{Finite-}G\textit{-Sets}$ 是 $G\textit{-Sets}$ 的全子范畴，
其对象是有限 $G$-集。$G\textit{-Sets}$（即 $G$-集范畴）定义于 Definition
\ref{pione-definition-G-set-continuous}。
\end{example}

\noindent
设 $G$ 为拓扑群。$G$ 的{\it profinite 完备化}是 profinite 群
$$
G^\wedge =
\lim_{U \subset G\text{ open, normal, finite index}} G/U
$$
及其 profinite 拓扑。注意，该极限是余滤过的，因为有限多个开、正规且
具有有限指数的子群之交仍具有这些性质。profinite 完备化的泛性质是：任意连续映射
$G \to H$，其中 $H$ 为 profinite 群，都典范地分解为
$G \to G^\wedge \to H$。

\begin{lemma}
\label{pione-lemma-single-out-profinite}
设 $G$ 为拓扑群。给函子 (\ref{pione-equation-forgetful}) 的自同构群赋予引理
\ref{pione-lemma-aut-inverse-limit} 所给的 profinite 拓扑，则该群是 $G$ 的
profinite 完备化。
\end{lemma}

\begin{proof}
用 $F_G$ 表示函子 (\ref{pione-equation-forgetful})。任意态射 $X \to Y$（属于
$\textit{Finite-}G\textit{-Sets}$）都与 $G$ 的作用可交换。
因此任意 $g \in G$ 都定义 $F_G$ 的一个自同构，从而得到典范群同态
$G \to \text{Aut}(F_G)$。注意，任意有限 $G$-集 $X$ 都是若干 $G$-集的
有限不交并，这些集合形如 $G/H_i$ 并带有典范 $G$-作用，其中 $H_i \subset G$
是有限指数开子群。于是 $U_i = \bigcap gH_ig^{-1}$ 是开、正规且具有有限指数的。
此外，$U_i$ 在 $G/H_i$ 上的作用平凡，故 $U = \bigcap U_i$ 在 $F_G(X)$ 上的
作用平凡。因此作用 $G \times F_G(X) \to F_G(X)$ 连续。由
$\text{Aut}(F_G)$ 上拓扑的泛性质，映射 $G \to \text{Aut}(F_G)$ 连续。
由引理 \ref{pione-lemma-aut-inverse-limit} 和 profinite 完备化的泛性质，诱导出连续群同态
$$
G^\wedge \longrightarrow \text{Aut}(F_G)
$$
此外，由于 $G/U$ 在 $G/U$ 上的作用忠实，该映射为单射。如果其像稠密，
则该映射为满射，因而由 Topology, Lemma \ref{topology-lemma-bijective-map}
可知它是同胚。

\medskip\noindent
设 $\gamma \in \text{Aut}(F_G)$ 且 $X \in \Ob(\mathcal{C})$。我们将证明存在
$g \in G$，使 $\gamma$ 与 $g$ 在 $F_G(X)$ 上诱导相同的作用。这将完成证明。
如前所述，$X$ 是若干 $G/H_i$ 的有限不交并。取上述 $U_i$ 和 $U$，则有限
$G$-集 $Y = G/U$ 满射到 $G/H_i$（对所有 $i$），故只需找到 $g \in G$，使 $\gamma$
与 $g$ 在 $F_G(G/U) = G/U$ 上诱导相同作用。设 $e \in G$ 为单位元，并设
$\gamma(eU) = g_0U$，其中 $g_0 \in G$。对任意 $g_1 \in G$，态射
$$
R_{g_1} : G/U \longrightarrow G/U,\quad gU \longmapsto gg_1U
$$
属于 $\textit{Finite-}G\textit{-Sets}$，并且与 $\gamma$ 的作用可交换。因此
$$
\gamma(g_1U) = \gamma(R_{g_1}(eU)) = R_{g_1}(\gamma(eU)) =
R_{g_1}(g_0U) = g_0g_1U
$$
所以 $g = g_0$ 符合要求。
\end{proof}

\noindent
回忆，正合函子是与所有有限极限和有限余极限可交换的函子。特别地，这种函子
与等化子、余等化子、纤维积、推出等构造可交换。

\begin{lemma}
\label{pione-lemma-second-fundamental-functor}
设 $G$ 为拓扑群。设
$F : \textit{Finite-}G\textit{-Sets} \to \textit{Sets}$
为正合函子，并且 $F(X)$ 对所有 $X$ 都有限。那么 $F$ 同构于函子
(\ref{pione-equation-forgetful})。
\end{lemma}

\begin{proof}
设 $X$ 为 $\textit{Finite-}G\textit{-Sets}$ 的非空对象。图
$$
\xymatrix{
X \ar[r] \ar[d] & \{*\} \ar[d] \\
\{*\} \ar[r] & \{*\}
}
$$
为余笛卡儿图。因此 $F(X)$ 非空。设 $U \subset G$ 为有限指数开正规子群。
注意
$$
G/U \times G/U = \coprod\nolimits_{gU \in G/U} G/U
$$
其中与 $gU$ 对应的直和项，对应于左端 $(eU, gU)$ 的轨道。于是
$$
F(G/U) \times F(G/U) = F(G/U \times G/U) = \coprod\nolimits_{gU \in G/U} F(G/U)
$$
故 $|F(G/U)| = |G/U|$，因为 $F(G/U)$ 非空。因此
$$
\lim_{U \subset G\text{ open, normal, finite idex}} F(G/U)
$$
非空（Categories, Lemma \ref{categories-lemma-nonempty-limit}）。选取该极限中的
元素 $\gamma = (\gamma_U)$。用 $F_G$ 表示函子 (\ref{pione-equation-forgetful})。
可将 $F_G$ 识别为函子
$$
X \longmapsto \colim_U \Mor(G/U, X)
$$
其中 $f : G/U \to X$ 对应于 $f(eU) \in X = F_G(X)$（略去细节）。因此元素
$\gamma$ 决定一个良定义的映射
$$
t : F_G \longrightarrow F
$$
具体而言，给定 $x \in X$，选取 $U$ 和 $f : G/U \to X$，使它把 $eU$
送到 $x$，再置 $t_X(x) = F(f)(\gamma_U)$。我们将证明，$t$ 诱导双射
$t_{G/U} : F_G(G/U) \to F(G/U)$（对任意 $U$）。由此容易推出 $t$ 是同构
（略去细节）。由于 $|F_G(G/U)| = |F(G/U)|$，只需证明 $t_{G/U}$ 为满射。
其像至少含有元素
$t_{G/U}(eU) = F(\text{id}_{G/U})(\gamma_U) = \gamma_U$.
对 $g \in G$，用 $R_g : G/U \to G/U$ 表示右乘。映射
$F(R_g) : F(G/U) \to F(G/U)$ 的不动点集等于 $F(\emptyset) = \emptyset$，
如果 $g \not \in U$；这是因为 $F$ 与等化子可交换。于是，如果
$g_1, \ldots, g_{|G/U|}$ 是 $G/U$
的一组代表元，则各元素 $F(R_{g_i})(\gamma_U)$ 两两不同，因而充满
$F(G/U)$。此时
$$
t_{G/U}(g_iU) = F(R_{g_i})(\gamma_U)
$$
证明完毕。
\end{proof}

\begin{example}
\label{pione-example-from-C-F-to-G-sets}
设 $\mathcal{C}$ 为范畴，而 $F : \mathcal{C} \to \textit{Sets}$ 为函子，
使得 $F(X)$ 对所有 $X \in \Ob(\mathcal{C})$ 都有限。由引理
\ref{pione-lemma-aut-inverse-limit}，$G = \text{Aut}(F)$ 典范地带有 profinite
拓扑群结构。得到函子
\begin{equation}
\label{pione-equation-remember}
\mathcal{C} \longrightarrow \textit{Finite-}G\textit{-Sets},\quad
X \longmapsto F(X)
\end{equation}
其中给 $F(X)$ 赋予所诱导的 $G$-作用。由我们对 $\text{Aut}(F)$ 上拓扑的构造，
该作用连续。
\end{example}

\noindent
定义 Galois 范畴的目的是刻画那些使函子 (\ref{pione-equation-remember}) 成为等价的
对 $(\mathcal{C}, F)$。我们如下定义 Galois 范畴。

\begin{definition}
\label{pione-definition-galois-category}
\begin{reference}
不同于 \cite[Expos\'e V, Definition 5.1]{SGA1} 中的定义。
可与 \cite[Definition 7.2.1]{BS} 比较。
\end{reference}
设 $\mathcal{C}$ 为范畴，而 $F : \mathcal{C} \to \textit{Sets}$ 为函子。如果
满足下列条件，则称对 $(\mathcal{C}, F)$ 为{\it Galois 范畴}：
\begin{enumerate}
\item $\mathcal{C}$ 有有限极限和有限余极限；
\item
\label{pione-item-connected-components}
$\mathcal{C}$ 的每个对象都是连通对象的有限（可以为空）余积；
\item $F(X)$ 对所有 $X \in \Ob(\mathcal{C})$ 都有限；
\item $F$ 反映同构\footnote{即，给定态射 $f$（属于 $\mathcal{C}$），如果 $F(f)$
是同构，则 $f$ 是同构。}，并且是正合的\footnote{这表示 $F$ 与有限极限和
有限余极限可交换；见 Categories, Section
\ref{categories-section-exact-functor}。}。
\end{enumerate}
这里称 $X \in \Ob(\mathcal{C})$ 连通，是指它不是始对象，并且对任意单态射
$Y \to X$，要么 $Y$ 是始对象，要么 $Y \to X$ 是同构。
\end{definition}

\noindent
{\bf 警告：}这个定义与大多数参考文献所给的定义不同（尽管最终会看出二者等价）。
具体而言，\cite[Expos\'e V, Definition 5.1]{SGA1} 把 Galois 范畴定义为
与 $\textit{Finite-}G\textit{-Sets}$ 等价的范畴，其中 $G$ 为某个 profinite 群。
随后 Grothendieck 用一组公理 (G1) -- (G6) 刻画 Galois 范畴；这些公理比上述
公理弱。我们作此选择，是为了强调有限极限和有限余极限的存在性以及函子 $F$
的正合性。其代价是稍后应用本节结果时必须多做一些工作。

\begin{lemma}
\label{pione-lemma-epi-mono}
设 $(\mathcal{C}, F)$ 为 Galois 范畴。设
$X \to Y \in \text{Arrows}(\mathcal{C})$。则
\begin{enumerate}
\item $F$ 忠实；
\item $X \to Y$ 是单态射
$\Leftrightarrow F(X) \to F(Y)$ 是单射；
\item $X \to Y$ 是满态射
$\Leftrightarrow F(X) \to F(Y)$ 是满射；
\item 对象 $A$（属于 $\mathcal{C}$）是始对象，当且仅当 $F(A) = \emptyset$；
\item 对象 $Z$（属于 $\mathcal{C}$）是终对象，当且仅当 $F(Z)$ 为单元集；
\item 如果 $X$ 和 $Y$ 连通，则 $X \to Y$ 是满态射；
\item
\label{pione-item-one-element}
如果 $X$ 连通，而 $a, b : X \to Y$ 为两个态射，则等式 $a = b$ 在如下条件下成立：
$F(a)$ 和 $F(b)$ 在 $F(X)$ 的某个元素上一致；
\item 如果 $X = \coprod_{i = 1, \ldots, n} X_i$ 且
$Y = \coprod_{j = 1, \ldots, m} Y_j$，其中 $X_i$、$Y_j$ 连通，则存在映射
$\alpha : \{1, \ldots, n\} \to \{1, \ldots, m\}$，使 $X \to Y$ 来自一族态射
$X_i \to Y_{\alpha(i)}$。
\end{enumerate}
\end{lemma}

\begin{proof}
(1) 的证明。设 $a, b : X \to Y$ 且 $F(a) = F(b)$。令 $E$ 为 $a$ 和 $b$
的等化子。于是 $F(E) = F(X)$；可知 $E = X$，因为 $F$ 反映同构。

\medskip\noindent
(2) 的证明。函子 $F$ 把态射 $X \to X \times_Y X$ 变为映射
$F(X) \to F(X) \times_{F(Y)} F(X)$，并且 $F$ 反映同构，故断言成立。

\medskip\noindent
(3) 的证明。函子 $F$ 把态射 $Y \amalg_X Y \to Y$ 变为映射
$F(Y) \amalg_{F(X)} F(Y) \to F(Y)$，并且 $F$ 反映同构，故断言成立。

\medskip\noindent
(4) 的证明。存在始对象 $A$，而且显然 $F(A) = \emptyset$。反之，如果对象
$X$ 满足 $F(X) = \emptyset$，则唯一映射 $A \to X$ 诱导双射
$F(A) \to F(X)$，故 $A \to X$ 是同构。

\medskip\noindent
(5) 的证明。存在终对象 $Z$，而且显然 $F(Z)$ 为单元集。反之，如果对象 $X$
的 $F(X)$ 为单元集，则唯一映射 $X \to Z$ 诱导双射 $F(X) \to F(Z)$，故
$X \to Z$ 是同构。

\medskip\noindent
(6) 的证明。等化子 $E$（对应两个映射 $Y \to Y \amalg_X Y$）不是
$\mathcal{C}$ 的始对象，因为 $X \to Y$ 经过 $E$ 分解，且
$F(X) \not = \emptyset$。因此 $E = Y$，断言得证。

\medskip\noindent
(\ref{pione-item-one-element}) 的证明。等化子 $E$ 是 $a$ 和 $b$ 的等化子，它带有单态射
$E \to X$，而 $F(E) \subset F(X)$ 是 $F(a)$ 和 $F(b)$ 取值相同的元素所成之集。
最后使用如下事实：要么 $E$ 是始对象，要么 $E = X$。

\medskip\noindent
(8) 的证明。对每个 $i, j$，$E_{ij} = X_i \times_Y Y_j$ 要么是始对象，
要么等于 $X_i$。选取 $s \in F(X_i)$，可知 $E_{ij} = X_i$ 当且仅当 $s$
映到 $F(Y_j) \subset F(Y)$ 的某个元素，故这种情形对唯一的
$j = \alpha(i)$ 成立。
\end{proof}

\noindent
由上述引理可知，给定 Galois 范畴中的连通对象 $X$；若该范畴记为
$(\mathcal{C}, F)$，则自同构群 $\text{Aut}(X)$ 的阶至多为 $|F(X)|$。事实上，给定
$s \in F(X)$ 和 $g \in \text{Aut}(X)$，由 (\ref{pione-item-one-element})，
$F(g)(s) = s$ 当且仅当 $g = \text{id}_X$。如果等号成立，就称 $X$ 是
{\it Galois 的}。等价地，$X$ 是 Galois 的，当且仅当它连通，并且
$\text{Aut}(X)$ 在 $F(X)$ 上传递作用。

\begin{lemma}
\label{pione-lemma-galois}
设 $(\mathcal{C}, F)$ 为 Galois 范畴。对任意连通对象 $X$（属于 $\mathcal{C}$），
存在 Galois 对象 $Y$ 和态射 $Y \to X$。
\end{lemma}

\begin{proof}
以下将不再明言地使用引理 \ref{pione-lemma-epi-mono} 的结果。令
$n = |F(X)|$。考虑 $X^n$ 及 $S_n$ 在其上的自然作用。设
$$
X^n = \coprod\nolimits_{t \in T} Z_t
$$
是连通对象分解。选取 $t$，使 $F(Z_t)$ 含有 $(s_1, \ldots, s_n)$，其中各
$s_i$ 两两不同。若 $(s'_1, \ldots, s'_n) \in F(Z_t)$ 是另一元素，则我们断言
各 $s'_i$ 也两两不同。否则，比如 $s'_i = s'_j$，则 $Z_t$ 是 $X^{n - 1}$
的某个连通分支在对角态射
$$
\Delta_{ij} : X^{n - 1} \longrightarrow X^n
$$
下的像。由于连通对象之间的态射是满态射，并且应用 $F$ 后诱导满射，
将推出 $s_i = s_j$，与假设矛盾。

\medskip\noindent
令 $G \subset S_n$ 为满足 $g(Z_t) = Z_t$ 的元素所成子群。考察 $S_n$ 在
$$
F(X)^n = F(X^n) = \coprod\nolimits_{t' \in T} F(Z_{t'})
$$
上的作用，可知
$G = \{g \in S_n \mid g(s_1, \ldots, s_n) \in F(Z_t)\}$。现选取第二个元素
$(s'_1, \ldots, s'_n) \in F(Z_t)$。上面已证明各 $s'_i$ 两两不同。因此可找到
$g \in S_n$，使 $g(s_1, \ldots, s_n) = (s'_1, \ldots, s'_n)$。换言之，
$G$ 在 $F(Z_t)$ 上的作用是传递的，证明完毕。
\end{proof}

\noindent
下面是一个关键引理。

\begin{lemma}
\label{pione-lemma-tame}
\begin{reference}
可与 \cite[Definition 7.2.4]{BS} 比较。
\end{reference}
设 $(\mathcal{C}, F)$ 为 Galois 范畴。设 $G = \text{Aut}(F)$ 如例
\ref{pione-example-from-C-F-to-G-sets} 所述。对任意连通对象 $X$（属于 $\mathcal{C}$），
$G$ 在 $F(X)$ 上的作用是传递的。
\end{lemma}

\begin{proof}
以下将不再明言地使用引理 \ref{pione-lemma-epi-mono} 的结果。令 $I$ 为
$\mathcal{C}$ 中 Galois 对象同构类所成之集。对每个 $i \in I$，令 $X_i$
为该同构类的代表。选取 $\gamma_i \in F(X_i)$（$i \in I$）。如下在 $I$ 上定义偏序：
$i \geq i'$ 当且仅当存在态射 $f_{ii'} : X_i \to X_{i'}$。给定这样的态射，
可以后复合一个自同构 $X_{i'} \to X_{i'}$，以保证
$F(f_{ii'})(\gamma_i) = \gamma_{i'}$。在这种规范化下，态射 $f_{ii'}$ 唯一。
注意 $I$ 是有向偏序集（Categories, Definition
\ref{categories-definition-directed-set}）：如果 $i_1, i_2 \in I$，则存在
Galois 对象 $Y$ 及态射 $Y \to X_{i_1} \times X_{i_2}$；这是对
$X_{i_1} \times X_{i_2}$ 的某个连通分支应用引理 \ref{pione-lemma-galois} 得到的。
于是 $Y \cong X_i$，其中某个 $i \in I$，且 $i \geq i_1$、$i \geq I_2$。

\medskip\noindent
我们断言函子 $F$ 同构于函子 $F'$，后者把 $X$ 送到
$$
F'(X) = \colim_I \Mor_\mathcal{C}(X_i, X)
$$
这是通过如下定义的函子变换 $t : F' \to F$ 实现的：给定 $f : X_i \to X$，
置 $t_X(f) = F(f)(\gamma_i)$。由 (\ref{pione-item-one-element}) 可知 $t_X$ 是单射。
为证满射，取 $\gamma \in F(X)$。由 Galois 范畴的定义，可立即归约到 $X$
连通的情形。再由引理 \ref{pione-lemma-galois}，可设 $X$ 是 Galois 的。此时
$X$ 同构于 $X_i$，其中 $i$ 是某个指标；并且可以选取同构 $X_i \to X$，使 $\gamma_i$ 映到
$\gamma$（由 Galois 对象的定义）。因此 $t$ 是同构。

\medskip\noindent
置 $A_i = \text{Aut}(X_i)$。我们断言，对 $i \geq i'$，存在典范映射
$h_{ii'} : A_i \to A_{i'}$，使得对所有 $a \in A_i$，图
$$
\xymatrix{
X_i \ar[d]_a \ar[r]_{f_{ii'}} & X_{i'} \ar[d]^{h_{ii'}(a)} \\
X_i \ar[r]^{f_{ii'}} & X_{i'}
}
$$
可交换。具体而言，令 $h_{ii'}(a) = a' : X_{i'} \to X_{i'}$ 为满足
$F(a')(\gamma_{i'}) = F(f_{ii'} \circ a)(\gamma_i)$.
的唯一自同构。如前，这使该图可交换，而且选择是唯一的。因此
$h_{i'i''} \circ h_{ii'} = h_{ii''}$
只要 $i \geq i' \geq i''$。由于 $F(X_i) \to F(X_{i'})$ 为满射，
可知 $A_i \to A_{i'}$ 为满射。取逆极限，得到群
$$
A = \lim_I A_i
$$
这是 profinite 群，因为各自同构群有限。由 Categories, Lemma
\ref{categories-lemma-nonempty-limit}，映射 $A \to A_i$ 对所有 $i$ 都是满射。

\medskip\noindent
由于 $A$ 的元素作用于逆系统 $X_i$，通过预复合得到 $A$ 在 $F'$ 上的右作用。
换言之，得到同态 $A^{opp} \to G$。由于 $A \to A_i$ 为满射，可知 $G$
在 $F(X_i)$ 上传递作用，对所有 $i$ 均如此。又由于每个连通对象都被某个 $X_i$ 支配，引理得证。
\end{proof}

\begin{proposition}
\label{pione-proposition-galois}
\begin{reference}
这是 \cite[Expos\'e V]{SGA1} 的弱版本。
证明取自 \cite[Theorem 7.2.5]{BS}。
\end{reference}
设 $(\mathcal{C}, F)$ 为 Galois 范畴。设 $G = \text{Aut}(F)$ 如例
\ref{pione-example-from-C-F-to-G-sets} 所述。函子
$F : \mathcal{C} \to \textit{Finite-}G\textit{-Sets}$
(\ref{pione-equation-remember}) 是等价。
\end{proposition}

\begin{proof}
以下将不再明言地使用引理 \ref{pione-lemma-epi-mono} 的结果。特别地，我们知道该函子
忠实。由引理 \ref{pione-lemma-tame}，对任意连通的 $X$，$G$ 在 $F(X)$ 上的作用
传递。因此，对一般的 $X$，函子 $F$ 把 $X$ 的连通分支（其存在性来自
Galois 范畴的一条公理）送到 $G$ 在 $F(X)$ 上的轨道。设 $X$ 和 $Y$ 为对象，
而 $s : F(X) \to F(Y)$ 为映射。于是
$\Gamma_s \subset F(X) \times F(Y)$（即 $s$ 的图像）是连通分支的并。因此存在
一个 $Z$，它是 $X \times Y$ 的若干连通分支之并，并带有单态射 $Z \to X \times Y$，
且 $F(Z) = \Gamma_s$。由于 $F(Z) \to F(X)$ 为双射，$Z \to X$ 是同构，
从而 $s = F(f)$，其中 $f : X \cong Z \to Y$ 为复合。因此 $F$ 全忠实。

\medskip\noindent
为完成证明，我们来证明 $F$ 本质满。只需证明 $G/H$ 位于本质像中，其中
$H \subset G$ 是任意有限指数开子群。由 $G$ 上拓扑的定义，存在有限族对象
$X_i$，使
$$
\Ker(G \longrightarrow \prod\nolimits_i \text{Aut}(F(X_i)))
$$
包含于 $H$。可设 $X_i$ 连通，对所有 $i$ 均如此。由引理 \ref{pione-lemma-galois}，
可以选取 Galois 对象 $Y$，使其映到 $\prod X_i$ 的某个连通分支。选取同构
$F(Y) = G/U$（位于 $G\textit{-sets}$ 中），其中 $U \subset G$ 为某个开子群。由于 $Y$
是 Galois 的，群 $\text{Aut}(Y) = \text{Aut}_{G\textit{-Sets}}(G/U)$
在 $F(Y) = G/U$ 上传递作用。这说明 $U$ 正规。由于 $F(Y)$ 满射到
$F(X_i)$，对每个 $i$ 均如此，可知 $U \subset H$。令 $M \subset \text{Aut}(Y)$ 为对应于
$$
(H/U)^{opp} \subset (G/U)^{opp} = \text{Aut}_{G\textit{-Sets}}(G/U)
= \text{Aut}(Y).
$$
的有限子群。置 $X = Y/M$，即 $X$ 是各箭头 $m : Y \to Y$（$m \in M$）
的余等化子。由于 $F$ 正合，可知 $F(X) = G/H$，证明完毕。
\end{proof}

\begin{lemma}
\label{pione-lemma-functoriality-galois}
设 $(\mathcal{C}, F)$ 和 $(\mathcal{C}', F')$ 为 Galois 范畴。设
$H : \mathcal{C} \to \mathcal{C}'$ 为正合函子。存在同构
$t : F' \circ H \to F$。$t$ 的选择决定连续同态
$h : G' = \text{Aut}(F') \to \text{Aut}(F) = G$ 以及 $2$-交换图
$$
\xymatrix{
\mathcal{C} \ar[r]_H \ar[d] & \mathcal{C}' \ar[d] \\
\textit{Finite-}G\textit{-Sets} \ar[r]^h &
\textit{Finite-}G'\textit{-Sets}
}
$$
映射 $h$ 与 $t$ 无关，仅允许相差 $G$ 的内自同构。反之，给定连续同态
$h : G' \to G$，存在正合函子 $H : \mathcal{C} \to \mathcal{C}'$ 和同构
$t$，按上述方式恢复 $h$。
\end{lemma}

\begin{proof}
由命题 \ref{pione-proposition-galois} 和引理 \ref{pione-lemma-single-out-profinite}，
可设 $\mathcal{C} = \textit{Finite-}G\textit{-Sets}$ 且 $F$ 为遗忘函子；
对 $\mathcal{C}'$ 同理。于是 $t$ 的存在性来自引理
\ref{pione-lemma-second-fundamental-functor}。映射 $h$ 来自经由 $t$ 的结构传递。
图的交换性显然。$h$ 在相差 $G$ 中某个元素的内共轭意义下唯一，这是因为
$t$ 的选择在相差 $G$ 的一个元素意义下唯一。最后一个断言是直接的。
\end{proof}





\section{函子与同态}
\label{pione-section-translation}

\noindent
设 $(\mathcal{C}, F)$、$(\mathcal{C}', F')$、$(\mathcal{C}'', F'')$
为 Galois 范畴。置 $G = \text{Aut}(F)$、$G' = \text{Aut}(F')$，以及
$G'' = \text{Aut}(F'')$。设 $H : \mathcal{C} \to \mathcal{C}'$
和 $H' : \mathcal{C}' \to \mathcal{C}''$ 为正合函子。
设 $h : G' \to G$ 和 $h' : G'' \to G'$ 为引理
\ref{pione-lemma-functoriality-galois} 所述的相应连续同态。本节考虑相应的 $2$-交换图
\begin{equation}
\label{pione-equation-translation}
\vcenter{
\xymatrix{
\mathcal{C} \ar[r]_H \ar[d] &
\mathcal{C}' \ar[r]_{H'} \ar[d] &
\mathcal{C}'' \ar[d] \\
\textit{Finite-}G\textit{-Sets} \ar[r]^h &
\textit{Finite-}G'\textit{-Sets} \ar[r]^{h'} &
\textit{Finite-}G''\textit{-Sets}
}
}
\end{equation}
并把序列 $1 \to G'' \to G' \to G \to 1$ 的正合性质与函子 $H$ 和 $H'$
的性质联系起来。

\begin{lemma}
\label{pione-lemma-functoriality-galois-surjective}
在图 (\ref{pione-equation-translation}) 中，下列条件等价：
\begin{enumerate}
\item $h : G' \to G$ 是满射；
\item $H : \mathcal{C} \to \mathcal{C}'$ 全忠实；
\item 如果 $X \in \Ob(\mathcal{C})$ 连通，则 $H(X)$ 连通；
\item 如果 $X \in \Ob(\mathcal{C})$ 连通，且存在态射
$*' \to H(X)$（位于 $\mathcal{C}'$ 中），则存在态射 $* \to X$；以及
\item 对任意对象 $X$（属于 $\mathcal{C}$），映射
$\Mor_\mathcal{C}(*, X) \to \Mor_{\mathcal{C}'}(*', H(X))$
为双射。
\end{enumerate}
这里 $*$ 和 $*'$ 分别是 $\mathcal{C}$ 和 $\mathcal{C}'$ 的终对象。
\end{lemma}

\begin{proof}
(5) $\Rightarrow$ (4) 和 (2) $\Rightarrow$ (5) 显然。

\medskip\noindent
假设 (3)。设 $X$ 为 $\mathcal{C}$ 的连通对象，并设 $*' \to H(X)$
为态射。由 (3)，$H(X)$ 连通，故 $*' \to H(X)$ 是同构。因此相应的
$G'$-集（对应于 $H(X)$）恰有一个元素，这意味着相应的 $G$-集
（对应于 $X$）有一个元素，进而意味着 $X$ 同构于 $\mathcal{C}$ 的终对象；
特别地，存在映射 $* \to X$。由此可见 (3) $\Rightarrow$ (4)。

\medskip\noindent
如果 (1) 成立，则函子
$\textit{Finite-}G\textit{-Sets} \to \textit{Finite-}G'\textit{-Sets}$
全忠实：在这种情形下，$G$-集之间的映射与 $G$ 的作用可交换，当且仅当它与
$G'$ 的作用可交换。因此 (1) $\Rightarrow$ (2)。

\medskip\noindent
如果 (1) 成立，则对 $G$-集 $X$，其 $G$-轨道与 $G'$-轨道一致。
因此 (1) $\Rightarrow$ (3)。

\medskip\noindent
为完成证明，只需证明 (4) 蕴含 (1)。如果 (1) 不成立，即 $h$ 不是满射，
则存在开子群 $U \subset G$，它包含 $h(G')$ 但不等于 $G$。于是有限
$G$-集 $M = G/U$ 上的作用是传递的，但 $G'$ 有一个不动点。对象
$X$（属于 $\mathcal{C}$，对应于 $M$）将与 (3) 矛盾。由此可见
(3) $\Rightarrow$ (1)，证明完毕。
\end{proof}

\begin{lemma}
\label{pione-lemma-composition-trivial}
在图 (\ref{pione-equation-translation}) 中，下列条件等价：
\begin{enumerate}
\item $h \circ h'$ 是平凡的；以及
\item $H' \circ H$ 的像由同构于终对象有限余积的对象组成。
\end{enumerate}
\end{lemma}

\begin{proof}
可以把 $H$ 和 $H'$ 替换为典范函子
$\textit{Finite-}G\textit{-Sets} \to \textit{Finite-}G'\textit{-Sets}
\to \textit{Finite-}G''\textit{-Sets}$；这些函子由 $h$ 和 $h'$ 决定。
于是所述命题变为：$G''$ 在每个
$G$-集上的作用平凡，当且仅当同态 $G'' \to G$ 平凡。这是显然的。
\end{proof}

\begin{lemma}
\label{pione-lemma-functoriality-galois-ses}
在图 (\ref{pione-equation-translation}) 中，下列条件等价：
\begin{enumerate}
\item 序列 $G'' \xrightarrow{h'} G' \xrightarrow{h} G \to 1$
在如下意义下正合：$h$ 是满射，$h \circ h'$ 平凡，并且 $\Ker(h)$ 是包含
$\Im(h')$ 的最小闭正规子群；
\item $H$ 全忠实，而且对象 $X'$（属于 $\mathcal{C}'$）位于 $H$ 的本质像中，
当且仅当 $H'(X')$ 同构于终对象的有限余积；以及
\item $H$ 全忠实，$H \circ H'$ 把每个对象送到终对象的有限余积，并且对
对象 $X'$（属于 $\mathcal{C}'$），如果 $H'(X')$ 是终对象的有限余积，则存在
对象 $X$（属于 $\mathcal{C}$）及满态射 $H(X) \to X'$。
\end{enumerate}
\end{lemma}

\begin{proof}
由引理 \ref{pione-lemma-functoriality-galois-surjective} 和
\ref{pione-lemma-composition-trivial}，可设 $H$ 全忠实，$h$ 是满射，
$H' \circ H$ 把对象送到终对象的不交并，并且 $h \circ h'$ 平凡。
令 $N \subset G'$ 为包含 $h'$ 的像的最小闭正规子群。显然
$N \subset \Ker(h)$。可设函子 $H$ 和 $H'$ 是典范函子
$\textit{Finite-}G\textit{-Sets} \to \textit{Finite-}G'\textit{-Sets}
\to \textit{Finite-}G''\textit{-Sets}$；这些函子由 $h$ 和 $h'$ 决定。

\medskip\noindent
假设 (2) 成立。这意味着，对有限 $G'$-集 $X'$，如果 $G''$ 的作用平凡，
则 $G'$ 的作用经过 $G$ 分解。把它应用于 $X' = G'/U'N$，其中 $U'$ 是
$G'$ 的充分小开子群。于是有 $\Ker(h) \subset U'N$，对所有 $U'$ 均如此。
由于 $N$ 闭，这蕴含 $\Ker(h) \subset N$，即 (1) 成立。

\medskip\noindent
假设 (1) 成立。这意味着 $N = \Ker(h)$。设 $X'$ 为有限 $G'$-集，且
$G''$ 的作用平凡。这意味着 $\Ker(G' \to \text{Aut}(X'))$ 是包含
$\Im(h')$ 的闭正规子群。因此它包含 $N = \Ker(h)$，而 $G'$ 在 $X'$
上的作用经过 $G$ 分解，即 (2) 成立。

\medskip\noindent
假设 (3) 成立。这意味着，对有限 $G'$-集 $X'$，如果 $G''$ 的作用平凡，
则存在 $G'$-集的满射 $X \to X'$，其中 $X$ 是 $G$-集。显然，这意味着
$G'$ 在 $X'$ 上的作用经过 $G$ 分解，即 (2) 成立。

\medskip\noindent
(2) $\Rightarrow$ (3) 立即成立。证明完毕。
\end{proof}

\begin{lemma}
\label{pione-lemma-functoriality-galois-injective}
在图 (\ref{pione-equation-translation}) 中，下列条件等价：
\begin{enumerate}
\item $h'$ 是单射；以及
\item 对每个连通对象 $X''$（属于 $\mathcal{C}''$），存在对象 $X'$
（属于 $\mathcal{C}'$）以及图
$$
X'' \leftarrow Y'' \rightarrow H(X')
$$
它位于 $\mathcal{C}''$ 中，其中 $Y'' \to X''$ 是满态射，而
$Y'' \to H(X')$ 是单态射。
\end{enumerate}
\end{lemma}

\begin{proof}
可以把 $H'$ 替换为有限 $G'$-集范畴与有限 $G''$-集范畴之间的相应函子。

\medskip\noindent
假设 $h' : G'' \to G'$ 是单射。设 $H'' \subset G''$ 为开子群。
由于 $G''$ 上的拓扑是由 $G'$ 诱导的拓扑，存在开子群 $H' \subset G'$，使
$(h')^{-1}(H') \subset H''$。于是所需图为
$$
G''/H'' \leftarrow G''/(h')^{-1}(H') \rightarrow G'/H'
$$
反之，假设 (2) 对函子
$\textit{Finite-}G'\textit{-Sets} \to \textit{Finite-}G''\textit{-Sets}$.
成立。设 $g'' \in \Ker(h')$。任取开子群 $H'' \subset G''$。
由假设，存在有限 $G'$-集 $X'$ 及图
$$
G''/H'' \leftarrow Y'' \rightarrow X'
$$
它由 $G''$-集构成，左箭头为满射，右箭头为单射。由于 $g''$ 位于 $h'$ 的核中，
可知 $g''$ 在 $X'$ 上作用平凡。因此 $g''$ 在 $Y''$ 上作用平凡，进而在
$G''/H''$ 上作用平凡。故 $g'' \in H''$。由于这对所有开子群都成立，可知
$g''$ 是单位元，正如所需。
\end{proof}

\begin{lemma}
\label{pione-lemma-functoriality-galois-normal}
在图 (\ref{pione-equation-translation}) 中，下列条件等价：
\begin{enumerate}
\item $h'$ 的像是正规的；以及
\item 对每个连通对象 $X'$（属于 $\mathcal{C}'$），如果存在从
$\mathcal{C}''$ 的终对象到 $H'(X')$ 的态射，则 $H'(X')$ 同构于终对象的有限余积。
\end{enumerate}
\end{lemma}

\begin{proof}
这转化为关于连续群同态 $h' : G'' \to G'$ 的如下陈述：$h'$ 的像正规，
当且仅当每个开子群 $U' \subset G'$，只要它包含 $h'(G'')$，也包含
$h'(G'')$ 的每个共轭。结论很容易由此推出；略去若干细节。
\end{proof}






\section{有限 \'etale 态射}
\label{pione-section-finite-etale}

\noindent
本节证明关于有限 \'etale 态射的足够多基本结果，以便构造 \'etale 基本群。

\medskip\noindent
设 $X$ 为概形。用记号 $\textit{F\'Et}_X$ 表示在 $X$ 上有限且 \'etale
的概形所成范畴。于是
\begin{enumerate}
\item $\textit{F\'Et}_X$ 的对象是有限 \'etale 态射
$Y \to X$，其目标为 $X$；以及
\item $\textit{F\'Et}_X$ 中从 $Y \to X$ 到 $Y' \to X$ 的态射，是使图
$Y \to Y'$
$$
\xymatrix{
Y \ar[rr] \ar[rd] & &  Y' \ar[ld] \\
& X
}
$$
可交换的态射。
\end{enumerate}
我们常把 $\textit{F\'Et}_X$ 的对象称为 $X$ 的{\it 有限 \'etale 覆盖}
（即使 $Y$ 为空）。事实上，在概形范畴上存在栈
$p : \textit{F\'Et} \to \Sch$，其在 $X$ 处的纤维正是刚定义的范畴
$\textit{F\'Et}_X$。见 Examples of Stacks, Section
\ref{examples-stacks-section-finite-etale}.

\begin{example}
\label{pione-example-finite-etale-geometric-point}
设 $k$ 为代数闭域，且 $X = \Spec(k)$。在这种情形下，
$\textit{F\'Et}_X$ 等价于有限集范畴。更一般地，当 $k$ 可分代数闭时仍然成立。
这是因为在 $k$ 上 \'etale 的概形是 $k$ 的有限可分扩域之谱的不交并；见
Morphisms, Lemma \ref{morphisms-lemma-etale-over-field}.
\end{example}

\begin{lemma}
\label{pione-lemma-finite-etale-covers-limits-colimits}
设 $X$ 为概形。范畴 $\textit{F\'Et}_X$ 有有限极限和有限余极限，并且对任意态射
$X' \to X$，基变换函子
$\textit{F\'Et}_X \to \textit{F\'Et}_{X'}$ 正合。
\end{lemma}

\begin{proof}
有限极限与左正合性。由 Categories, Lemma
\ref{categories-lemma-finite-limits-exist}，只需证明
$\textit{F\'Et}_X$ 有终对象和纤维积。这是显然的，因为 $X$ 上所有概形的范畴
有终对象（即 $X$）和纤维积。此外，在 $X$ 上有限 \'etale 的概形之纤维积仍
在 $X$ 上有限 \'etale。显然，基变换与这些运算可交换，因而基变换左正合
（Categories, Lemma
\ref{categories-lemma-characterize-left-exact}).

\medskip\noindent
有限余极限与右正合性。由 Categories, Lemma
\ref{categories-lemma-colimits-exist}，只需证明 $\textit{F\'Et}_X$
有有限余积和余等化子。有限余积由不交并给出（空余积是空概形）。
设 $a, b : Z \to Y$ 为 $\textit{F\'Et}_X$ 的两个态射。由于
$Z \to X$ 和 $Y \to X$ 有限 \'etale，可以写成
$Z = \underline{\Spec}(\mathcal{C})$ 和 $Y = \underline{\Spec}(\mathcal{B})$，
其中某些有限局部自由 $\mathcal{O}_X$-代数 $\mathcal{C}$ 和 $\mathcal{B}$
给出上述表示。态射 $a, b$ 诱导两个映射
$a^\sharp, b^\sharp : \mathcal{B} \to \mathcal{C}$。令
$\mathcal{A} = \text{Eq}(a^\sharp, b^\sharp)$ 为它们的等化子。如果
$$
\underline{\Spec}(\mathcal{A}) \longrightarrow X
$$
有限 \'etale，则显然这就是余等化子（毕竟，$\textit{F\'Et}_X$ 的任意对象
都可写成某个 $\mathcal{O}_X$-代数层的相对谱）。把 $X$ 替换为某个
\'etale 覆盖的成员后即可验证这一点（Descent, Lemmas
\ref{descent-lemma-descending-property-finite} 和
\ref{descent-lemma-descending-property-etale}）。因此由 \'Etale Morphisms,
Lemma \ref{etale-lemma-finite-etale-etale-local}，可以假设
$Y = \coprod_{i = 1, \ldots, n} X$ 和 $Z = \coprod_{j = 1, \ldots, m} X$。
此时
$$
\mathcal{C} = \prod\nolimits_{1 \leq j \leq m} \mathcal{O}_X
\quad\text{且}\quad
\mathcal{B} = \prod\nolimits_{1 \leq i \leq n} \mathcal{O}_X
$$
进一步替换为某个开覆盖的成员后，可设 $a, b$ 对应于映射
$a_s, b_s : \{1, \ldots, m\} \to \{1, \ldots, n\}$；也就是说，直和项
$X$（属于 $Z$ 并对应于指标 $j$）映到直和项 $X$（属于 $Y$ 并对应于指标
$a_s(j)$、相应地 $b_s(j)$），分别在态射 $a$、$b$ 下。
令 $\{1, \ldots, n\} \to T$ 为 $a_s, b_s$ 的余等化子。于是
$$
\mathcal{A} = \prod\nolimits_{t \in T} \mathcal{O}_X
$$
其谱显然在 $X$ 上有限 \'etale。略去它与基变换相容的验证。
因此基变换是右正合函子。
\end{proof}

\begin{remark}
\label{pione-remark-colimits-commute-forgetful}
设 $X$ 为概形。考虑自然函子
$F_1 : \textit{F\'Et}_X \to \Sch$ 和 $F_2 : \textit{F\'Et}_X \to \Sch/X$。
则
\begin{enumerate}
\item 函子 $F_1$ 和 $F_2$ 与有限余极限可交换。
\item 函子 $F_2$ 与有限极限可交换。
\item 函子 $F_1$ 与连通有限极限可交换，即与等化子和纤维积可交换。
\end{enumerate}
关于极限的结果立即来自引理
\ref{pione-lemma-finite-etale-covers-limits-colimits} 证明中的讨论以及
Categories, Lemma \ref{categories-lemma-connected-limit-over-X}。
显然，$F_1$ 和 $F_2$ 与有限余积可交换。由 Categories, Lemma
\ref{categories-lemma-characterize-left-exact}
的对偶，只需证明 $F_1$ 和 $F_2$ 与余等化子可交换。在引理
\ref{pione-lemma-finite-etale-covers-limits-colimits} 的证明中，我们看到
$\textit{F\'Et}_X$ 中的余等化子在 \'etale 局部具有如下形式：
$$
\xymatrix{
\coprod_{j \in J} U \ar@<1ex>[r]^a \ar@<-1ex>[r]_b &
\coprod_{i \in I} U \ar[r] &
\coprod_{t \in \text{Coeq}(a, b)} U
}
$$
这在概形范畴中显然是余等化子。因此，结论来自余等化子性质是 fpqc 局部的这一事实；
其精确表述见 Descent, Lemma \ref{descent-lemma-coequalizer-fpqc-local}。
\end{remark}

\begin{lemma}
\label{pione-lemma-internal-hom-finite-etale}
设 $X$ 为概形。给定 $U, V$ 在 $X$ 上有限 \'etale，存在概形 $W$ 在 $X$
上有限 \'etale，使
$$
\Mor_X(X, W) = \Mor_X(U, V)
$$
并且在任意基变换后同一结论仍然成立。
\end{lemma}

\begin{proof}
由 More on Morphisms, Lemma
\ref{more-morphisms-lemma-hom-from-finite-locally-free-separated-lqf}，
存在概形 $W$，它表示 $\mathit{Mor}_X(U, V)$。（这里使用 Morphisms, Lemmas
\ref{morphisms-lemma-etale-locally-quasi-finite} 所说明的 \'etale 态射局部拟有限，
以及有限态射是分离的。）这个概形在任意基变换后显然仍满足该公式。
为完成证明，必须证明 $W \to X$ 有限 \'etale。把 $X$ 替换为某个
\'etale 覆盖的成员后即可验证这一点（Descent, Lemmas
\ref{descent-lemma-descending-property-finite} 和
\ref{descent-lemma-descending-property-separated}）。因此由 \'Etale Morphisms,
Lemma \ref{etale-lemma-finite-etale-etale-local}，可设
$U = \coprod_{i = 1, \ldots, n} X$ 和
$V = \coprod_{j = 1, \ldots, m} X$。在这种情形下，直接检验可得（略去细节）
$W = \coprod_{\alpha : \{1, \ldots, n\} \to \{1, \ldots, m\}} X$
，证明完毕。
\end{proof}

\noindent
设 $X$ 为概形。$X$ 的一个{\it 几何点}是态射
$\Spec(k) \to X$，其中 $k$ 代数闭。这样的点通常记作 $\overline{x}$，
即用带上划线的小写字母表示。我们常用 $\overline{x}$ 同时表示概形
$\Spec(k)$ 和该态射，并用 $\kappa(\overline{x})$ 表示 $k$。称
$\overline{x}$ {\it 位于} $x$ {\it 之上}，是指 $x \in X$ 是
$\overline{x}$ 的像。进一步讨论见
\'Etale Cohomology, Section \ref{etale-cohomology-section-stalks}.
给定 $\overline{x}$ 和 \'etale 态射 $U \to X$，可以考虑
$$
|U_{\overline{x}}| : \text{the underlying set of points of the
scheme }U_{\overline{x}} = U \times_X \overline{x}
$$
由于 $U_{\overline{x}}$ 作为 $\overline{x}$ 上的概形，是若干个
$\overline{x}$ 的不交并（Morphisms, Lemma
\ref{morphisms-lemma-etale-over-field}），也可以把该集合描述为
$$
|U_{\overline{x}}| =
\left\{
\begin{matrix}
\text{commutative} \\
\text{diagrams}
\end{matrix}
\vcenter{
\xymatrix{
\overline{x} \ar[rd]_{\overline{x}} \ar[r]_{\overline{u}} & U \ar[d] \\
& X
}
}
\right\}
$$
对应 $U \mapsto |U_{\overline{x}}|$ 是一个函子，常记作
$F_{\overline{x}}$。

\begin{lemma}
\label{pione-lemma-finite-etale-connected-galois-category}
设 $X$ 为连通概形。设 $\overline{x}$ 为几何点。函子
$$
F_{\overline{x}} : \textit{F\'Et}_X \longrightarrow \textit{Sets},\quad
Y \longmapsto |Y_{\overline{x}}|
$$
定义了一个 Galois 范畴（Definition \ref{pione-definition-galois-category}）。
\end{lemma}

\begin{proof}
把 $\textit{F\'Et}_{\overline{x}}$ 与有限集范畴识别后（Example
\ref{pione-example-finite-etale-geometric-point}），可知函子
$F_{\overline{x}}$ 正是态射 $\overline{x} \to X$ 的基变换函子。因此，由引理
\ref{pione-lemma-finite-etale-covers-limits-colimits} 可知 $\textit{F\'Et}_X$
有有限极限和有限余极限，并且 $F_{\overline{x}}$ 正合。还将使用如下事实：
$\textit{F\'Et}_X$ 中的有限极限与 $X$ 上概形范畴中的相应有限极限一致；见
Remark \ref{pione-remark-colimits-commute-forgetful}。

\medskip\noindent
如果 $Y' \to Y$ 是 $\textit{F\'Et}_X$ 中的单态射，则
$Y' \to Y' \times_Y Y'$ 是同构，因而 $Y' \to Y$ 是概形的单态射。
于是 $Y' \to Y$ 是开浸入（\'Etale Morphisms, Theorem
\ref{etale-theorem-etale-radicial-open}）。由于 $Y'$ 在 $X$ 上有限，而
$Y$ 在 $X$ 上分离，态射 $Y' \to Y$ 有限（Morphisms, Lemma
\ref{morphisms-lemma-finite-permanence}），因而是闭的（Morphisms, Lemma
\ref{morphisms-lemma-finite-proper}），所以它是 $Y$ 的开闭子概形之包含。
由此，$Y$ 是范畴 $\textit{F\'Et}_X$ 的连通对象（如 Definition
\ref{pione-definition-galois-category}），当且仅当 $Y$ 作为概形连通。于是由
Topology, Lemma \ref{topology-lemma-finite-fibre-connected-components}，
$Y$ 无论作为概形还是在 Definition \ref{pione-definition-galois-category}
的意义下，都是其连通分支的有限余积。

\medskip\noindent
设 $Y \to Z$ 为 $\textit{F\'Et}_X$ 中的态射，它诱导双射
$F_{\overline{x}}(Y) \to F_{\overline{x}}(Z)$。必须证明 $Y \to Z$
是同构。由上可设 $Z$ 连通。由于 $Y \to Z$ 有限 \'etale，因而有限局部自由，
只需证明 $Y \to Z$ 是次数为 $1$ 的有限局部自由态射。这在 $Z$ 中位于
$\overline{x}$ 之上的任意点的某个邻域内成立；由于 $Z$ 连通且次数局部常值，
结论成立。
\end{proof}



\section{基本群}
\label{pione-section-fundamental-groups}

\noindent
本节定义 Grothendieck 的代数基本群。由于引理
\ref{pione-lemma-finite-etale-connected-galois-category}，下述定义是有意义的。

\begin{definition}
\label{pione-definition-fundamental-group}
设 $X$ 为连通概形。设 $\overline{x}$ 为 $X$ 的几何点。概形 $X$ 以
$\overline{x}$ 为{\it 基点}的{\it 基本群}是群
$$
\pi_1(X, \overline{x}) = \text{Aut}(F_{\overline{x}})
$$
即纤维函子
$F_{\overline{x}} : \textit{F\'Et}_X \to \textit{Sets}$
的自同构群，并赋予引理 \ref{pione-lemma-aut-inverse-limit} 所给的典范 profinite 拓扑。
\end{definition}

\noindent
把上述结果与 Section \ref{pione-section-galois} 的内容结合，得到下述定理。

\begin{theorem}
\label{pione-theorem-fundamental-group}
设 $X$ 为连通概形。设 $\overline{x}$ 为 $X$ 的几何点。
\begin{enumerate}
\item 纤维函子 $F_{\overline{x}}$ 定义范畴等价
$$
\textit{F\'Et}_X \longrightarrow
\textit{Finite-}\pi_1(X, \overline{x})\textit{-Sets}
$$
\item 给定第二个几何点 $\overline{x}'$（属于 $X$），存在同构
$t : F_{\overline{x}} \to F_{\overline{x}'}$。它给出与 (1) 中等价相容的同构
$\pi_1(X, \overline{x}) \to \pi_1(X, \overline{x}')$。在相差内共轭的意义下，
该同构与 $t$ 无关。
\item 给定连通概形之间的态射 $f : X \to Y$，记
$\overline{y} = f \circ \overline{x}$。存在典范连续同态
$$
f_* : \pi_1(X, \overline{x}) \to \pi_1(Y, \overline{y})
$$
使图
$$
\xymatrix{
\textit{F\'Et}_Y \ar[r]_{\text{base change}} \ar[d]_{F_{\overline{y}}} &
\textit{F\'Et}_X \ar[d]^{F_{\overline{x}}} \\
\textit{Finite-}\pi_1(Y, \overline{y})\textit{-Sets} \ar[r]^{f_*} &
\textit{Finite-}\pi_1(X, \overline{x})\textit{-Sets}
}
$$
可交换。
\end{enumerate}
\end{theorem}

\begin{proof}
(1) 来自引理 \ref{pione-lemma-finite-etale-connected-galois-category} 和命题
\ref{pione-proposition-galois}。(2) 是引理 \ref{pione-lemma-functoriality-galois}
的特殊情形。对于 (3)，注意图
$$
\xymatrix{
\textit{F\'Et}_Y \ar[r] \ar[d]_{F_{\overline{y}}} &
\textit{F\'Et}_X \ar[d]^{F_{\overline{x}}} \\
\textit{Sets} \ar@{=}[r] & \textit{Sets}
}
$$
可交换（实际上可交换，而不只是 $2$-交换），因为
$\overline{y} = f \circ \overline{x}$。因此可以把引理
\ref{pione-lemma-functoriality-galois} 与隐含的函子变换一起应用，从而得到 (3)。
\end{proof}

\begin{lemma}
\label{pione-lemma-fundamental-group-Galois-group}
设 $K$ 为域，并置 $X = \Spec(K)$。设 $\overline{K}$ 为代数闭包，并以
$\overline{x} : \Spec(\overline{K}) \to X$ 表示相应的几何点。设
$K^{sep} \subset \overline{K}$ 为可分闭包。
\begin{enumerate}
\item 引理 \ref{pione-lemma-sheaves-point} 的函子诱导等价
$$
\textit{F\'Et}_X \longrightarrow
\textit{Finite-}\text{Gal}(K^{sep}/K)\textit{-Sets}.
$$
它与 $F_{\overline{x}}$ 及函子
$\textit{Finite-}\text{Gal}(K^{sep}/K)\textit{-Sets} \to \textit{Sets}$.
相容。
\item 它诱导典范同构
$$
\text{Gal}(K^{sep}/K) \longrightarrow \pi_1(X, \overline{x})
$$
，这是 profinite 拓扑群的同构。
\end{enumerate}
\end{lemma}

\begin{proof}
引理 \ref{pione-lemma-sheaves-point} 的函子与函子 $F_{\overline{x}}$ 相同，
因为对任意 $Y$，若它在 $X$ 上 \'etale，则有
$$
\Mor_X(\Spec(\overline{K}), Y) = \Mor_X(\Spec(K^{sep}), Y)
$$
事实上，如引理 \ref{pione-lemma-sheaves-point} 的证明所示，
$Y = \coprod_{i \in I} \Spec(L_i)$，其中 $L_i/K$ 是 $K$ 的有限可分扩张。
因此任意 $K$-代数同态 $L_i \to \overline{K}$ 都经过 $K^{sep}$ 分解。
还要注意，$F_{\overline{x}}(Y)$ 有限，当且仅当 $I$ 有限，又当且仅当
$Y \to X$ 有限 \'etale。这证明了 (1)。

\medskip\noindent
(2) 是 (1)、引理 \ref{pione-lemma-functoriality-galois} 和引理
\ref{pione-lemma-single-out-profinite} 的形式推论。（另见下述评注。）
\end{proof}

\begin{remark}
\label{pione-remark-variance}
在引理 \ref{pione-lemma-fundamental-group-Galois-group} 的情形下，下面给出同构
$\text{Gal}(K^{sep}/K) \to
\pi_1(X, \overline{x}) = \text{Aut}(F_{\overline{x}})$ 的更明确构造。注意
$\text{Gal}(K^{sep}/K) = \text{Aut}(\overline{K}/K)$
，因为 $\overline{K}$ 是 $K^{sep}$ 的完美化。由于
$F_{\overline{x}}(Y) = \Mor_X(\Spec(\overline{K}), Y)$，可以考虑映射
$$
\text{Aut}(\overline{K}/K) \times F_{\overline{x}}(Y) \to F_{\overline{x}}(Y),
\quad
(\sigma, \overline{y}) \mapsto
\sigma \cdot \overline{y} = \overline{y} \circ \Spec(\sigma)
$$
这是一个作用，因为
$$
\sigma\tau \cdot \overline{y} =
\overline{y} \circ \Spec(\sigma\tau) =
\overline{y} \circ \Spec(\tau) \circ \Spec(\sigma) =
\sigma \cdot (\tau \cdot \overline{y})
$$
该作用关于 $Y \in \textit{F\'Et}_X$ 具有函子性，从而得到所需映射。
\end{remark}





\section{连通概形的 Galois 覆盖}
\label{pione-section-finite-etale-under-galois}

\noindent
设 $X$ 为带有几何点 $\overline{x}$ 的连通概形。由于
$F_{\overline{x}} : \textit{F\'Et}_X \to \textit{Sets}$ 是 Galois 范畴
（Lemma \ref{pione-lemma-finite-etale-connected-galois-category}），Section
\ref{pione-section-galois} 的内容适用。本节把若干术语和结果明确转移到概形与有限
\'etale 态射的情形。

\medskip\noindent
称有限 \'etale 态射 $Y \to X$ 为{\it Galois 覆盖}，如果 $Y$ 定义了
$\textit{F\'Et}_X$ 的 Galois 对象。对有限 \'etale 态射 $Y \to X$，置
$G = \text{Aut}_X(Y)$，则下列条件等价：
\begin{enumerate}
\item $Y$ 是 $X$ 的 Galois 覆盖；
\item $Y$ 连通，且 $|G|$ 等于 $Y \to X$ 的次数；
\item $Y$ 连通，且 $G$ 在 $F_{\overline{x}}(Y)$ 上传递作用；以及
\item $Y$ 连通，且 $G$ 在 $F_{\overline{x}}(Y)$ 上单传递作用。
\end{enumerate}
这立即来自 Section \ref{pione-section-galois} 中的讨论。

\medskip\noindent
对任意有限 \'etale 态射 $f : Y \to X$，如果 $Y$ 连通，则存在有限
\'etale Galois 覆盖 $Y' \to X$，它支配 $Y$（Lemma \ref{pione-lemma-galois}）。

\medskip\noindent
$\textit{F\'Et}_X$ 的 Galois 对象经由等价
$$
F_{\overline{x}} : \textit{F\'Et}_X \to
\textit{Finite-}\pi_1(X, \overline{x})\textit{-Sets}
$$
（Theorem \ref{pione-theorem-fundamental-group}），对应于有限
$\pi_1(X, \overline{x})\textit{-Sets}$，它们形如
$G = \pi_1(X, \overline{x})/H$，其中 $H$ 为正规开子群。等价地，
如果 $G$ 是有限群，且 $\pi_1(X, \overline{x}) \to G$ 是连续满射，则把
$G$ 看作 $\pi_1(X, \overline{x})$-集时，它对应于 Galois 覆盖。

\medskip\noindent
如果 $Y_i \to X$（$i = 1, 2$）是 Galois 群为 $G_i$ 的有限 \'etale
Galois 覆盖，则存在有限 \'etale Galois 覆盖 $Y \to X$，其 Galois 群是
$G_1 \times G_2$ 的子群。事实上，取相应的连续同态
$\pi_1(X, \overline{x}) \to G_i$，并令 $G$ 为诱导连续同态
$\pi_1(X, \overline{x}) \to G_1 \times G_2$ 的像。









\section{基本群的拓扑不变性}
\label{pione-section-topological-invariance}

\noindent
本节的主要结果是，连通概形的泛同胚会诱导基本群之间的同构。
见命题 \ref{pione-proposition-universal-homeomorphism}。

\medskip\noindent
我们通常不直接证明两个概形具有相同的基本群，而是证明它们的有限 \'etale
覆盖范畴相同。当然，只要它们连通，这就蕴含它们的基本群相同。

\begin{lemma}
\label{pione-lemma-what-equivalence-gives}
设 $f : X \to Y$ 是拟紧拟分离概形的态射，使得基变换函子
$\textit{F\'Et}_Y \to \textit{F\'Et}_X$ 是范畴等价。此时
\begin{enumerate}
\item $f$ 诱导同胚 $\pi_0(X) \to \pi_0(Y)$；
\item 若 $X$ 连通，等价地，若 $Y$ 连通，则
$\pi_1(X, \overline{x}) = \pi_1(Y, \overline{y})$。
\end{enumerate}
\end{lemma}

\begin{proof}
设 $Y = Y_0 \amalg Y_1$ 为到两个非空开闭子概形的不交并分解。我们断言
$f(X)$ 与两个 $Y_i$ 都相交。事实上，若非如此，不妨设
$f(X) \subset Y_1$，则可考虑有限 \'etale 态射 $V = Y_1 \to Y$。
它不是同构，但 $V \times_Y X \to X$ 是同构，矛盾。

\medskip\noindent
设 $X = X_0 \amalg X_1$ 为到两个开闭子概形的不交并分解。考虑有限
\'etale 态射 $U = X_1 \to X$。于是有 $U = X \times_Y V$，其中
$V \to Y$ 是某个有限 \'etale 态射。态射 $V \to Y$ 的次数局部常值，
故可得如下到开闭子概形的分解
$Y = \coprod_{d \geq 0} Y_d$，使得 $V \to Y$ 的次数为 $d$ 的部分恰为
$Y_d$。因为 $f^{-1}(Y_d) = \emptyset$ 在 $d > 1$ 时成立，由上面的论证
可知，$Y_d = \emptyset$ 在 $d > 1$ 时成立。进而可知，
$f^{-1}(Y_i) = X_i$ 对 $i = 0, 1$ 成立。

\medskip\noindent
于是 $f^{-1}$ 在 $Y$ 的开闭子集的集合与 $X$ 的开闭子集的集合之间诱导
双射。注意 $X$ 和 $Y$ 都是谱空间，见《性质》引理
\ref{properties-lemma-quasi-compact-quasi-separated-spectral}。
由《拓扑》引理 \ref{topology-lemma-connected-component-intersection}，
谱空间的开闭子集格决定其连通分支的集合。因此
$\pi_0(X) \to \pi_0(Y)$ 是双射。由于 $\pi_0(X)$ 和 $\pi_0(Y)$ 都是
profinite 空间（《拓扑》引理 \ref{topology-lemma-pi0-profinite}），
由《拓扑》引理 \ref{topology-lemma-bijective-map} 可知
$\pi_0(X) \to \pi_0(Y)$ 是同胚。这证明了 (1)。断言 (2) 是立即的。
\end{proof}

\noindent
下一个引理说明，Hensel 对的基本群就是相应闭子集的基本群。

\begin{lemma}
\label{pione-lemma-gabber}
设 $(A, I)$ 是 Hensel 对。置 $X = \Spec(A)$ 且 $Z = \Spec(A/I)$。函子
$$
\textit{F\'Et}_X \longrightarrow \textit{F\'Et}_Z,\quad
U \longmapsto U \times_X Z
$$
是范畴等价。
\end{lemma}

\begin{proof}
这只是《代数补篇》引理
\ref{more-algebra-lemma-finite-etale-equivalence} 的几何转述。
\end{proof}

\noindent
下一个引理说明，增厚的基本群与原概形的基本群相同。我们将在下面关于泛同胚
的较强命题的证明中使用这一点。

\begin{lemma}
\label{pione-lemma-thickening}
设 $X \subset X'$ 是概形的增厚。函子
$$
\textit{F\'Et}_{X'} \longrightarrow \textit{F\'Et}_X,\quad
U' \longmapsto U' \times_{X'} X
$$
是范畴等价。
\end{lemma}

\begin{proof}
关于增厚的讨论见《态射补篇》第
\ref{more-morphisms-section-thickenings} 节。设 $U' \to X'$ 是
\'etale 态射，并且 $U = U' \times_{X'} X \to X$ 是有限 \'etale
态射。则 $U' \to X'$ 也是有限 \'etale 态射。例如，这可由《态射补篇》
引理 \ref{more-morphisms-lemma-properties-that-extend-over-thickenings}
得到。现在，若 $X \subset X'$ 是有限阶增厚，则结合这一观察与《\'Etale
态射》定理 \ref{etale-theorem-remarkable-equivalence} 即可证明本引理。
下面我们将对一般的增厚证明本引理，不过建议读者跳过这个证明。

\medskip\noindent
设 $X' = \bigcup X_i'$ 是仿射开覆盖。置
$X_i = X \times_{X'} X_i'$，$X_{ij}' = X'_i \cap X'_j$，
$X_{ij} = X \times_{X'} X_{ij}'$, $X_{ijk}' = X'_i \cap X'_j \cap X'_k$,
$X_{ijk} = X \times_{X'} X_{ijk}'$。
假设对每个增厚
$X_i \subset X'_i$、$X_{ij} \subset X_{ij}'$ 和
$X_{ijk} \subset X_{ijk}'$ 都能证明该定理。则由概形的相对粘合，可得
$X \subset X'$ 的结论，见《构造》第
\ref{constructions-section-relative-glueing} 节。注意，概形
$X_i'$、$X_{ij}'$、$X_{ijk}'$ 作为仿射概形的开子概形，都是分离的。
把上述论证再重复一次，便可归约到概形
$X'_i$、$X_{ij}'$、$X_{ijk}'$ 都仿射的情形。

\medskip\noindent
在仿射情形下，有 $X' = \Spec(A')$ 和 $X = \Spec(A'/I')$，其中
$I'$ 是局部幂零理想。于是 $(A', I')$ 是 Hensel 对（《代数补篇》引理
\ref{more-algebra-lemma-locally-nilpotent-henselian}），结论便由引理
\ref{pione-lemma-gabber} 得出（在此情形下，该引理的证明容易得多）。
\end{proof}

\noindent
证明下一个命题的“正确”方式，是从 \'etale 景的不变性推出它，见《\'Etale
上同调》定理 \ref{etale-cohomology-theorem-topological-invariance}。

\begin{proposition}
\label{pione-proposition-universal-homeomorphism}
设 $f : X \to Y$ 是概形的泛同胚。则
$$
\textit{F\'Et}_Y \longrightarrow \textit{F\'Et}_X,\quad
V \longmapsto V \times_Y X
$$
是范畴等价。因此，若 $X$ 和 $Y$ 连通，则 $f$ 诱导基本群的同构
$\pi_1(X, \overline{x}) \to \pi_1(Y, \overline{y})$。
\end{proposition}

\begin{proof}
回忆，泛同胚恰好就是整、泛单射且满射的态射，见《态射》引理
\ref{morphisms-lemma-universal-homeomorphism}。特别地，由《态射》引理
\ref{morphisms-lemma-universally-injective}，对角态射
$\Delta : X \to X \times_Y X$ 是增厚。因此，由引理
\ref{pione-lemma-thickening} 可知，给定有限 \'etale 态射 $U \to X$，存在唯一的同构
$$
\varphi : U \times_Y X \to X \times_Y U
$$
；这是 $X \times_Y X$ 上两个有限 \'etale 概形之间的同构，并且沿 $\Delta$ 拉回后，
成为恒等态射 $\text{id} : U \to U$；该态射定义在 $X$ 上。由于
$X \to X \times_Y X \times_Y X$ 也是增厚（它是双射闭浸入），可知
$(U, \varphi)$ 是相对于 $X/Y$ 的下降资料。由《\'Etale 态射》命题
\ref{etale-proposition-effective} 可知，$U = X \times_Y V$，其中
$V \to Y$ 拟紧、分离且 \'etale。我们略去 $V \to Y$ 有限的证明（提示：
态射 $U \to V$ 是满射，而 $U \to Y$ 是整态射）。因此
$\textit{F\'Et}_Y \to \textit{F\'Et}_X$ 是本质满的。

\medskip\noindent
与上面同样论证可知，给定 $V_1 \to Y$ 和 $V_2 \to Y$ 这两个
$\textit{F\'Et}_Y$ 中的对象，任意态射
$a : X \times_Y V_1 \to X \times_Y V_2$（在 $X$ 上）都与典范下降资料相容。因此，由
《\'Etale 态射》引理 \ref{etale-lemma-fully-faithful-cases}，$a$ 下降为
态射 $V_1 \to V_2$；该态射定义在 $Y$ 上。
\end{proof}










\section{固有概形的有限 \'etale 覆盖}
\label{pione-section-finite-etale-over-proper}

\noindent
本节将证明，Hensel 局部环上的连通固有概形与其特殊纤维具有相同的基本群。
我们还将证明该结果在 Hensel 对情形下的一个变体。

\medskip\noindent
我们还将证明，代数闭域 $k$ 上的连通固有概形的基本群在把 $k$ 换成
一个代数闭域扩张后保持不变。

\medskip\noindent
我们不在连通情形下陈述并证明这些结果，而是在一般情形下证明它们，
并把利用引理 \ref{pione-lemma-what-equivalence-gives} 推出基本群结论的工作
留给读者。

\begin{lemma}
\label{pione-lemma-finite-etale-on-proper-over-henselian}
设 $A$ 是 Hensel 局部环，$X$ 是 $A$ 上的固有概形，其闭纤维为 $X_0$。
则函子
$$
\textit{F\'Et}_X \to \textit{F\'Et}_{X_0},\quad
U \longmapsto U_0 = U \times_X X_0
$$
是范畴等价。
\end{lemma}

\begin{proof}
这里给出的证明是应用代数化与逼近的一个例子。我们分若干步进行。

\medskip\noindent
当 $A$ 是完备 Noether 局部环时的本质满性。令
$X_n = X \times_{\Spec(A)} \Spec(A/\mathfrak m^{n + 1})$。
由《\'Etale 态射》定理 \ref{etale-theorem-remarkable-equivalence}，嵌入
$$
X_0 \to X_1 \to X_2 \to \ldots
$$
在 $X_0$ 上的 \'etale 概形范畴与在 $X_n$ 上的 \'etale 概形范畴之间
诱导范畴等价。此外，若 $U_n \to X_n$ 对应于有限 \'etale 态射
$U_0 \to X_0$，则 $U_n \to X_n$ 也是有限的；例如，这可由《态射补篇》
引理 \ref{more-morphisms-lemma-thicken-property-morphisms-cartesian} 得到。
在这种情形下，态射 $U_0 \to \Spec(A/\mathfrak m)$ 是固有的，因为
$X_0$ 在 $A/\mathfrak m$ 上固有。因此可应用 Grothendieck
代数化定理（采用《概形上同调》引理
\ref{coherent-lemma-algebraize-formal-scheme-finite-over-proper} 的形式），
得到有限态射 $U \to X$，其在 $X_0$ 上的限制还原为 $U_0$。由《态射补篇》
引理 \ref{more-morphisms-lemma-check-smoothness-on-infinitesimal-nbhds}
可知，$U \to X$ 在 $U_0$ 的每一点处都是 \'etale 的。然而，因为 $U$
的每一点都特化到 $U_0$ 的某一点（这是因为 $U$ 在 $A$ 上固有），
可知 $U \to X$ 是 \'etale 态射。由此，该函子本质满。

\medskip\noindent
当 $A$ 是完备 Noether 局部环时的全忠实性。设 $U \to X$ 和 $V \to X$
是有限 \'etale 态射，且 $\varphi_0 : U_0 \to V_0$ 是 $X_0$ 上的态射。
考虑态射
$$
\Gamma_{\varphi_0} : U_0 \longrightarrow U_0 \times_{X_0} V_0
$$
它既是有限 \'etale 态射，也是闭浸入。把本质满性应用于
$X = U \times_X V$，可得有限 \'etale 态射 $W \to U \times_X V$，
其特殊纤维同构于 $\Gamma_{\varphi_0}$。考虑投影 $W \to U$。依构造，
它是有限 \'etale 态射，并且在 $U_0$ 上是同构。由《\'Etale 态射》引理
\ref{etale-lemma-finite-etale-one-point}，$W \to U$ 在 $U_0$ 的一个开邻域
上是同构。因此它处处是同构，而复合 $\varphi : U \cong W \to V$
正是所需的 $\varphi_0$ 的提升。

\medskip\noindent
当 $A$ 是 Hensel Noether 局部 G-环时的本质满性。设 $U_0 \to X_0$
是有限 \'etale 态射。令 $A^\wedge$ 为 $A$ 关于极大理想的完备化，
并令 $X^\wedge$ 为 $X$ 到 $A^\wedge$ 的基变换。由上面的结果，存在
有限 \'etale 态射 $V \to X^\wedge$，其特殊纤维为 $U_0$。写成
$A^\wedge = \colim A_i$，其中 $A \to A_i$ 为有限型。由《极限》引理
\ref{limits-lemma-descend-finite-presentation}，存在某个 $i$ 以及有限表现态射
$U_i \to X_{A_i}$，其到 $X^\wedge$ 的基变换为 $V$。增大 $i$ 后，
可设 $U_i \to X_{A_i}$ 有限且 \'etale（《极限》引理
\ref{limits-lemma-descend-finite-finite-presentation} 和
\ref{limits-lemma-descend-etale}）。写
$$
A_i = A[x_1, \ldots, x_n]/(f_1, \ldots, f_m)
$$
则环同态 $A_i \to A^\wedge$ 可重新解释为一个解
$(a_1, \ldots, a_n)$，它是 $A^\wedge$ 中方程组 $f_j = 0$ 的解。
由《环同态的光滑化》定理
\ref{smoothing-theorem-approximation-property}，可把这个解逼近到例如
$11$ 阶；所得解为 $(b_1, \ldots, b_n)$，位于 $A$ 中。把它重新翻译成
环同态，便得到 $A$-代数同态 $A_i \to A$；它与原同态
$A_i \to A^\wedge$ 给出同一个闭点（因为 $11 > 1$）。因此，
$U \to X$，即沿这个环同态对 $V \to X_{A_i}$ 作的基变换，是有限 \'etale 态射，
其特殊纤维同构于 $U_0$。

\medskip\noindent
当 $A$ 是 Hensel Noether 局部 G-环时的全忠实性。它可由本质满性推出，
论证与 $A$ 是完备 Noether 环时完全相同。

\medskip\noindent
一般情形。设 $(A, \mathfrak m)$ 是 Hensel 局部环。置 $S = \Spec(A)$，
并以 $s \in S$ 表示闭点。由《极限》引理
\ref{limits-lemma-proper-limit-of-proper-finite-presentation-noetherian}，
可把 $X \to \Spec(A)$ 写成固有态射 $X_i \to S_i$ 的余滤过极限，其中
$S_i$ 在 $\mathbf{Z}$ 上为有限型。对每个 $i$，令 $s_i \in S_i$ 为 $s$
的像。由于 $S = \lim S_i$ 且 $A = \mathcal{O}_{S, s}$，有
$A = \colim \mathcal{O}_{S_i, s_i}$。环 $A_i = \mathcal{O}_{S_i, s_i}$
是 Noether 局部 G-环（《代数补篇》命题
\ref{more-algebra-proposition-ubiquity-G-ring}）。由《代数补篇》引理
\ref{more-algebra-lemma-henselization-colimit} 可知
$A = \colim A_i^h$。由《代数补篇》引理
\ref{more-algebra-lemma-henselization-G-ring}，各环 $A_i^h$ 都是 G-环。
因此 $A = \colim A_i^h$，并且作为概形有
$$
X = \lim (X_i \times_{S_i} \Spec(A_i^h))
$$
。$X$ 上的有限 \'etale 概形范畴，是
$X_i \times_{S_i} \Spec(A_i^h)$ 上的有限 \'etale 概形范畴的极限（由《极限》引理
\ref{limits-lemma-descend-finite-presentation}、
\ref{limits-lemma-descend-finite-finite-presentation} 和
\ref{limits-lemma-descend-etale})
。同样的结论也适用于
$X_0 = \lim (X_i \times_{S_i} s_i)$ 上的有限 \'etale 概形。因此，
$X / \Spec(A)$ 的结果可从上面已经处理过的
$(X_i \times_{S_i} \Spec(A_i^h)) / \Spec(A_i^h)$ 的结果形式地推出。
\end{proof}

\begin{lemma}
\label{pione-lemma-finite-etale-on-proper-over-henselian-pair}
设 $(A, I)$ 是 Hensel 对，$X$ 是 $A$ 上的固有概形。置
$X_0 = X \times_{\Spec(A)} \Spec(A/I)$。则函子
$$
\textit{F\'Et}_X \to \textit{F\'Et}_{X_0},\quad
U \longmapsto U_0 = U \times_X X_0
$$
是范畴等价。
\end{lemma}

\begin{proof}
本引理的证明与引理
\ref{pione-lemma-finite-etale-on-proper-over-henselian} 的证明完全相同。

\medskip\noindent
当 $A$ 是 Noether 环且 $I$-进完备时的本质满性。令
$X_n = X \times_{\Spec(A)} \Spec(A/I^{n + 1})$。由《\'Etale 态射》定理
\ref{etale-theorem-remarkable-equivalence}，嵌入
$$
X_0 \to X_1 \to X_2 \to \ldots
$$
在 $X_0$ 上的 \'etale 概形范畴与在 $X_n$ 上的 \'etale 概形范畴之间
诱导范畴等价。此外，若 $U_n \to X_n$ 对应于有限 \'etale 态射
$U_0 \to X_0$，则 $U_n \to X_n$ 也是有限的；例如，这可由《态射补篇》
引理 \ref{more-morphisms-lemma-thicken-property-morphisms-cartesian} 得到。
在这种情形下，态射 $U_0 \to \Spec(A/I)$ 是固有的，因为
$X_0$ 在 $A/I$ 上固有。因此可应用 Grothendieck 代数化定理
（采用《概形上同调》引理
\ref{coherent-lemma-algebraize-formal-scheme-finite-over-proper} 的形式），
得到有限态射 $U \to X$，其在 $X_0$ 上的限制还原为 $U_0$。由《态射补篇》
引理 \ref{more-morphisms-lemma-check-smoothness-on-infinitesimal-nbhds}
可知，$U \to X$ 在 $U_0$ 的每一点处都是 \'etale 的。然而，因为 $U$
的每一点都特化到 $U_0$ 的某一点（这是因为 $U$ 在 $A$ 上固有），
可知 $U \to X$ 是 \'etale 态射。由此，该函子本质满。

\medskip\noindent
当 $A$ 是 Noether 环且 $I$-进完备时的全忠实性。设 $U \to X$ 和
$V \to X$ 是有限 \'etale 态射，且 $\varphi_0 : U_0 \to V_0$ 是
$X_0$ 上的态射。考虑态射
$$
\Gamma_{\varphi_0} : U_0 \longrightarrow U_0 \times_{X_0} V_0
$$
它既是有限 \'etale 态射，也是闭浸入。把本质满性应用于
$X = U \times_X V$，可得有限 \'etale 态射 $W \to U \times_X V$，
其特殊纤维同构于 $\Gamma_{\varphi_0}$。考虑投影 $W \to U$。依构造，
它是有限 \'etale 态射，并且在 $U_0$ 上是同构。由《\'Etale 态射》引理
\ref{etale-lemma-finite-etale-one-point}，$W \to U$ 在 $U_0$ 的一个开邻域
上是同构。因此它处处是同构，而复合 $\varphi : U \cong W \to V$
正是所需的 $\varphi_0$ 的提升。

\medskip\noindent
当 $(A, I)$ 是 Hensel 对且 $A$ 是 Noether G-环时的本质满性。
设 $U_0 \to X_0$ 是有限 \'etale 态射。令 $A^\wedge$ 为 $A$ 关于 $I$
的完备化。注意，$A^\wedge$ 是 $IA^\wedge$-进完备的 Noether 环，
见《代数》引理 \ref{algebra-lemma-completion-complete} 和
\ref{algebra-lemma-completion-Noetherian-Noetherian}。令 $X^\wedge$
为 $X$ 到 $A^\wedge$ 的基变换。由上面的结果，存在有限 \'etale 态射
$V \to X^\wedge$，其特殊纤维为 $U_0$。写成
$A^\wedge = \colim A_i$，其中 $A \to A_i$ 为有限型。由《极限》引理
\ref{limits-lemma-descend-finite-presentation}，存在某个 $i$ 以及有限表现态射
$U_i \to X_{A_i}$，其到 $X^\wedge$ 的基变换为 $V$。增大 $i$ 后，
可设 $U_i \to X_{A_i}$ 有限且 \'etale（《极限》引理
\ref{limits-lemma-descend-finite-finite-presentation} 和
\ref{limits-lemma-descend-etale}）。写
$$
A_i = A[x_1, \ldots, x_n]/(f_1, \ldots, f_m)
$$
则环同态 $A_i \to A^\wedge$ 可重新解释为一个解
$(a_1, \ldots, a_n)$，它是 $A^\wedge$ 中方程组 $f_j = 0$ 的解。
由《环同态的光滑化》引理
\ref{smoothing-lemma-henselian-pair}，可把这个解逼近到例如 $11$ 阶；
所得解为 $(b_1, \ldots, b_n)$，位于 $A$ 中。把它重新翻译成
环同态，便得到 $A$-代数同态 $A_i \to A$；它与原同态
$A_i \to A^\wedge$ 给出同一个闭点（因为 $11 > 1$）。因此，
$U \to X$，即沿这个环同态对 $V \to X_{A_i}$ 作的基变换，是有限 \'etale 态射，
其特殊纤维同构于 $U_0$。

\medskip\noindent
当 $(A, I$ 是 Hensel 对且 $A$ 是 Noether G-环时的全忠实性。它可由
本质满性推出，论证与 $A$ 是完备 Noether 环时完全相同。

\medskip\noindent
一般情形。设 $(A, I)$ 是 Hensel 对。置 $S = \Spec(A)$，并记
$S_0 = \Spec(A/I)$。由《极限》引理
\ref{limits-lemma-proper-limit-of-proper-finite-presentation-noetherian}，
可把 $X \to \Spec(A)$ 写成固有态射 $X_i \to S_i$ 的余滤过极限，其中
$S_i$ 仿射且在 $\mathbf{Z}$ 上为有限型。写 $S_i = \Spec(A_i)$，并以
$I_i \subset A_i$ 表示理想 $I$ 沿同态 $A_i \to A$ 的逆像。置
$S_{i, 0} = \Spec(A_i/I_i)$。由于 $S = \lim S_i$，有
$A = \colim A_i$。因而还有 $I = \colim I_i$ 及 $A/I = \colim A_i/I_i$。
环 $A_i$ 是 Noether G-环（《代数补篇》命题
\ref{more-algebra-proposition-ubiquity-G-ring}）。以 $(A_i^h, I_i^h)$
表示对 $(A_i, I_i)$ 的 Hensel 化。由《代数补篇》引理
\ref{more-algebra-lemma-henselization-colimit} 可知
$A = \colim A_i^h$。由《代数补篇》引理
\ref{more-algebra-lemma-henselization-pair-G-ring}，各环 $A_i^h$ 都是
G-环。因此 $A = \colim A_i^h$，并且作为概形有
$$
X = \lim (X_i \times_{S_i} \Spec(A_i^h))
$$
。$X$ 上的有限 \'etale 概形范畴，是
$X_i \times_{S_i} \Spec(A_i^h)$ 上的有限 \'etale 概形范畴的极限（由《极限》引理
\ref{limits-lemma-descend-finite-presentation}、
\ref{limits-lemma-descend-finite-finite-presentation} 和
\ref{limits-lemma-descend-etale})
。同样的结论也适用于
$X_0 = \lim (X_i \times_{S_i} S_{i, 0})$ 上的有限 \'etale 概形。因此，
$X / \Spec(A)$ 的结果可从上面已经处理过的
$(X_i \times_{S_i} \Spec(A_i^h)) / \Spec(A_i^h)$ 的结果形式地推出。
\end{proof}

\begin{lemma}
\label{pione-lemma-finite-etale-invariant-over-proper}
设 $k'/k$ 是代数闭域的扩张，$X$ 是 $k$ 上的固有概形。则函子
$$
U \longmapsto U_{k'}
$$
在 $X$ 上的有限 \'etale 概形范畴与在 $X_{k'}$ 上的有限 \'etale
概形范畴之间给出范畴等价。
\end{lemma}

\begin{proof}
我们来证明该函子本质满。设 $U' \to X_{k'}$ 是有限 \'etale 态射。
把 $k' = \colim A_i$ 写成有限型 $k$-代数的滤过余极限。由《极限》引理
\ref{limits-lemma-descend-finite-presentation}，存在某个 $i$ 以及有限表现态射
$U_i \to X_{A_i}$，其到 $X_{k'}$ 的基变换为 $U'$。增大 $i$ 后，
可设 $U_i \to X_{A_i}$ 有限且 \'etale（《极限》引理
\ref{limits-lemma-descend-finite-finite-presentation} 和
\ref{limits-lemma-descend-etale}）。
由于 $k$ 代数闭，可找到一个 $k$-值点 $t$，它位于 $\Spec(A_i)$ 中。令
$U = (U_i)_t$ 为 $U_i$ 在 $t$ 上的纤维。令 $A_i^h$ 为
$(A_i)_{\mathfrak m}$ 的 Hensel 化，其中 $\mathfrak m$ 是与点 $t$
对应的极大理想。由引理
\ref{pione-lemma-finite-etale-on-proper-over-henselian} 可知，
$(U_i)_{A_i^h} = U \times \Spec(A_i^h)$；这是 $X_{A_i^h}$ 上的概形同构。
现在，由于 $A_i^h$ 在 $A_i$
上是代数的（例如见《环同态的光滑化》例
\ref{smoothing-example-describe-henselian} 中的讨论），而且 $k'$ 代数闭，
可以找到环同态 $A_i^h \to k'$，它扩张给定的包含
$A_i \subset k'$。因此 $U'$ 同构于 $U$ 的基变换。
全忠实性的证明完全相同。
\end{proof}








\section{局部连通性}
\label{pione-section-unibranch}

\noindent
本节研究映射 $\pi_1(U) \to \pi_1(X)$ 何时为满射，其中 $U$ 是概形
$X$ 的稠密开子概形。我们将看到，粗略地说，只要 $U \cap B$ 连通，
其中 $B$ 是点 $x \in X \setminus U$ 附近任意小的“球”，上述映射就是满射。

\begin{lemma}
\label{pione-lemma-dense-faithful}
设 $f : X \to Y$ 是概形态射。若 $f(X)$ 在 $Y$ 中稠密，则基变换函子
$\textit{F\'Et}_Y \to \textit{F\'Et}_X$ 忠实。
\end{lemma}

\begin{proof}
因为有限 \'etale 覆盖的范畴具有内部 Hom（引理
\ref{pione-lemma-internal-hom-finite-etale}），只需证明以下命题：给定 $W$
（它在 $Y$ 上有限 \'etale），以及态射 $s : X \to W$（在 $X$ 上），
截面 $t : Y \to W$ 若满足 $s = t \circ f$，则至多有一个。设有两个截面
$t_1, t_2 : Y \to W$，使得 $s = t_1 \circ f = t_2 \circ f$。
$t_1$ 与 $t_2$ 的等化子在 $Y$ 中闭（《概形》引理
\ref{schemes-lemma-where-are-they-equal}），而 $f(X)$ 在 $Y$ 中稠密，
故 $t_1$ 与 $t_2$ 在 $Y_{red}$ 上相同。于是 $t_1$ 与 $t_2$ 具有相同的像；
该像是 $W$ 的开闭子概形，并同构地映到 $Y$（《\'Etale 态射》命题
\ref{etale-proposition-properties-sections}）。因此二者相等。
\end{proof}

\noindent
下一个引理假设严格 Hensel 化的穿孔谱连通。例如，局部环几何单枝时，
这一条件成立，见《代数补篇》引理
\ref{more-algebra-lemma-geometrically-unibranch}。《性质》引理
\ref{properties-lemma-geometrically-unibranch} 给出了一个部分逆命题。

\begin{lemma}
\label{pione-lemma-same-etale-extensions}
设 $(A, \mathfrak m)$ 是局部环。置 $X = \Spec(A)$，并令
$U = X \setminus \{\mathfrak m\}$。若 $A$ 的严格 Hensel 化的穿孔谱连通，则
$$
\textit{F\'Et}_X \longrightarrow \textit{F\'Et}_U,\quad
Y \longmapsto Y \times_X U
$$
是全忠实函子。
\end{lemma}

\begin{proof}
先设 $A$ 严格 Hensel。在这种情形下，有限 \'etale 覆盖 $Y$（在 $X$
上）都同构于有限多个 $X$ 的不交并。因此，只需证明：态射
$U \to U \amalg \ldots \amalg U$（在 $U$ 上）唯一延拓为态射
$X \to X \amalg \ldots \amalg X$（在 $X$ 上）。若 $U$ 连通（特别地，非空），
这显然成立。

\medskip\noindent
一般情形。因为有限 \'etale 覆盖的范畴具有内部 Hom（引理
\ref{pione-lemma-internal-hom-finite-etale}），只需证明：给定 $Y$（它在 $X$ 上有限
\'etale），任意态射 $s : U \to Y$（在 $X$ 上）都延拓为态射
$t : X \to Y$（在 $X$ 上）。令 $A^{sh}$ 为 $A$ 的严格 Hensel 化，并记
$X^{sh} = \Spec(A^{sh})$、$U^{sh} = U \times_X X^{sh}$、
$Y^{sh} = Y \times_X X^{sh}$。由第一段及对 $A$ 的假设，可把基变换
$s^{sh} : U^{sh} \to Y^{sh}$（它来自 $s$）延拓为
$t^{sh} : X^{sh} \to Y^{sh}$。置 $A' = A^{sh} \otimes_A A^{sh}$。
于是，两个拉回 $t'_1, t'_2$（来自 $t^{sh}$，拉到 $X' = \Spec(A')$）
都是拉回 $s'$（来自 $s$，拉到 $U' = U \times_X X'$）的延拓。由于
$A \to A'$ 平坦，$U' \subset X'$ 在拓扑意义下稠密；这是由
$A \to A'$ 的降链性质（《代数》引理
\ref{algebra-lemma-flat-going-down}）得到的。因此，由引理
\ref{pione-lemma-dense-faithful}，$t'_1 = t'_2$。
所以 $t^{sh}$ 下降为态射 $t : X \to Y$；例如，这可由《下降》引理
\ref{descent-lemma-fpqc-universal-effective-epimorphisms} 得到。
\end{proof}

\noindent
鉴于引理 \ref{pione-lemma-same-etale-extensions}，自然要问：一个环（及其严格
Hensel 化）的穿孔谱何时连通。Hartshorne 的一个著名引理给出了充分条件，
见《局部上同调》引理
\ref{local-cohomology-lemma-depth-2-connected-punctured-spectrum}。

\begin{lemma}
\label{pione-lemma-quasi-compact-dense-open-connected-at-infinity-Noetherian}
设 $X$ 是概形，$U \subset X$ 是稠密开子概形。假设
\begin{enumerate}
\item $X$ 的底拓扑空间是 Noether 空间；
\item 对每个 $x \in X \setminus U$，$\mathcal{O}_{X, x}$ 的严格 Hensel
化的穿孔谱都连通。
\end{enumerate}
则 $\textit{F\'Et}_X \to \textit{F\'et}_U$ 全忠实。
\end{lemma}

\begin{proof}
设 $Y_1, Y_2$ 在 $X$ 上有限 \'etale，并设
$\varphi : (Y_1)_U \to (Y_2)_U$ 是 $U$ 上的态射。须证明 $\varphi$
唯一提升为态射 $Y_1 \to Y_2$（在 $X$ 上）。唯一性由引理
\ref{pione-lemma-dense-faithful} 得出。

\medskip\noindent
设 $x \in X \setminus U$ 是 $X \setminus U$ 的一个不可约分支的泛点。置
$V = U \times_X \Spec(\mathcal{O}_{X, x})$。由 $x$ 的选取方式，它就是
$\Spec(\mathcal{O}_{X, x})$ 的穿孔谱。由引理
\ref{pione-lemma-same-etale-extensions}，态射
$\varphi_V : (Y_1)_V \to (Y_2)_V$ 唯一延拓为态射
$(Y_1)_{\Spec(\mathcal{O}_{X, x})} \to (Y_2)_{\Spec(\mathcal{O}_{X, x})}$。
由《极限》引理 \ref{limits-lemma-glueing-near-point}，可找到开子概形
$U \subset U'$（它包含 $x$），以及态射
$\varphi' : (Y_1)_{U'} \to (Y_2)_{U'}$，它延拓 $\varphi$。由于 $X$ 的底拓扑空间是
Noether 空间，对 $\varphi$ 已定义的开子概形的补集作 Noether 归纳，
便完成证明。
\end{proof}

\begin{lemma}
\label{pione-lemma-retrocompact-dense-open-connected-at-infinity-closed}
设 $X$ 是概形，$U \subset X$ 是稠密开子概形。假设
\begin{enumerate}
\item $U \to X$ 拟紧；
\item $X \setminus U$ 的每一点都是闭点；
\item 对每个 $x \in X \setminus U$，$\mathcal{O}_{X, x}$ 的严格 Hensel
化的穿孔谱都连通。
\end{enumerate}
则 $\textit{F\'Et}_X \to \textit{F\'et}_U$ 全忠实。
\end{lemma}

\begin{proof}
设 $Y_1, Y_2$ 在 $X$ 上有限 \'etale，并设
$\varphi : (Y_1)_U \to (Y_2)_U$ 是 $U$ 上的态射。须证明 $\varphi$
唯一提升为态射 $Y_1 \to Y_2$（在 $X$ 上）。唯一性由引理
\ref{pione-lemma-dense-faithful} 得出。

\medskip\noindent
设 $x \in X \setminus U$。置 $V = U \times_X \Spec(\mathcal{O}_{X, x})$。
由于 $X \setminus U$ 的每一点都闭，$V$ 就是
$\Spec(\mathcal{O}_{X, x})$ 的穿孔谱。由引理
\ref{pione-lemma-same-etale-extensions}，态射
$\varphi_V : (Y_1)_V \to (Y_2)_V$ 唯一延拓为态射
$(Y_1)_{\Spec(\mathcal{O}_{X, x})} \to (Y_2)_{\Spec(\mathcal{O}_{X, x})}$。
由《极限》引理 \ref{limits-lemma-glueing-near-point}（这里用到 $U$
在 $X$ 中逆紧），可找到开子概形 $U \subset U'_x$（它包含 $x$），
以及态射 $\varphi'_x : (Y_1)_{U'_x} \to (Y_2)_{U'_x}$，它延拓
$\varphi$。注意，对任意两点
$x, x' \in X \setminus U$，态射 $\varphi'_x$ 与 $\varphi'_{x'}$
在 $U'_x \cap U'_{x'}$ 上相同，因为 $U$ 在该开子概形中稠密（引理
\ref{pione-lemma-dense-faithful}）。因此可按所需把 $\varphi$ 延拓到
$\bigcup U'_x = X$。
\end{proof}

\begin{lemma}
\label{pione-lemma-quasi-compact-dense-open-connected-at-infinity}
设 $X$ 是概形，$U \subset X$ 是稠密开子概形。假设
\begin{enumerate}
\item $X$ 的每个拟紧开子概形只有有限多个不可约分支；
\item 对每个 $x \in X \setminus U$，$\mathcal{O}_{X, x}$ 的严格 Hensel
化的穿孔谱都连通。
\end{enumerate}
则 $\textit{F\'Et}_X \to \textit{F\'et}_U$ 全忠实。
\end{lemma}

\begin{proof}
设 $Y_1, Y_2$ 在 $X$ 上有限 \'etale，并设
$\varphi : (Y_1)_U \to (Y_2)_U$ 是 $U$ 上的态射。须证明 $\varphi$
唯一提升为态射 $Y_1 \to Y_2$（在 $X$ 上）。唯一性由引理
\ref{pione-lemma-dense-faithful} 得出。为证明存在性，我们将说明可扩大 $U$
（若 $U \not = X$），再用 Zorn 引理完成证明。

\medskip\noindent
设 $x \in X \setminus U$ 是 $X \setminus U$ 的一个不可约分支的泛点。置
$V = U \times_X \Spec(\mathcal{O}_{X, x})$。由 $x$ 的选取方式，它就是
$\Spec(\mathcal{O}_{X, x})$ 的穿孔谱。由引理
\ref{pione-lemma-same-etale-extensions}，态射
$\varphi_V : (Y_1)_V \to (Y_2)_V$（唯一）延拓为态射
$(Y_1)_{\Spec(\mathcal{O}_{X, x})} \to (Y_2)_{\Spec(\mathcal{O}_{X, x})}$。
选取仿射邻域 $W \subset X$（它包含 $x$）。因为 $U \cap W$ 在 $W$ 中稠密，
它包含各泛点 $\eta_1, \ldots, \eta_n$，即 $W$ 的各泛点。选取仿射开子概形
$W' \subset W \cap U$，使其包含 $\eta_1, \ldots, \eta_n$。置
$V' = W' \times_X \Spec(\mathcal{O}_{X, x})$。把《极限》引理
\ref{limits-lemma-glueing-near-point} 应用于 $x \in W \supset W'$，
可找到开子概形 $W' \subset W'' \subset W$，其中 $x \in W''$，
以及态射 $\varphi'' : (Y_1)_{W''} \to (Y_2)_{W''}$，它与
$\varphi$ 在 $W'$ 上相同。由于 $W'$ 在 $W'' \cap U$ 中稠密，由引理
\ref{pione-lemma-dense-faithful} 可知，$\varphi$ 与 $\varphi''$ 在
$U \cap W'$ 上相同。因此 $\varphi$ 与 $\varphi''$ 粘合为态射
$\varphi'$，它定义在 $U' = U \cup W''$ 上，并与 $\varphi$ 在 $U$
相同。注意 $x \in U'$，所以我们已把 $\varphi$ 延拓到一个严格更大的开子概形。

\medskip\noindent
考虑集合 $\mathcal{S}$，其元素是二元组 $(U', \varphi')$，满足
$U \subset U'$，且 $\varphi'$ 是 $\varphi$ 的延拓。按显然的方式在
$\mathcal{S}$ 上赋予偏序。若 $(U'_i, \varphi'_i)$ 是全序子集，则它有
最大元 $(U', \varphi')$：只需取 $U' = \bigcup U'_i$，并令
$\varphi' : (Y_1)_{U'} \to (Y_2)_{U'}$ 为与 $\varphi'_i$ 在
$U'_i$ 上相同的态射。因此可应用 Zorn 引理，$\mathcal{S}$ 有极大元。
由上面的论证，这个极大元就是 $\varphi$ 在整个 $X$ 上的延拓。
\end{proof}

\begin{lemma}
\label{pione-lemma-local-exact-sequence}
设 $(A, \mathfrak m)$ 是局部环。置 $X = \Spec(A)$，并令
$U = X \setminus \{\mathfrak m\}$。令 $U^{sh}$ 为严格 Hensel 化
$A^{sh}$（即 $A$ 的严格 Hensel 化）的穿孔谱。假设 $U$ 拟紧且 $U^{sh}$
连通。则序列
$$
\pi_1(U^{sh}, \overline{u}) \to \pi_1(U, \overline{u}) \to
\pi_1(X, \overline{u}) \to 1
$$
按引理 \ref{pione-lemma-functoriality-galois-ses} 第 (1) 部分的意义是正合的。
\end{lemma}

\begin{proof}
映射 $\pi_1(U) \to \pi_1(X)$ 由引理
\ref{pione-lemma-same-etale-extensions} 和
\ref{pione-lemma-functoriality-galois-surjective} 可知是满射。

\medskip\noindent
记 $X^{sh} = \Spec(A^{sh})$。设 $Y \to X$ 是有限 \'etale 态射。则
$Y^{sh} = Y \times_X X^{sh} \to X^{sh}$ 是有限 \'etale 态射。由于
$A^{sh}$ 严格 Hensel，$Y^{sh}$ 同构于有限多个 $X^{sh}$ 的不交并。因此
$Y \times_X U^{sh}$ 亦然。于是复合
$\pi_1(U^{sh}) \to \pi_1(U) \to \pi_1(X)$ 平凡，见引理
\ref{pione-lemma-composition-trivial}。

\medskip\noindent
为完成证明，按照引理 \ref{pione-lemma-functoriality-galois-ses}，只需证明：
给定有限 \'etale 态射 $V \to U$，若 $V \times_U U^{sh}$ 是若干个
$U^{sh}$ 的不交并，则可以找到有限 \'etale 态射 $Y \to X$，使得
$V \cong Y \times_X U$ 在 $U$ 上成立。该假设蕴含存在有限 \'etale 态射
$Y^{sh} \to X^{sh}$ 以及同构
$V \times_U U^{sh} \cong Y^{sh} \times_{X^{sh}} U^{sh}$。考虑图表
$$
\xymatrix{
U \ar[d] & U^{sh} \ar[d] \ar[l] &
U^{sh} \times_U U^{sh} \ar[d] \ar@<1ex>[l] \ar@<-1ex>[l] &
U^{sh} \times_U U^{sh} \times_U U^{sh}
\ar[d] \ar@<1ex>[l] \ar[l] \ar@<-1ex>[l] \\
X & X^{sh} \ar[l] &
X^{sh} \times_X X^{sh} \ar@<1ex>[l] \ar@<-1ex>[l] &
X^{sh} \times_X X^{sh} \times_X X^{sh} \ar@<1ex>[l] \ar[l] \ar@<-1ex>[l]
}
$$
由于按假设 $U \subset X$ 拟紧，所有向下的箭头都是拟紧开浸入。设
$\xi \in X^{sh} \times_X X^{sh}$ 是不属于
$U^{sh} \times_U U^{sh}$ 的点。则 $\xi$ 位于闭点 $x^{sh}$（在
$X^{sh}$ 中）上方。考虑局部环同态
$$
A^{sh} = \mathcal{O}_{X^{sh}, x^{sh}} \to
\mathcal{O}_{X^{sh} \times_X X^{sh}, \xi}
$$
它由第一投影 $X^{sh} \times_X X^{sh}$ 决定。这个同态是一些局部同态的
滤过余极限，而这些局部同态由 \'etale 环映射局部化而得。由于 $A^{sh}$
严格 Hensel，可知该同态是同构。由于对补集中的每个 $\xi$ 都是如此，
这些点之间不存在特殊化关系，因而每个这样的 $\xi$ 都是闭点（这也可以
直接证明）。点 $\xi$ 处的局部环同构于 $A^{sh}$，故它严格 Hensel，且其
穿孔谱连通。对于点 $\xi$（它属于
$X^{sh} \times_X X^{sh} \times_X X^{sh}$ 而不属于
$U^{sh} \times_U U^{sh} \times_U U^{sh}$），同理可得这一结论。
由引理 \ref{pione-lemma-retrocompact-dense-open-connected-at-infinity-closed}，沿各
竖直箭头的拉回在有限 \'etale 概形范畴上诱导全忠实函子。因此，
$V \times_U U^{sh}$ 相对于 fpqc 覆盖 $\{U^{sh} \to U\}$ 的典范下降数据，
转化为 $Y^{sh}$ 相对于 fpqc 覆盖 $\{X^{sh} \to X\}$ 的下降数据。由于
$Y^{sh} \to X^{sh}$ 有限，因而仿射，该下降数据有效（《下降》引理
\ref{descent-lemma-affine}）。于是得到仿射态射 $Y \to X$，以及与下降数据
相容的同构 $Y \times_X X^{sh} \to Y^{sh}$。由下降数据的全忠实性
（如《下降》引理 \ref{descent-lemma-refine-coverings-fully-faithful}），得到
同构 $V \to U \times_X Y$。最后，$Y \to X$ 有限 \'etale，因为
$Y^{sh} \to X^{sh}$ 如此；见《下降》引理
\ref{descent-lemma-descending-property-etale} 和
\ref{descent-lemma-descending-property-finite}。
\end{proof}

\noindent
设 $X$ 是不可约概形，$\eta \in X$ 是其泛点。典范态射 $\eta \to X$
诱导典范映射
\begin{equation}
\label{pione-equation-inclusion-generic-point}
\text{Gal}(\kappa(\eta)^{sep}/\kappa(\eta)) = \pi_1(\eta, \overline{\eta})
\longrightarrow \pi_1(X, \overline{\eta})
\end{equation}
左侧的等同来自引理 \ref{pione-lemma-fundamental-group-Galois-group}。

\begin{lemma}
\label{pione-lemma-irreducible-geometrically-unibranch}
设 $X$ 是不可约几何单枝概形。对任意非空开子概形 $U \subset X$，典范映射
$$
\pi_1(U, \overline{u}) \longrightarrow \pi_1(X, \overline{u})
$$
是满射。映射 (\ref{pione-equation-inclusion-generic-point})
$\pi_1(\eta, \overline{\eta}) \to \pi_1(X, \overline{\eta})$
也是满射。
\end{lemma}

\begin{proof}
由引理 \ref{pione-lemma-thickening}，可以作约化后代替 $X$。所以可设
$X$ 是整概形。由引理 \ref{pione-lemma-functoriality-galois-surjective}，本引理的
断言等价于函子 $\textit{F\'Et}_X \to \textit{F\'Et}_U$ 和
$\textit{F\'Et}_X \to \textit{F\'Et}_\eta$ 全忠实。

\medskip\noindent
关于 $\textit{F\'Et}_X \to \textit{F\'Et}_U$ 的结论来自引理
\ref{pione-lemma-quasi-compact-dense-open-connected-at-infinity}，以及如下事实：若
局部环 $A$ 几何单枝，则其严格 Hensel 化的谱不可约。见《代数补篇》引理
\ref{more-algebra-lemma-geometrically-unibranch}。

\medskip\noindent
注意，剩余域 $\kappa(\eta) = \mathcal{O}_{X, \eta}$ 是
$\mathcal{O}_X(U)$ 的滤过余极限，其中 $U \subset X$ 遍历非空仿射开子概形。
因此，$\textit{F\'Et}_\eta$ 是相应范畴 $\textit{F\'Et}_U$ 的余极限，其中
$U$ 如上；见《极限》引理 \ref{limits-lemma-descend-finite-presentation}、
\ref{limits-lemma-descend-finite-finite-presentation} 和
\ref{limits-lemma-descend-etale}。
再作形式论证，便可知函子 $\textit{F\'Et}_X \to \textit{F\'Et}_\eta$ 的
全忠实性可由各函子 $\textit{F\'Et}_X \to \textit{F\'Et}_U$ 的全忠实性推出。
\end{proof}

\begin{lemma}
\label{pione-lemma-exact-sequence-finite-nr-closed-pts}
设 $X$ 是概形。设 $x_1, \ldots, x_n \in X$ 是有限多个闭点，满足
\begin{enumerate}
\item $U = X \setminus \{x_1, \ldots, x_n\}$ 连通，且是 $X$ 的逆紧开子概形；
\item 对每个 $i$，穿孔谱 $U_i^{sh}$（它属于 $\mathcal{O}_{X, x_i}$ 的
严格 Hensel 化）连通。
\end{enumerate}
则映射 $\pi_1(U) \to \pi_1(X)$ 是满射，且其核是 $\pi_1(U)$ 的最小闭正规
子群，它包含各映射 $\pi_1(U_i^{sh}) \to \pi_1(U)$ 的像，其中
$i = 1, \ldots, n$。
\end{lemma}

\begin{proof}
满射性由引理 \ref{pione-lemma-retrocompact-dense-open-connected-at-infinity-closed} 和
\ref{pione-lemma-functoriality-galois-surjective} 得出。可以考虑映射序列
$$
\pi_1(U)  \to \ldots \to
\pi_1(X \setminus \{x_1, x_2\}) \to \pi_1(X \setminus \{x_1\}) \to \pi_1(X)
$$
群论论证表明，只需在 $n = 1$ 时证明关于核的断言（略去细节）。记
$x = x_1$、$U^{sh} = U_1^{sh}$，置 $A = \mathcal{O}_{X, x}$，并令
$A^{sh}$ 为其严格 Hensel 化。考虑图表
$$
\xymatrix{
U \ar[d] &
\Spec(A) \setminus \{\mathfrak m\} \ar[l] \ar[d] &
U^{sh} \ar[d] \ar[l] \\
X & \Spec(A) \ar[l] & \Spec(A^{sh}) \ar[l]
}
$$
由引理 \ref{pione-lemma-functoriality-galois-ses}，须证明：若有限 \'etale 态射
$V \to U$ 拉回到 $U^{sh}$ 后成为平凡覆盖，则它延拓为 $X$ 上的有限
\'etale 概形。由引理 \ref{pione-lemma-local-exact-sequence}，我们已知 $A$ 的
穿孔谱上有限 \'etale 概形的相应断言。另一方面，由《极限》引理
\ref{limits-lemma-glueing-near-closed-point}，$X$ 上有限表示概形等同于把
$U$ 上和 $A$ 上的有限表示概形沿 $A$ 的穿孔谱粘合而得的概形。证明完毕。
\end{proof}









\section{正规概形的基本群}
\label{pione-section-normal}

\noindent
设 $X$ 是整的几何单枝概形。上一节已经看到，$X$ 的基本群是函数域
$K$（即 $X$ 的函数域）的 Galois 群之商。由于该映射连续，其核是 Galois 群的闭正规
子群。因此，由 Galois 理论，该核对应于 Galois 扩张 $M/K$（《域》定理
\ref{fields-theorem-inifinite-galois-theory}）。本节将确定 $M$；这里假设
$X$ 是正规整概形。

\medskip\noindent
设 $X$ 是函数域为 $K$ 的正规整概形。设 $L/K$ 是有限扩张。考虑正规化
$Y \to X$，即 $X$ 在态射 $\Spec(L) \to X$ 中的正规化，其定义见《态射》
第 \ref{morphisms-section-normalization-X-in-Y} 节。在此情形下，称
{\it $X$ 在 $L$ 中不分歧}，如果 $Y \to X$ 是概形的非分歧态射。我们将在
引理 \ref{pione-lemma-unramified} 中阐明这一条件。注意，$Y \to X$ 在
$\Spec(K)$ 上的概形论纤维是 $\Spec(L)$。因此，域扩张 $L/K$ 是可分的，
只要 $X$ 在 $L$ 中不分歧；见《态射》引理
\ref{morphisms-lemma-unramified-over-field}。

\begin{lemma}
\label{pione-lemma-unramified-in-L}
在上述情形下，下列条件等价：
\begin{enumerate}
\item $X$ 在 $L$ 中不分歧；
\item $Y \to X$ 是 \'etale 态射；
\item $Y \to X$ 是有限 \'etale 态射。
\end{enumerate}
\end{lemma}

\begin{proof}
注意，$Y \to X$ 是整态射。在每种情形下，按照定义，态射 $Y \to X$
局部有限型。因此，由《态射》引理 \ref{morphisms-lemma-finite-integral}，
在每种情形下 $Y \to X$ 都有限。特别地，(2) 与 (3) 等价。
\'Etale 态射非分歧，故 (2) 蕴含 (1)。

\medskip\noindent
反之，假设 $Y \to X$ 非分歧。由于正规概形几何单枝（《性质》引理
\ref{properties-lemma-normal-geometrically-unibranch}），由《态射补篇》引理
\ref{more-morphisms-lemma-global-unramified-dominant-unibranch-is-etale} 可知
态射 $Y \to X$ 是 \'etale 态射。下一段还将给出直接证明。

\medskip\noindent
设 $x \in X$。可以选取 \'etale 邻域 $(U, u) \to (X, x)$，使得
$$
Y \times_X U = \coprod V_j \longrightarrow U
$$
是若干闭浸入的不交并；见《\'Etale 态射》引理
\ref{etale-lemma-finite-unramified-etale-local}。通过缩小，可以假设 $U$
拟紧。于是 $U$ 只有有限多个不可约分支（《下降》引理
\ref{descent-lemma-locally-finite-nr-irred-local-fppf}）。由于 $U$ 正规
（《下降》引理 \ref{descent-lemma-normal-local-smooth}），$U$ 的不可约分支
都是开闭的（《性质》引理
\ref{properties-lemma-normal-locally-finite-nr-irreducibles}），故可以假设
$U$ 不可约。于是 $U$ 是整概形，其泛点 $\xi$ 映到 $X$ 的泛点。另一方面，
已知 $Y \times_X U$ 是 $U$ 在 $\Spec(L) \times_X U$ 中的正规化；见
《态射补篇》引理
\ref{more-morphisms-lemma-normalization-smooth-localization}。
$\Spec(L) \times_X U$ 的每一点都映到 $\xi$。因此，由《态射》引理
\ref{morphisms-lemma-normalization-generic}，每个 $V_j$ 都包含一个映到
$\xi$ 的点。态射 $V_j \to U$ 是同构，因为 $U = \overline{\{\xi\}}$。
所以 $Y \times_X U \to U$ 是 \'etale 态射。由《下降》引理
\ref{descent-lemma-descending-property-etale}，可知 $Y \to X$ 在
$U \to X$ 的像（即 $x$ 的一个开邻域）上是 \'etale 态射。
\end{proof}

\begin{lemma}
\label{pione-lemma-finite-etale-covering-normal-unramified}
设 $X$ 是函数域为 $K$ 的正规整概形。设 $Y \to X$ 是有限 \'etale 态射。
若 $Y$ 连通，则 $Y$ 是正规整概形，且 $Y$ 是 $X$ 在 $Y$ 的函数域中的
正规化。
\end{lemma}

\begin{proof}
由《下降》引理 \ref{descent-lemma-normal-local-smooth}，概形 $Y$ 正规。
由于 $Y \to X$ 平坦，由《态射》引理
\ref{morphisms-lemma-generalizations-lift-flat}，$Y$ 的每个泛点都映到 $X$
的泛点。由于 $Y \to X$ 有限，$Y$ 只有有限多个不可约分支。因此，由
《性质》引理 \ref{properties-lemma-normal-locally-finite-nr-irreducibles}，
$Y$ 是有限多个正规整概形的不交并。所以，若 $Y$ 连通，则 $Y$ 是正规
整概形。

\medskip\noindent
设 $L$ 是 $Y$ 的函数域，并令 $Y' \to X$ 为 $X$ 在 $L$ 中的正规化。
由《态射》引理 \ref{morphisms-lemma-characterize-normalization}，得到分解
$Y' \to Y \to X$，并且 $Y' \to Y$ 是 $Y$ 在 $L$ 中的正规化。由于 $Y$
正规，显然 $Y' = Y$（这也可由《态射》引理
\ref{morphisms-lemma-finite-birational-over-normal} 得出）。
\end{proof}

\begin{proposition}
\label{pione-proposition-normal}
设 $X$ 是函数域为 $K$ 的正规整概形。则典范映射
(\ref{pione-equation-inclusion-generic-point})
$$
\text{Gal}(K^{sep}/K) = \pi_1(\eta, \overline{\eta})
\longrightarrow \pi_1(X, \overline{\eta})
$$
等同于商映射 $\text{Gal}(K^{sep}/K) \to \text{Gal}(M/K)$，其中
$M \subset K^{sep}$ 是有限子扩张 $L$ 之并，其中 $X$ 在 $L$ 中不分歧。
\end{proposition}

\begin{proof}
正规概形 $X$ 几何单枝（《性质》引理
\ref{properties-lemma-normal-geometrically-unibranch}）。因此，引理
\ref{pione-lemma-irreducible-geometrically-unibranch} 可应用于 $X$。所以
$\pi_1(\eta, \overline{\eta}) \to \pi_1(X, \overline{\eta})$ 是满射，且
交换图表的上方水平箭头
$$
\xymatrix{
\textit{F\'Et}_X \ar[r] \ar[d] \ar[rd]_c & \textit{F\'Et}_\eta \ar[d] \\
\textit{Finite-}\pi_1(X, \overline{\eta})\textit{-sets} \ar[r] &
\textit{Finite-}\text{Gal}(K^{sep}/K)\textit{-sets}
}
$$
全忠实。左侧竖直箭头是定理 \ref{pione-theorem-fundamental-group} 的等价，
右侧竖直箭头是引理 \ref{pione-lemma-fundamental-group-Galois-group} 的等价。
下方水平箭头由本命题中的映射诱导。由引理
\ref{pione-lemma-unramified-in-L} 和
\ref{pione-lemma-finite-etale-covering-normal-unramified}，可知 $c$ 的本质像由
同构于下列形式集合的 $\text{Gal}(K^{sep}/K)\textit{-Sets}$ 构成：
$$
S = \Hom_K(\prod\nolimits_{i = 1, \ldots, n} L_i, K^{sep}) =
\coprod\nolimits_{i = 1, \ldots, n} \Hom_K(L_i, K^{sep})
$$
其中 $L_i/K$ 是有限可分扩张，且 $X$ 在 $L_i$ 中不分歧。因此，若
$M \subset K^{sep}$ 如该引理的陈述所述，则 $\text{Gal}(K^{sep}/M)$
恰为 $\text{Gal}(K^{sep}/K)$ 中在 $c$ 的本质像的每个对象上平凡作用的
子群。另一方面，$c$ 的本质像恰为这样的 $S$ 所成的范畴：
$\text{Gal}(K^{sep}/K)$ 的作用经由满射
$\text{Gal}(K^{sep}/K) \to \pi_1(X, \overline{\eta})$ 分解。因此，
$\text{Gal}(K^{sep}/M)$ 是其核。于是 $\text{Gal}(K^{sep}/M)$ 是正规子群，
$M/K$ 是 Galois 扩张，并且有短正合列
$$
1 \to \text{Gal}(K^{sep}/M) \to
\text{Gal}(K^{sep}/K) \to
\text{Gal}(M/K) \to 1
$$
这是由 Galois 理论得到的（《域》定理
\ref{fields-theorem-inifinite-galois-theory} 和引理
\ref{fields-lemma-ses-infinite-galois}）。证明完毕。
\end{proof}

\begin{lemma}
\label{pione-lemma-local-exact-sequence-normal}
设 $(A, \mathfrak m)$ 是正规局部环。置 $X = \Spec(A)$。令 $A^{sh}$ 为
$A$ 的严格 Hensel 化。令 $K$ 和 $K^{sh}$ 分别为 $A$ 和 $A^{sh}$ 的
分式域。则序列
$$
\pi_1(\Spec(K^{sh})) \to \pi_1(\Spec(K)) \to \pi_1(X) \to 1
$$
按引理 \ref{pione-lemma-functoriality-galois-ses} 第 (1) 部分的意义是正合的。
\end{lemma}

\begin{proof}
注意，$A^{sh}$ 是正规整环，见《代数补篇》引理
\ref{more-algebra-lemma-henselization-normal}。映射
$\pi_1(\Spec(K)) \to \pi_1(X)$ 由命题 \ref{pione-proposition-normal} 可知是满射。

\medskip\noindent
记 $X^{sh} = \Spec(A^{sh})$。设 $Y \to X$ 是有限 \'etale 态射。则
$Y^{sh} = Y \times_X X^{sh} \to X^{sh}$ 是有限 \'etale 态射。由于
$A^{sh}$ 严格 Hensel，$Y^{sh}$ 同构于若干个 $X^{sh}$ 的不交并。因此，
$Y \times_X \Spec(K^{sh})$ 亦然。于是复合
$\pi_1(\Spec(K^{sh})) \to \pi_1(X)$ 平凡，见引理
\ref{pione-lemma-composition-trivial}。

\medskip\noindent
为完成证明，按照引理 \ref{pione-lemma-functoriality-galois-ses}，只需证明：
给定有限 \'etale 态射 $V \to \Spec(K)$，若
$V \times_{\Spec(K)} \Spec(K^{sh})$ 是若干个 $\Spec(K^{sh})$ 的不交并，
则可以找到有限 \'etale 态射 $Y \to X$，使得
$V \cong Y \times_X \Spec(K)$ 在 $\Spec(K)$ 上成立。记
$V = \Spec(L)$，则 $L$ 是 $K$ 的若干有限可分扩张之有限直积。令
$B \subset L$ 为 $A$ 在 $L$ 中的整闭包。若 $A \to B$ 是 \'etale 态射，
则可以取 $Y = \Spec(B)$，证明即告完成。由《代数》引理
\ref{algebra-lemma-integral-closure-commutes-smooth}（以及一个略去的极限
论证），可知 $B \otimes_A A^{sh}$ 是 $A^{sh}$ 在
$L^{sh} = L \otimes_K K^{sh}$ 中的整闭包。我们的假设是 $L^{sh}$ 为若干个
$K^{sh}$ 的直积，因而 $B^{sh}$ 为若干个 $A^{sh}$ 的直积。所以
$A^{sh} \to B^{sh}$ 是 \'etale 态射。由于 $A \to A^{sh}$ 忠实平坦，
$A \to B$ 是 \'etale 态射（《下降》引理
\ref{descent-lemma-descending-property-etale}），正如所需。
\end{proof}







\section{群作用与整闭包}
\label{pione-section-group-actions-integral}

\noindent
本节继续《代数补篇》第
\ref{more-algebra-section-group-actions-integral} 节的讨论。
回顾正规局部环按定义是整环。

\begin{lemma}
\label{pione-lemma-get-algebraic-closure}
设 $A$ 是正规整环，其分式域 $K$ 可分代数闭。设
$\mathfrak p \subset A$ 是非零素理想。则剩余域
$\kappa(\mathfrak p)$ 代数闭。
\end{lemma}

\begin{proof}
反设该引理不成立以导出矛盾。于是存在首一不可约多项式
$P(T) \in \kappa(\mathfrak p)[T]$，其次数 $d > 1$。将 $P$ 替换为
$a^d P(a^{-1}T)$，其中 $a \in A$ 适当（以清除分母）之后，可以假设
$P$ 是 $Q$ 的像，其中后者是 $A[T]$ 中的首一多项式。注意，$Q$ 在
$K[T]$ 中不可约。
事实上，由《代数》引理 \ref{algebra-lemma-polynomials-divide}，$K$ 上的一个
因式分解会给出 $A$ 上的一个因式分解；将其模 $\mathfrak p$ 约化，就会得到
$P$ 的因式分解。由于 $K$ 可分闭，$Q$ 不是可分多项式
（《域》定义 \ref{fields-definition-separable}）。于是 $K$ 的特征为
$p > 0$，且 $Q$ 的一次项为零（《域》定义
\ref{fields-definition-separable}）。然而，此时可将 $Q$ 替换为
$Q + a T$，其中 $a \in \mathfrak p$ 非零，从而得到矛盾。
\end{proof}

\begin{lemma}
\label{pione-lemma-normal-local-domain-separablly-closed-fraction-field}
分式域可分闭的正规局部环是严格 Hensel 环。
\end{lemma}

\begin{proof}
设 $(A, \mathfrak m, \kappa)$ 是正规局部环，其分式域 $K$ 可分闭。
若 $A = K$，则结论成立。否则，由引理
\ref{pione-lemma-get-algebraic-closure}，剩余域 $\kappa$ 代数闭，因而只需验证
$A$ 是 Hensel 环。设 $f \in A[T]$ 首一，并设 $a_0 \in \kappa$ 是重数为
$1$ 的根，这里所约化的多项式为 $\overline{f} \in \kappa[T]$。设
$f = \prod f_i$ 是其在 $K[T]$ 中的因式分解。由《代数》引理
\ref{algebra-lemma-polynomials-divide}，有 $f_i \in A[T]$。因此，$a_0$ 是
某个 $f_i$ 的根，其中指标为 $i$。将 $f$ 替换为 $f_i$ 后，可以假设
$f$ 不可约。此时，导数 $f'$ 不可能在 $A[T]$ 中为零，因为 $a_0$ 是单根；
又可断定 $f$ 必为一次多项式，因为 $K$ 可分代数闭。因此 $A$ 是 Hensel 环，见《代数》
定义 \ref{algebra-definition-henselian}。
\end{proof}

\begin{lemma}
\label{pione-lemma-inertia-base-change}
设有限群 $G$ 作用于环 $R$。设 $R^G \to A$ 是环同态。设
$\mathfrak q' \subset A \otimes_{R^G} R$ 是位于素理想
$\mathfrak q \subset R$ 之上的素理想。则
$$
I_\mathfrak q = \{\sigma \in G \mid
\sigma(\mathfrak q) = \mathfrak q\text{ 且 }
\sigma \bmod \mathfrak q = \text{id}_{\kappa(\mathfrak q)}\}
$$
等于
$$
I_{\mathfrak q'} = \{\sigma \in G \mid
\sigma(\mathfrak q') = \mathfrak q'\text{ 且 }
\sigma \bmod \mathfrak q' = \text{id}_{\kappa(\mathfrak q')}\}
$$
\end{lemma}

\begin{proof}
由于 $\mathfrak q$ 是 $\mathfrak q'$ 的逆像，且
$\kappa(\mathfrak q) \subset \kappa(\mathfrak q')$，可得
$I_{\mathfrak q'} \subset I_\mathfrak q$。反过来，若
$\sigma \in I_\mathfrak q$，则该 $\sigma$ 在纤维环
$A \otimes_{R^G} \kappa(\mathfrak q)$ 上平凡作用。因此，$\sigma$ 固定所有
位于 $\mathfrak q$ 之上的素理想，并在它们的剩余域上诱导恒等映射。
\end{proof}

\begin{lemma}
\label{pione-lemma-inertia-invariants-etale}
设有限群 $G$ 作用于环 $R$。设 $\mathfrak q \subset R$ 是素理想。置
$$
I = \{\sigma \in G \mid \sigma(\mathfrak q) = \mathfrak q
\text{ 且 } \sigma \bmod \mathfrak q = \text{id}_\mathfrak q\}
$$
则 $R^G \to R^I$ 在 $R^I \cap \mathfrak q$ 处为 \'etale 的。
\end{lemma}

\begin{proof}
证明的策略是利用 \'etale 局部化，约化到 $R \to R^I$ 在
$R^I \cap \mathfrak p$ 处为局部同构的情形。设 $R^G \to A$ 是 \'etale
环同态。我们断言：如果对于 $G$ 在 $A \otimes_{R^G} R$ 上的作用以及
某个素理想 $\mathfrak q'$（它是 $A \otimes_{R^G} R$ 中位于
$\mathfrak q$ 之上的素理想），该结论成立，那么原结论亦成立。

\medskip\noindent
为验证这一点，注意，由于 $R^G \to A$ 平坦，有
$A = (A \otimes_{R^G} R)^G$，见《代数补篇》引理
\ref{more-algebra-lemma-base-change-invariants}。由引理
\ref{pione-lemma-inertia-base-change}，群 $I$ 不变。再次应用《代数补篇》引理
\ref{more-algebra-lemma-base-change-invariants} 可知
$A \otimes_{R^G} R^I = (A \otimes_{R^G} R)^I$
（因为 $R^I \to A \otimes_{R^G} R^I$ 平坦）。因此
$$
\xymatrix{
\Spec((A \otimes_{R^G} R)^I) \ar[d] \ar[r] & \Spec(R^I) \ar[d] \\
\Spec(A) \ar[r] & \Spec(R^G)
}
$$
是笛卡儿图，且水平箭头为 \'etale 的。因此，若左侧竖直箭头在某个开邻域
$W$（它是 $(A \otimes_{R^G} R)^I \cap \mathfrak q'$ 的邻域）上为
\'etale 的，则右侧竖直箭头在 $W$ 于 $\Spec(R^I)$ 中的（开）像的各点处
为 \'etale 的，见《下降》引理 \ref{descent-lemma-smooth-permanence}。
特别地，态射 $\Spec(R^I) \to \Spec(R^G)$ 在
$R^I \cap \mathfrak q$ 处为 \'etale 的。

\medskip\noindent
令 $\mathfrak p = R^G \cap \mathfrak q$。由《代数补篇》引理
\ref{more-algebra-lemma-one-orbit}，$\Spec(R) \to \Spec(R^G)$ 在
$\mathfrak p$ 上的纤维是有限的。再由《代数补篇》引理
\ref{more-algebra-lemma-one-orbit-geometric}，这些点处的剩余域扩张是
代数且正规的，并具有有限自同构群。因此，可以把《态射补篇》引理
\ref{more-morphisms-lemma-etale-makes-integral-split} 应用于整同态
$R^G \to R$ 和素理想 $\mathfrak p$。结合上述断言，可约化到如下情形：
$R = A_1 \times \ldots \times A_n$，每个 $A_i$ 恰有一个素理想
$\mathfrak q_i$ 位于 $\mathfrak p$ 之上，且剩余域扩张
$\kappa(\mathfrak q_i)/\kappa(\mathfrak p)$ 纯不可分。当然，
$\mathfrak q$ 是这些素理想之一，不妨设 $\mathfrak q = \mathfrak q_1$。

\medskip\noindent
$G$ 未必置换各因子 $A_i$（若 $A_i$ 的谱连通，则会如此；例如，若
$R^G$ 是局部环）。可按如下方式修正这一点；我们建议读者自行推演，或许可以
用幂等元代替拓扑来论证。回顾上述乘积分解给出 $\Spec(R)$ 相应的开闭子集
$U_i$ 的不交并分解。由于 $G$ 有限，可以把这个覆盖细化为有限不交并分解
$\Spec(R) = \coprod_{j \in J} W_j$，其中 $W_j$ 是开闭子集，使得对所有
$j \in J$，存在
$j' \in J$ 满足 $\sigma(W_j) = W_{j'}$。在这些 $W_j$ 中，不与
$\{\mathfrak q_1, \ldots, \mathfrak q_n\}$ 相交者的并是一个不与
$\mathfrak p$ 上纤维相交的闭子集。它因而映到 $\Spec(R^G)$ 中一个不含
$\mathfrak p$ 的闭子集，这是因为
$\Spec(R) \to \Spec(R^G)$ 是闭映射。因此，将 $R^G$ 替换为一个主局部化之后（上述断言
允许这样做），可以假设每个 $W_j$ 都与某个点 $\mathfrak q_i$ 相交。此时，
置 $U_i = W_j$，其中 $\mathfrak q_i \in W_j$。相应的乘积分解
$R = A_1 \times \ldots \times A_n$ 具有 $G$ 置换各因子 $A_i$ 的性质。

\medskip\noindent
所以，可以假设有一个乘积分解 $R = A_1 \times \ldots \times A_n$，
它与 $G$-作用相容；其中每个 $A_i$ 恰有一个素理想 $\mathfrak q_i$ 位于
$\mathfrak p$ 之上，且域扩张
$\kappa(\mathfrak q_i)/\kappa(\mathfrak p)$ 纯不可分。记
$A' = A_2 \times \ldots \times A_n$，于是
$$
R = A_1 \times A'
$$
由于 $\mathfrak q = \mathfrak q_1$，每个 $\sigma \in I$ 都保持上述乘积分解。
因此
$$
R^I = (A_1)^I \times (A')^I
$$
注意，有 $I = D = \{\sigma \in G \mid \sigma(\mathfrak q) = \mathfrak q\}$，
因为 $\kappa(\mathfrak q)/\kappa(\mathfrak p)$ 纯不可分。
由于 $G$ 在 $\mathfrak p$ 之上的素理想上的作用是传递的（《代数补篇》引理
\ref{more-algebra-lemma-one-orbit}），可知 $I$ 在 $G$ 中的指数为 $n$，并且可写成
$G = eI \amalg \sigma_2I \amalg \ldots \amalg \sigma_nI$，使得
$A_i = \sigma_i(A_1)$ 对 $i = 2, \ldots, n$ 成立。由此得到
$$
R^G = (A_1)^I.
$$
所以映射 $R^G \to R^I$ 在 $R^I \cap \mathfrak q$ 处为 \'etale 的，
证明完毕。
\end{proof}

\noindent
下一个引理推广了《代数补篇》引理
\ref{more-algebra-lemma-inertial-invariants-unramified}。

\begin{lemma}
\label{pione-lemma-inertial-invariants-unramified}
设 $A$ 是分式域为 $K$ 的正规整环。设 $L/K$ 是一个（可能无限的）Galois
扩张。令 $G = \text{Gal}(L/K)$，并令 $B$ 为 $A$ 在 $L$ 中的整闭包。
设 $\mathfrak q \subset B$。置
$$
I = \{\sigma \in G \mid
\sigma(\mathfrak q) = \mathfrak q \text{ 且 }
\sigma \bmod \mathfrak q = \text{id}_{\kappa(\mathfrak q)}\}
$$
则 $(B^I)_{B^I \cap \mathfrak q}$ 是 \'etale $A$-代数的滤过余极限。
\end{lemma}

\begin{proof}
可将 $L$ 写成 $K$ 的有限 Galois 扩张的滤过余极限。因此，只需在 $L/K$
为有限 Galois 扩张时证明该引理，见《代数》引理
\ref{algebra-lemma-colimit-colimit-etale}。因为 $A = B^G$（$A$ 在
$K = L^G$ 中整闭），故结论由引理 \ref{pione-lemma-inertia-invariants-etale} 得出。
\end{proof}







\section{分歧理论}
\label{pione-section-ramification}

\noindent
本节继续 More on Algebra, Section
\ref{more-algebra-section-ramification} 中的讨论，并将其与我们对概形基本群的讨论联系起来。

\medskip\noindent
设 $(A, \mathfrak m, \kappa)$ 是分式域为 $K$ 的正规局部环。选取一个可分代数闭包
$K^{sep}$。设 $A^{sep}$ 为 $A$ 在 $K^{sep}$ 中的整闭包，并选取极大理想
$\mathfrak m^{sep} \subset A^{sep}$。设 $A \subset A^h \subset A^{sh}$
分别为 Hensel 化和严格 Hensel 化。注意，$A^h$ 和 $A^{sh}$ 也都是正规环
（More on Algebra, Lemma \ref{more-algebra-lemma-henselization-normal}）。以
$K^h$ 和 $K^{sh}$ 表示它们的分式域。由 Lemma
\ref{pione-lemma-normal-local-domain-separablly-closed-fraction-field}，
$(A^{sep})_{\mathfrak m^{sep}}$ 是严格 Hensel 环，故可以选取一个 $A$-代数映射
$A^{sh} \to (A^{sep})_{\mathfrak m^{sep}}$。具体来说，先选取一个
$\kappa$-嵌入\footnote{这是可行的，因为 $\kappa(\mathfrak m^{sh})$ 是
$\kappa$ 的可分代数闭包，而由 Lemma \ref{pione-lemma-get-algebraic-closure}，
$\kappa(\mathfrak m^{sep})$ 是 $\kappa$ 的代数闭包。}
$\kappa(\mathfrak m^{sh}) \to \kappa(\mathfrak m^{sep})$，再由 Algebra,
Lemma \ref{algebra-lemma-strictly-henselian-functorial} 将其（唯一地）延拓为
一个 $A$-代数同态。由此得到交换图
$$
\xymatrix{
K^{sep} & K^{sh} \ar[l] & K^h \ar[l] & K \ar[l] \\
(A^{sep})_{\mathfrak m^{sep}} \ar[u] &
A^{sh} \ar[u] \ar[l] &
A^h \ar[u] \ar[l] &
A \ar[u] \ar[l]
}
$$
可以取这些环之谱的基本群。当然，因为 $K^{sep}$、
$(A^{sep})_{\mathfrak m^{sep}}$ 和 $A^{sh}$ 都是严格 Hensel 的，所得的群均为平凡群。
因此有意义的部分是
\begin{equation}
\label{pione-equation-inertia-diagram-pione}
\vcenter{
\xymatrix{
\pi_1(U^{sh}) \ar[r] \ar[rd]_1 & \pi_1(U^h) \ar[d] \ar[r] & \pi_1(U) \ar[d] \\
& \pi_1(X^h) \ar[r] & \pi_1(X)
}
}
\end{equation}
这里 $X^h$ 和 $X$ 分别是 $A^h$ 和 $A$ 的谱，而 $U^{sh}$、$U^h$、$U$
分别是 $K^{sh}$、$K^h$、$K$ 的谱。标号 $1$ 表示该映射是平凡的；这是因为它经过平凡群
$\pi_1(X^{sh})$ 分解。另一方面，profinite 群
$G = \text{Gal}(K^{sep}/K)$ 作用在 $A^{sep}$ 上，因而可以作如下定义：
$$
D = \{\sigma \in G \mid \sigma(\mathfrak m^{sep}) = \mathfrak m^{sep}\}
\supset
I = \{\sigma \in D \mid \sigma \bmod \mathfrak m^{sep} =
\text{id}_{\kappa(\mathfrak m^{sep})}\}
$$
这两个群有时分别称为{\it 分解群}和{\it 惯性群}，尤其是在 $A$ 为离散赋值环时。

\begin{lemma}
\label{pione-lemma-identify-inertia}
在上述情形中，通过同构 $\pi_1(U) = \text{Gal}(K^{sep}/K)$，交换图
(\ref{pione-equation-inertia-diagram-pione}) 转化为交换图
$$
\xymatrix{
I \ar[r] \ar[rd]_1 & D \ar[d] \ar[r] & \text{Gal}(K^{sep}/K) \ar[d] \\
& \text{Gal}(\kappa(\mathfrak m^{sh})/\kappa) \ar[r] & \text{Gal}(M/K)
}
$$
其中 $K^{sep}/M/K$ 是关于 $A$ 非分歧的最大子扩张。此外，竖直箭头均为满射，
左侧竖直箭头的核为 $I$，右侧竖直箭头的核是
$\text{Gal}(K^{sep}/K)$ 中包含 $I$ 的最小闭正规子群。
\end{lemma}

\begin{proof}
按构造，群 $D$ 作用于 $(A^{sep})_{\mathfrak m^{sep}}$，且该作用在 $A$ 上。
给定剩余域上的映射后，
$A^{sh} \to (A^{sep})_{\mathfrak m^{sep}}$ 具有唯一性（Algebra, Lemma
\ref{algebra-lemma-strictly-henselian-functorial}），故映射
$A^{sh} \to (A^{sep})_{\mathfrak m^{sep}}$ 的像包含于
$((A^{sep})_{\mathfrak m^{sep}})^I$。另一方面，Lemma
\ref{pione-lemma-inertial-invariants-unramified} 表明
$((A^{sep})_{\mathfrak m^{sep}})^I$ 是 $A$ 的 \'etale 扩张的滤过余极限。
由于 $A^{sh}$ 是最大的这种扩张，遂有
$A^{sh} = ((A^{sep})_{\mathfrak m^{sep}})^I$，从而 $K^{sh} = (K^{sep})^I$。

\medskip\noindent
回顾 $I$ 是一个满射 $D \to \text{Aut}(\kappa(\mathfrak m^{sep})/\kappa)$
的核，见 More on Algebra, Lemma
\ref{more-algebra-lemma-one-orbit-geometric-galois}。我们有
$\text{Aut}(\kappa(\mathfrak m^{sep})/\kappa) =
\text{Gal}(\kappa(\mathfrak m^{sh})/\kappa)$，因为由上文可知，这些域分别是
$\kappa$ 的代数闭包与可分代数闭包。另一方面，由 Algebra, Lemma
\ref{algebra-lemma-henselian-functorial} 中的唯一性，$A^{sh}$ 在 $A$ 上的任意自同构
也是 $A^{sh}$ 在 $A^h$ 上的自同构。此外，$A^{sh}$ 是有限 \'etale 扩张
$A^h \subset A'$ 的余极限，而这些扩张与有限可分扩张成 $1$ 对 $1$ 的对应，
后者形如 $\kappa'/\kappa$，
见 Algebra, Remark \ref{algebra-remark-construct-sh-from-h}。因此
$$
\text{Aut}(A^{sh}/A) = \text{Aut}(A^{sh}/A^h) =
\text{Gal}(\kappa(\mathfrak m^{sh})/\kappa)
$$
设 $\kappa'/\kappa$ 是 Galois 群为 $G$ 的有限 Galois 扩张。由 Algebra, Lemma
\ref{algebra-lemma-henselian-cat-finite-etale}，设 $A^h \subset A'$ 是与
$\kappa \subset \kappa'$ 对应的有限 \'etale 扩张。考察分式域与次数可得
$(A')^G = A^h$（略去一个小细节）。取余极限即得
$(A^{sh})^{\text{Gal}(\kappa(\mathfrak m^{sh})/\kappa)} = A^h$。
综合以上各点，可得 $A^h = ((A^{sep})_{\mathfrak m^{sep}})^D$，从而
$K^h = (K^{sep})^D$。

\medskip\noindent
因为 $U$、$U^h$、$U^{sh}$ 分别是域 $K$、$K^h$、$K^{sh}$ 的谱，所以由 Lemma
\ref{pione-lemma-fundamental-group-Galois-group} 可知两个交换图的上行彼此对应。由 Lemma
\ref{pione-lemma-gabber} 有
$\pi_1(X^h) = \text{Gal}(\kappa(\mathfrak m^{sh})/\kappa)$。上文已经指出正合列
$1 \to I \to D \to \text{Gal}(\kappa(\mathfrak m^{sh})/\kappa) \to 1$
的正合性。由 Proposition \ref{pione-proposition-normal} 可知
$\pi_1(X) = \text{Gal}(M/K)$。最后，关于
$\text{Gal}(K^{sep}/K) \to \text{Gal}(M/K) = \pi_1(X)$ 之核的断言由 Lemma
\ref{pione-lemma-local-exact-sequence-normal} 得出。证毕。
\end{proof}

\noindent
设 $X$ 是函数域为 $K$ 的正规整概形，$K^{sep}$ 是 $K$ 的一个可分代数闭包，
而 $X^{sep} \to X$ 是 $X$ 在 $K^{sep}$ 中的正规化。由于
$G = \text{Gal}(K^{sep}/K)$ 作用于 $K^{sep}$，便得到 $G$ 在 $X^{sep}$ 上的右作用。
对 $y \in X^{sep}$，定义
$$
D_y = \{\sigma \in G \mid \sigma(y) = y\} \supset
I_y = \{\sigma \in D \mid \sigma \bmod \mathfrak m_y =
\text{id}_{\kappa(y)} \}
$$
与上文类似。另一方面，对 $x \in X$，设 $\mathcal{O}_{X, x}^{sh}$ 是一个严格
Hensel 化，设 $K_x^{sh}$ 是 $\mathcal{O}_{X, x}^{sh}$ 的分式域，并选取一个
$K$-嵌入 $K_x^{sh} \to K^{sep}$。

\begin{lemma}
\label{pione-lemma-normal-pione-quotient-inertia}
设 $X$ 是函数域为 $K$ 的正规整概形。沿用上述记号，
$\text{Gal}(K^{sep}/K) = \pi_1(\Spec(K))$ 的下列三个子群相等：
\begin{enumerate}
\item 满射
$\text{Gal}(K^{sep}/K) \longrightarrow \pi_1(X)$,
的核；
\item 包含所有 $I_y$（其中 $y \in X^{sep}$）的最小闭正规子群；以及
\item 包含所有 $\text{Gal}(K^{sep}/K_x^{sh})$（其中 $x \in  X$）的最小闭正规子群。
\end{enumerate}
\end{lemma}

\begin{proof}
(2) 与 (3) 的等价性来自 Lemma \ref{pione-lemma-identify-inertia}：$I_y$ 与
$\text{Gal}(K^{sep}/K_x^{sh})$ 共轭，如果 $y$ 位于 $x$ 上方。由 Lemma
\ref{pione-lemma-local-exact-sequence-normal} 可知
$\text{Gal}(K^{sep}/K_x^{sh})$ 到 $\pi_1(\Spec(\mathcal{O}_{X, x}))$ 的映射是平凡的，
因此 (2) 与 (3) 中的子群 $N \subset G = \text{Gal}(K^{sep}/K)$ 包含于
$G \longrightarrow \pi_1(X)$ 的核。

\medskip\noindent
为证明反向包含关系，因 $N$ 正规，只需证明：若 $N \subset U \subset G$ 且 $U$ 是开正规子群，
则商映射 $G \to G/U$ 经过 $\pi_1(X)$ 分解。换言之，若 $L/K$ 是与 $U$ 对应的 Galois
扩张，则必须证明 $X$ 在 $L$ 中非分歧（Section \ref{pione-section-normal}，特别是
Proposition \ref{pione-proposition-normal}）。只需在 $X$ 仿射时证明这一点（如此便可在以下证明中引用代数结论）。
设 $Y \to X$ 是 $X$ 在 $L$ 中的正规化。包含映射 $L \subset K^{sep}$ 诱导态射
$\pi : X^{sep} \to Y$。对 $y \in X^{sep}$，$\pi(y)$ 在 $\text{Gal}(L/K)$ 中的惯性群
是 $I_y$ 在 $\text{Gal}(L/K)$ 中的像；这由 More on Algebra, Lemma
\ref{more-algebra-lemma-one-orbit-geometric-galois-compare} 得出。由于 $N \subset U$，
所有这些惯性群都平凡。应用 Lemma \ref{pione-lemma-inertia-invariants-etale} 可知
$Y \to X$ 是 \'etale 的。（另一种做法是使用 Lemma
\ref{pione-lemma-local-exact-sequence-normal}，看出 $Y$ 在
$\Spec(\mathcal{O}_{X, x})$ 上的拉回对所有 $x \in X$ 都是 \'etale 的，
再多做一些工作便可得出结论。）
\end{proof}

\begin{example}
\label{pione-example-bigger-codim}
设 $X$ 是函数域为 $K$ 的正规整 Noether 概形。分歧轨迹的纯性（见下文）告诉我们，
若 $X$ 正则，则在 Lemma \ref{pione-lemma-normal-pione-quotient-inertia} 中，只需考虑惯性群
$I = \pi_1(\Spec(K_x^{sh}))$，其中 $x$ 是余维为 $1$ 的 $X$ 中的点。
但一般而言这并不充分。事实上，令
$Y = \mathbf{A}_k^n = \Spec(k[t_1, \ldots, t_n])$，其中 $k$ 是特征不为 $2$ 的域。
令 $G = \{\pm 1\}$ 是阶为 $2$ 的群，它通过各坐标上的乘法作用于 $Y$。令
$$
X = \Spec(k[t_it_j, i, j \in \{1, \ldots, n\}])
$$
嵌入 $k[t_it_j] \subset k[t_1, \ldots, t_n]$ 定义了一个次数为 $2$ 的态射
$Y \to X$；除极大理想 $\mathfrak m = (t_it_j)$ 上方外，它处处非分歧，
而该理想给出一个余维为 $n$ 的 $X$ 中的点。
\end{example}

\begin{lemma}
\label{pione-lemma-unramified}
设 $X$ 是函数域为 $K$ 的整正规概形。设 $L/K$ 是有限扩张，并设 $Y \to X$
是 $X$ 在 $L$ 中的正规化。下列条件等价：
\begin{enumerate}
\item 按 Section \ref{pione-section-normal} 中的定义，$X$ 在 $L$ 中非分歧；
\item $Y \to X$ 是概形的非分歧态射；
\item $Y \to X$ 是概形的 \'etale 态射；
\item $Y \to X$ 是概形的有限 \'etale 态射；
\item 对 $x \in X$，投影
$Y \times_X \Spec(\mathcal{O}_{X, x}) \to \Spec(\mathcal{O}_{X, x})$
非分歧；
\item 将 (5) 中的局部环换成 $\mathcal{O}_{X, x}^h$ 后，结论仍成立；
\item 将 (5) 中的局部环换成 $\mathcal{O}_{X, x}^{sh}$ 后，结论仍成立；
\item 对 $x \in X$，概形论纤维 $Y_x$ 在 $x$ 上是次数 $\geq [L : K]$ 的 \'etale 概形。
\end{enumerate}
若 $L/K$ 是 Galois 群为 $G$ 的 Galois 扩张，则上述条件还等价于
\begin{enumerate}
\item[(9)] 对 $y \in Y$，群
$I_y = \{g \in G \mid g(y) = y\text{ 且 }
g \bmod \mathfrak m_y = \text{id}_{\kappa(y)}\}$ 平凡。
\end{enumerate}
\end{lemma}

\begin{proof}
(1) 与 (2) 的等价性就是 (1) 的定义。(2)、(3)、(4) 的等价性由 Lemma
\ref{pione-lemma-unramified-in-L} 给出。容易证明 (4) $\Rightarrow$ (5)、
(5) $\Rightarrow$ (6)、(6) $\Rightarrow$ (7)。

\medskip\noindent
假设 (7) 成立。注意 $\mathcal{O}_{X, x}^{sh}$ 是正规局部整环
（More on Algebra, Lemma \ref{more-algebra-lemma-henselization-normal}）。
令 $L^{sh} = L \otimes_K K_x^{sh}$，其中 $K_x^{sh}$ 是
$\mathcal{O}_{X, x}^{sh}$ 的分式域。于是
$L^{sh} = \prod_{i = 1, \ldots, n} L_i$，且 $L_i/K_x^{sh}$ 是有限可分扩张。
由 Algebra, Lemma \ref{algebra-lemma-integral-closure-commutes-smooth}
（以及一个略去的极限论证）可知，$Y \times_X \Spec(\mathcal{O}_{X, x}^{sh})$
是 $\Spec(\mathcal{O}_{X, x}^{sh})$ 在 $L^{sh}$ 中的整闭包。因此，由 Lemma
\ref{pione-lemma-unramified-in-L}（分别应用于各因子 $L_i$，它们来自 $L^{sh}$），可知
$Y \times_X \Spec(\mathcal{O}_{X, x}^{sh}) \to \Spec(\mathcal{O}_{X, x}^{sh})$
是有限 \'etale 的。考察一般点可知其次数等于 $[L : K]$，故 (8) 成立。

\medskip\noindent
假设 (8) 成立。设 $x \in X$，并设概形论纤维 $Y_x$ 在 $x$ 上是次数
$\geq [L : K]$ 的 \'etale 概形。这意味着 $Y$ 有
$\geq [L : K]$ 个位于 $x$ 上方的几何点。我们将证明 $Y \to X$ 在 $x$ 的某个邻域上是有限
\'etale 的，这便证明了 (1)。
为此可以假设 $X = \Spec(R)$，点 $x$ 对应于素理想 $\mathfrak p \subset R$，
并且 $Y = \Spec(S)$。应用 More on Morphisms, Lemma
\ref{more-morphisms-lemma-etale-makes-integral-split}，可找到一个 \'etale 邻域
$(U, u) \to (X, x)$，使得
$Y \times_X U = V_1 \amalg \ldots \amalg V_m$，并且 $V_i$ 有唯一一个点 $v_i$，
它位于 $u$ 上方，而 $\kappa(v_i)/\kappa(u)$ 是纯不可分扩张。必要时缩小 $U$，
可以假设 $U$ 是一般点为 $\xi$ 的正规整概形（使用 Descent, Lemmas
\ref{descent-lemma-locally-finite-nr-irred-local-fppf}、
\ref{descent-lemma-normal-local-smooth}，以及 Properties, Lemma
\ref{properties-lemma-normal-locally-finite-nr-irreducibles}）。
由上面对几何点的说明可知 $m \geq [L : K]$。另一方面，由 More on Morphisms, Lemma
\ref{more-morphisms-lemma-normalization-smooth-localization} 可知
$\coprod V_i \to U$ 是 $U$ 在 $\Spec(L) \times_X U$ 中的正规化。因为
$K \subset \kappa(\xi)$ 是有限可分扩张，可以写成
$\Spec(L) \times_X U = \Spec(\prod_{i = 1, \ldots, n} L_i)$，
其中 $L_i/\kappa(\xi)$ 有限，且 $[L : K] = \sum [L_i : \kappa(\xi)]$。
由于 $V_j$ 对每个 $j$ 都非空，并且 $m \geq [L : K]$，所以
$m = n$，且 $[L_i : \kappa(\xi)] = 1$ 对所有 $i$ 都成立。于是
$V_j \to U$ 是同构，特别地是 \'etale 的，故 $Y \times_X U \to U$ 是
\'etale 的。由 Descent, Lemma \ref{descent-lemma-descending-property-etale}，
可知 $Y \to X$ 在 $U \to X$ 的像（即 $x$ 的一个开邻域）上是 \'etale 的。

\medskip\noindent
假设 $L/K$ 是 Galois 扩张且 (9) 成立。由 Lemma
\ref{pione-lemma-inertial-invariants-unramified}，$Y \to X$ 是 \'etale 的。
略去 (1) 蕴含 (9) 的证明。
\end{proof}

\noindent
对于离散赋值环的无限 Galois 扩张，还可以多说一点。为此引入如下记号。
称整数子集 $S \subset \mathbf{N}$ {\it 乘法有向}，如果 $1 \in S$，且对
$n, m \in S$，存在 $k \in S$ 使得 $n | k$ 且 $m | k$。在 $S$ 上定义偏序：
$n \geq_S m$ 当且仅当 $m | n$。给定域 $\kappa$，得到有限群的逆系统
$\{\mu_n(\kappa)\}_{n \in S}$，其转移映射为
$$
\mu_n(\kappa) \longrightarrow \mu_m(\kappa),\quad
\zeta \longmapsto \zeta^{n/m}
$$
其中 $n \geq_S m$。于是可以构造 profinite 群
$$
\lim_{n \in S} \mu_n(\kappa)
$$
注意，该极限是余滤过的（因为 $S$ 有向）。此构造关于 $\kappa$ 具有函子性。
特别地，$\text{Aut}(\kappa)$ 作用于这个 profinite 群。例如，若
$S = \{1, n\}$，则得到 $\mu_n(\kappa)$。若
$S = \{1, \ell, \ell^2, \ell^3, \ldots\}$，其中 $\ell$ 是不同于 $\kappa$
之特征的素数，则得到 $\lim_n \mu_{\ell^n}(\kappa)$；它有时称为
$\ell$-进 Tate 模，即 $\kappa$ 的乘法群的 Tate 模（参见 More on Algebra, Example
\ref{more-algebra-example-spectral-sequence-principal}）。

\begin{lemma}
\label{pione-lemma-structure-decomposition}
设 $A$ 是分式域为 $K$ 的离散赋值环。设 $L/K$ 是一个（可能无限的）Galois 扩张，
$B$ 是 $A$ 在 $L$ 中的整闭包，$\mathfrak m$ 是 $B$ 的一个极大理想。令
$G = \text{Gal}(L/K)$、
$D = \{\sigma \in G \mid \sigma(\mathfrak m) = \mathfrak m\}$，以及
$I = \{\sigma \in D \mid \sigma \bmod \mathfrak m =
\text{id}_{\kappa(\mathfrak m)}\}$.
分解群 $D$ 属于如下典范正合列：
$$
1 \to I \to D \to \text{Aut}(\kappa(\mathfrak m)/\kappa_A) \to 1
$$
惯性群 $I$ 属于如下典范正合列：
$$
1 \to P \to I \to I_t \to 1
$$
其中
\begin{enumerate}
\item 野惯性群 $P$ 是 $D$ 的正规子群；
\item $P$ 是 pro-$p$-群，如果 $\kappa_A$ 的特征为 $p > 1$；并且
$P = \{1\}$，如果 $\kappa_A$ 的特征为零；
\item 存在一个乘法有向子集 $S \subset \mathbf{N}$，使得
$\kappa(\mathfrak m)$ 含有一个本原 $n$ 次单位根，对每个 $n \in S$
（$S$ 的元素均与 $p$ 互素）；
\item 存在一个典范满射
$$
\theta_{can} : I \to \lim_{n \in S} \mu_n(\kappa(\mathfrak m))
$$
其核为 $P$，且满足
$\theta_{can}(\tau \sigma \tau^{-1}) = \tau(\theta_{can}(\sigma))$
（其中 $\tau \in D$、$\sigma \in I$），并诱导同构
$I_t \to \lim_{n \in S} \mu_n(\kappa(\mathfrak m))$；这里的商群称为温和惯性群。
\end{enumerate}
\end{lemma}

\begin{proof}
这基本上只是重述 More on Algebra, Section
\ref{more-algebra-section-ramification} 中所证明的有限 Galois 扩张结果。
映射 $D \to \text{Aut}(\kappa(\mathfrak m)/\kappa)$ 的满射性就是
More on Algebra, Lemma \ref{more-algebra-lemma-one-orbit-geometric-galois}。
由此得到第一个正合列。

\medskip\noindent
为构造第二个短正合列，令 $\Lambda$ 为有限 Galois 子扩张之集；也就是说，
$\lambda \in \Lambda$ 对应于 $L/L_\lambda/K$。令
$G_\lambda = \text{Gal}(L_\lambda/K)$。回顾 $G_\lambda$ 构成转移映射均为满射的有限群逆系统，
并且 $G = \lim_{\lambda \in \Lambda} G_\lambda$，见 Fields, Lemma
\ref{fields-lemma-infinite-galois-limit}。令 $B_\lambda$ 为 $A$ 在 $L_\lambda$ 中的整闭包。
再令 $\mathfrak m_\lambda = \mathfrak m \cap B_\lambda$，并以
$P_\lambda, I_\lambda, D_\lambda$ 分别表示 $\mathfrak m_\lambda$ 的
野惯性群、惯性群和分解群，见 More on Algebra, Lemma
\ref{more-algebra-lemma-galois-inertia}。对 $\lambda \geq \lambda'$，限制映射给出交换图
$$
\xymatrix{
P_\lambda \ar[d] \ar[r] &
I_\lambda \ar[d] \ar[r] &
D_\lambda \ar[d] \ar[r] &
G_\lambda \ar[d] \\
P_{\lambda'} \ar[r] &
I_{\lambda'} \ar[r] &
D_{\lambda'} \ar[r] &
G_{\lambda'}
}
$$
其中竖直映射均为满射，见 More on Algebra, Lemma
\ref{more-algebra-lemma-compare-inertia}。

\medskip\noindent
由定义立刻可知，$I = \lim I_\lambda$ 且 $D = \lim D_\lambda$；
这里使用的是上述同构 $G = \lim G_\lambda$。因为 $L = \colim L_\lambda$，所以
$B = \colim B_\lambda$ 且
$\kappa(\mathfrak m) = \colim \kappa(\mathfrak m_\lambda)$。
系统 $D_\lambda$ 的转移映射与映射
$D_\lambda \to \text{Aut}(\kappa(\mathfrak m_\lambda)/\kappa)$ 相容
（见 More on Algebra, Lemma \ref{more-algebra-lemma-compare-inertia}），
故映射 $D \to \text{Aut}(\kappa(\mathfrak m)/\kappa)$ 是诸映射
$D_\lambda \to \text{Aut}(\kappa(\mathfrak m_\lambda)/\kappa)$ 的极限。

\medskip\noindent
存在典范映射
$$
\theta_{\lambda, can} :
I_\lambda
\longrightarrow
\mu_{n_\lambda}(\kappa(\mathfrak m_\lambda))
$$
其中 $n_\lambda = |I_\lambda|/|P_\lambda|$，群
$\mu_{n_\lambda}(\kappa(\mathfrak m_\lambda))$ 的阶为 $n_\lambda$，并且
$\theta_{\lambda, can}(\tau \sigma \tau^{-1}) =
\tau(\theta_{\lambda, can}(\sigma))$ 对 $\tau \in D_\lambda$ 和
$\sigma \in I_\lambda$ 成立，而且得到交换图
$$
\xymatrix{
I_\lambda \ar[r]_-{\theta_{\lambda, can}} \ar[d] &
\mu_{n_\lambda}(\kappa(\mathfrak m_\lambda))
\ar[d]^{(-)^{n_\lambda/n_{\lambda'}}} \\
I_{\lambda'} \ar[r]^-{\theta_{\lambda', can}} &
\mu_{n_{\lambda'}}(\kappa(\mathfrak m_{\lambda'}))
}
$$
见 More on Algebra, Remark
\ref{more-algebra-remark-canonical-inertia-character}。

\medskip\noindent
令 $S \subset \mathbf{N}$ 为整数 $n_\lambda$ 的集合。由于 $\Lambda$ 有向，
可知 $S$ 是乘法有向的。对上面交换图中的映射 $\theta_{\lambda, can}$ 取极限，得到
$$
\theta_{can} : I \to \lim_{n \in S} \mu_n(\kappa(\mathfrak m)).
$$
此映射连续（略去一个小细节）。由于系统 $I_\lambda$ 的转移映射都是满射，
且 $\Lambda$ 有向，投影 $I \to I_\lambda$ 均为满射。对每个 $\lambda$，交换图
$$
\xymatrix{
I \ar[d] \ar[r]_-{\theta_{can}} &
\lim_{n \in S} \mu_n(\kappa(\mathfrak m)) \ar[d] \\
I_{\lambda} \ar[r]^-{\theta_{\lambda, can}} &
\mu_{n_\lambda}(\kappa(\mathfrak m_\lambda))
}
$$
成立。因此 $\theta_{can}$ 的像满射到有限群
$\mu_{n_\lambda}(\kappa(\mathfrak m)) =
\mu_{n_\lambda}(\kappa(\mathfrak m_\lambda))$；该群的阶为 $n_\lambda$
（见上文）。于是 $\theta_{can}$ 的像稠密。另一方面，$\theta_{can}$ 连续且源是
profinite 群。因此，一个拓扑论证表明 $\theta_{can}$ 是满射。

\medskip\noindent
等式 $\theta_{can}(\tau \sigma \tau^{-1}) = \tau(\theta_{can}(\sigma))$
对 $\tau \in D$、$\sigma \in I$ 成立；这是由映射 $\theta_{\lambda, can}$
的相应性质，以及映射 $D \to \text{Aut}(\kappa(\mathfrak m))$ 同各映射
$D_\lambda \to \text{Aut}(\kappa(\mathfrak m_\lambda))$ 的相容性得出的。
令 $P = \Ker(\theta_{can})$，这表明 $P$ 是 $D$ 的正规子群。令 $I_t = I/P$；
得到同构 $I_t \to \lim_{n \in S} \mu_n(\kappa(\mathfrak m))$，
这是由 $\theta_{can}$ 的满射性推出的。

\medskip\noindent
为完成证明，我们证明 $P = \lim P_\lambda$，这说明 $P$ 是 pro-$p$-群。回顾温和惯性群
$I_{\lambda, t} = I_\lambda/P_\lambda$ 的阶为 $n_\lambda$。由于转移映射
$P_\lambda \to P_{\lambda'}$ 均为满射，且 $\Lambda$ 有向，得到短正合列
$$
1 \to \lim P_\lambda \to I \to \lim I_{\lambda, t} \to 1
$$
（略去细节）。因为对每个 $\lambda$，映射 $\theta_{\lambda, can}$ 诱导同构
$I_{\lambda, t} \cong \mu_{n_\lambda}(\kappa(\mathfrak m))$
，所需结论随即得出。
\end{proof}

\begin{lemma}
\label{pione-lemma-structure-decomposition-separable-closure}
设 $A$ 是分式域为 $K$ 的离散赋值环，$K^{sep}$ 是 $K$ 的一个可分闭包，
$A^{sep}$ 是 $A$ 在 $K^{sep}$ 中的整闭包，而 $\mathfrak m^{sep}$ 是
$A^{sep}$ 的一个极大理想。令 $\mathfrak m = \mathfrak m^{sep} \cap A$、
$\kappa = A/\mathfrak m$，以及
$\overline{\kappa} = A^{sep}/\mathfrak m^{sep}$。于是
$\overline{\kappa}$ 是 $\kappa$ 的代数闭包。令
$G = \text{Gal}(K^{sep}/K)$、
$D = \{\sigma \in G \mid \sigma(\mathfrak m^{sep}) = \mathfrak m^{sep}\}$，以及
$I = \{\sigma \in D \mid \sigma \bmod \mathfrak m^{sep} =
\text{id}_{\kappa(\mathfrak m^{sep})}\}$.
分解群 $D$ 属于如下典范正合列：
$$
1 \to I \to D \to \text{Gal}(\kappa^{sep}/\kappa) \to 1
$$
其中 $\kappa^{sep} \subset \overline{\kappa}$ 是 $\kappa$ 的可分闭包。
惯性群 $I$ 属于如下典范正合列：
$$
1 \to P \to I \to I_t \to 1
$$
其中
\begin{enumerate}
\item 野惯性群 $P$ 是 $D$ 的正规子群；
\item $P$ 是 pro-$p$-群，如果 $\kappa_A$ 的特征为 $p > 1$；并且
$P = \{1\}$，如果 $\kappa_A$ 的特征为零；
\item 存在一个典范满射
$$
\theta_{can} : I \to \lim_{n\text{ prime to }p} \mu_n(\kappa^{sep})
$$
其核为 $P$，且满足
$\theta_{can}(\tau \sigma \tau^{-1}) = \tau(\theta_{can}(\sigma))$
（其中 $\tau \in D$、$\sigma \in I$），并诱导同构
$I_t \to \lim_{n\text{ prime to }p} \mu_n(\kappa^{sep})$。
\end{enumerate}
\end{lemma}

\begin{proof}
由 Lemma \ref{pione-lemma-get-algebraic-closure}，域 $\overline{\kappa}$ 是
$\kappa$ 的代数闭包。大部分断言立即由 Lemma
\ref{pione-lemma-structure-decomposition} 的相应部分得出。例如，因为
$\text{Aut}(\overline{\kappa}/\kappa) = \text{Gal}(\kappa^{sep}/\kappa)$，
便得到第一个正合列。
因此，唯一还需证明的断言是：若 $S$ 如 Lemma
\ref{pione-lemma-structure-decomposition} 中所述，则极限
$\lim_{n \in S} \mu_n(\overline{\kappa})$ 等于
$\lim_{n\text{ prime to }p} \mu_n(\kappa^{sep})$。为此，只需证明每个整数 $n$
只要与 $p$ 互素，就整除 $S$ 的某个元素。
令 $\pi \in A$ 是一个素元，并考虑分裂域 $L$，即多项式 $X^n - \pi$ 的分裂域。
因为该多项式可分，$L$ 是 $K$ 的有限 Galois 扩张。选取嵌入
$L \to K^{sep}$。注意，若 $B$ 是 $A$ 在 $L$ 中的整闭包，则
$A \to B_{\mathfrak m^{sep} \cap B}$ 的分歧指数可被 $n$ 整除
（因为 $\pi$ 有一个 $n$ 次根位于 $B$ 中；事实上该分歧指数等于 $n$，但此处用不到）。
于是，由 Lemma \ref{pione-lemma-structure-decomposition} 的证明中 $S$ 的构造可知，
$n$ 整除 $S$ 的某个元素。
\end{proof}








\section{几何基本群与算术基本群}
\label{pione-section-galois-action}

\noindent
本节研究如下比较：概形 $X$ 在域 $k$ 上的基本群与 $X_{\overline{k}}$
的基本群之间有何关系，其中 $\overline{k}$ 是 $k$ 的代数闭包。

\begin{lemma}
\label{pione-lemma-limit}
设 $I$ 是有向集。设 $X_i$ 是以 $I$ 为指标、转移态射均为仿射态射的拟紧拟分离概形逆系统。
依照 Limits, Section \ref{limits-section-limits}，令 $X = \lim X_i$。
则有范畴等价
$$
\colim \textit{F\'Et}_{X_i} = \textit{F\'Et}_X
$$
若 $X_i$ 对所有充分大的 $i$ 都连通，并且 $\overline{x}$ 是 $X$ 的一个几何点，则
$$
\pi_1(X, \overline{x}) = \lim \pi_1(X_i, \overline{x})
$$
\end{lemma}

\begin{proof}
该范畴等价由 Limits, Lemmas
\ref{limits-lemma-descend-finite-presentation},
\ref{limits-lemma-descend-finite-finite-presentation} 和
\ref{limits-lemma-descend-etale} 得出。第二个断言是范畴断言的形式推论。
\end{proof}

\begin{lemma}
\label{pione-lemma-perfection}
设 $k$ 是完美闭包为 $k^{perf}$ 的域。设 $X$ 是 $k$ 上的连通概形。
则 $X_{k^{perf}}$ 连通，且
$\pi_1(X_{k^{perf}}) \to \pi_1(X)$ 是同构。
\end{lemma}

\begin{proof}
这是基本群拓扑不变性的一个特例，见 Proposition
\ref{pione-proposition-universal-homeomorphism}。要看出
$\Spec(k^{perf}) \to \Spec(k)$ 是泛同胚，可以使用 Algebra, Lemma
\ref{algebra-lemma-radicial-integral-bijective}。
\end{proof}

\begin{lemma}
\label{pione-lemma-ses-field}
设 $k$ 是代数闭包为 $\overline{k}$ 的域。设 $X$ 是 $k$ 上拟紧且拟分离的概形。
若基变换 $X_{\overline{k}}$ 连通，则有短正合列
$$
1 \to \pi_1(X_{\overline{k}}) \to \pi_1(X) \to \pi_1(\Spec(k)) \to 1
$$
其中各项均为 profinite 拓扑群。
\end{lemma}

\begin{proof}
$\textit{F\'Et}_{\Spec(k)}$ 的连通对象形如
$\Spec(k') \to \Spec(k)$，其中 $k'/k$ 是有限可分扩张。于是
$X_{\Spec{k'}}$ 连通，因为态射 $X_{\overline{k}} \to X_{\Spec(k')}$ 是满射，
且按假设 $X_{\overline{k}}$ 连通。因此，由 Lemma
\ref{pione-lemma-functoriality-galois-surjective}，
$\pi_1(X) \to \pi_1(\Spec(k))$ 是满射。

\medskip\noindent
继续之前，注意可以假设 $k$ 是完美域。事实上，由 Lemma \ref{pione-lemma-perfection} 有
$\pi_1(X_{k^{perf}}) = \pi_1(X)$ 且
$\pi_1(\Spec(k^{perf})) = \pi_1(\Spec(k))$。

\medskip\noindent
显然，函子的复合
$\textit{F\'Et}_{\Spec(k)} \to \textit{F\'Et}_X \to
\textit{F\'Et}_{X_{\overline{k}}}$ 将对象送到 $X_{\Spec(\overline{k})}$ 的若干副本的不交并。
因此复合
$\pi_1(X_{\overline{k}}) \to \pi_1(X) \to \pi_1(\Spec(k))$
由 Lemma \ref{pione-lemma-composition-trivial} 可知是平凡同态。

\medskip\noindent
设 $U \to X$ 是有限 \'etale 态射，且 $U$ 连通。注意
$U \times_X X_{\overline{k}} = U_{\overline{k}}$。假设
$U_{\overline{k}} \to X_{\overline{k}}$ 有一个截面
$s : X_{\overline{k}} \to U_{\overline{k}}$。于是 $s(X_{\overline{k}})$ 是
$U_{\overline{k}}$ 的一个开连通分支。对
$\sigma \in \text{Gal}(\overline{k}/k)$，以 $s^\sigma$ 表示 $s$ 沿
$\Spec(\sigma)$ 的基变换。由于 $U_{\overline{k}} \to X_{\overline{k}}$ 是有限
\'etale 态射，它只有有限多个截面。因此
$$
\overline{T} = \bigcup s^\sigma(X_{\overline{k}})
$$
是有限并，且 $\overline{T}$ 是一个在
$\text{Gal}(\overline{k}/k)$ 作用下稳定的开闭子集。由 Varieties, Lemma
\ref{varieties-lemma-closed-fixed-by-Galois} 可知，$\overline{T}$ 是某个闭子集
$T \subset U$ 的原像。由于 $U_{\overline{k}} \to U$ 是开映射
（Morphisms, Lemma \ref{morphisms-lemma-scheme-over-field-universally-open}），
$T$ 也开。又因 $U$ 连通，故 $T = U$。因此 $U_{\overline{k}}$ 是
$X_{\overline{k}}$ 的若干副本的（有限）不交并。由 Lemma
\ref{pione-lemma-functoriality-galois-normal} 可知
$\pi_1(X_{\overline{k}}) \to \pi_1(X)$ 的像是正规子群。

\medskip\noindent
设 $V \to X_{\overline{k}}$ 是有限 \'etale 覆盖。回顾 $\overline{k}$ 是 $k$
的有限可分扩张之并。由 Lemma \ref{pione-lemma-limit}，可找到有限可分扩张 $k'/k$
和有限 \'etale 态射 $U \to X_{k'}$，使得
$V = X_{\overline{k}} \times_{X_{k'}} U =
U \times_{\Spec(k')} \Spec(\overline{k})$.
于是复合 $U \to X_{k'} \to X$ 是有限 \'etale 态射，并且
$U \times_{\Spec(k)} \Spec(\overline{k})$ 包含
$V = U \times_{\Spec(k')} \Spec(\overline{k})$ 作为开闭子概形。（这是因为
$\Spec(\overline{k})$ 是
$\Spec(k') \times_{\Spec(k)} \Spec(\overline{k})$ 的一个开闭子概形，
它由乘法映射 $k' \otimes_k \overline{k} \to \overline{k}$ 给出。）
由 Lemma \ref{pione-lemma-functoriality-galois-injective} 可知
$\pi_1(X_{\overline{k}}) \to \pi_1(X)$ 是单射。

\medskip\noindent
最后，必须证明：对任意有限 \'etale 态射 $U \to X$，若
$U_{\overline{k}}$ 是 $X_{\overline{k}}$ 的若干副本的不交并，则存在有限
\'etale 态射 $V \to \Spec(k)$ 和满射 $V \times_{\Spec(k)} X \to U$。
见 Lemma \ref{pione-lemma-functoriality-galois-ses}。像上面一样应用 Lemma
\ref{pione-lemma-limit}，可找到有限可分扩张 $k'/k$，使得存在同构
$U_{k'} \cong \coprod_{i = 1, \ldots, n} X_{k'}$.
因此令 $V = \coprod_{i = 1, \ldots, n} \Spec(k')$ 即得结论。
\end{proof}






\section{同伦正合列}
\label{pione-section-homotopy-exact-sequence}

\noindent
本节讨论如下结果。设 $f : X \to S$ 是有限表示的平坦且固有的态射，
其几何纤维均连通且约化。假设 $S$ 连通，并设 $\overline{s}$ 是 $S$
的一个几何点。则有基本群的正合列
$$
\pi_1(X_{\overline{s}}) \to \pi_1(X) \to \pi_1(S) \to 1
$$
见 Proposition \ref{pione-proposition-first-homotopy-sequence}。

\begin{lemma}
\label{pione-lemma-stein-factorization-etale}
\begin{reference}
\cite[Expose X, Proposition 1.2, p. 262]{SGA1}.
\end{reference}
设 $f : X \to S$ 是概形的固有态射。设 $X \to S' \to S$ 是 $f$ 的 Stein 分解，
见 More on Morphisms, Theorem
\ref{more-morphisms-theorem-stein-factorization-general}。若 $f$ 有限表示且平坦，
并具有几何约化纤维，则 $S' \to S$ 是有限 \'etale 态射。
\end{lemma}

\begin{proof}
这由 Derived Categories of Schemes, Lemma
\ref{perfect-lemma-proper-flat-geom-red} 以及 More on Morphisms, Theorem
\ref{more-morphisms-theorem-stein-factorization-general} 中的信息得出。
\end{proof}

\begin{proposition}
\label{pione-proposition-first-homotopy-sequence}
设 $f : X \to S$ 是有限表示的平坦且固有的态射，其几何纤维均连通且约化。
假设 $S$ 连通，并设 $\overline{s}$ 是 $S$ 的一个几何点。则有基本群的正合列
$$
\pi_1(X_{\overline{s}}) \to \pi_1(X) \to \pi_1(S) \to 1
$$
\end{proposition}

\begin{proof}
设 $Y \to X$ 是有限 \'etale 态射。考虑 Stein 分解
$$
\xymatrix{
Y \ar[d] \ar[r] & X \ar[d] \\
T \ar[r] & S
}
$$
它是 $Y \to S$ 的 Stein 分解。由 Lemma
\ref{pione-lemma-stein-factorization-etale}，态射 $T \to S$ 是有限 \'etale 的。
由此得到函子 $\textit{F\'Et}_X \to \textit{F\'Et}_S$。对任意有限 \'etale 态射
$U \to S$，态射 $Y \to U \times_S X$ 在 $X$ 上给出，当且仅当态射
$Y \to U$ 在 $S$ 上给出；这样的态射唯一地经过 Stein 分解，亦即对应于唯一的态射
$T \to U$
（这是因为 More on Morphisms, Lemma
\ref{more-morphisms-lemma-stein-universally-closed} 将 Stein 分解构造成相对正规化，
再应用 Morphisms, Lemma
\ref{morphisms-lemma-characterize-normalization} 的分解性质）。因此函子
$\textit{F\'Et}_X \to \textit{F\'Et}_S$ 与
$\textit{F\'Et}_S \to \textit{F\'Et}_X$ 互为伴随。注意，
$U \times_S X \to S$ 的 Stein 分解是 $U$，因为 $U \times_S X \to U$
的各纤维都几何连通。

\medskip\noindent
由上述讨论和 Categories, Lemma
\ref{categories-lemma-adjoint-fully-faithful}，可知
$\textit{F\'Et}_S \to \textit{F\'Et}_X$ 全忠实；也就是说，
$\pi_1(X) \to \pi_1(S)$ 是满射（Lemma
\ref{pione-lemma-functoriality-galois-surjective}）。

\medskip\noindent
显然，复合
$\textit{F\'Et}_S \to \textit{F\'Et}_X \to \textit{F\'Et}_{X_{\overline{s}}}$
将任意 $U$ 送到 $X_{\overline{s}}$ 的若干副本的不交并。因此由 Lemma
\ref{pione-lemma-composition-trivial}，
$\pi_1(X_{\overline{s}}) \to \pi_1(X) \to \pi_1(S)$ 是平凡同态。

\medskip\noindent
设 $Y \to X$ 是有限 \'etale 态射，$Y$ 连通，并且
$Y \times_X X_{\overline{s}}$ 含有一个连通分支 $Z$，它同构于
$X_{\overline{s}}$。考虑上述 Stein 分解 $T$。设
$\overline{t} \in T_{\overline{s}}$ 是与纤维 $Z$ 对应的点。注意 $T$ 连通
（因为它是连通概形的像），而且由上面的满射性，$T \times_S X$ 连通。
现在考虑分解
$$
\pi : Y \longrightarrow T \times_S X
$$
设 $\overline{x} \in X_{\overline{s}}$ 是任意闭点。注意
$\kappa(\overline{t}) = \kappa(\overline{s}) = \kappa(\overline{x})$
是代数闭域。于是 $\pi$ 在 $(\overline{t}, \overline{x})$ 上的纤维仅含一个点，
即唯一的点 $\overline{z} \in Z$；它对应于
$\overline{x} \in X_{\overline{s}}$，这一对应通过同构
$Z \to X_{\overline{s}}$ 给出。因此有限 \'etale 态射 $\pi$ 的次数为 $1$，
这是在 $(\overline{t}, \overline{x})$ 的一个邻域上。由于
$T \times_S X$ 连通，该态射处处次数为 $1$，从而 $Y \cong T \times_S X$。
所以 $Y \times_X X_{\overline{s}}$ 完全分裂。综合以上各点，Lemmas
\ref{pione-lemma-functoriality-galois-ses} 和
\ref{pione-lemma-functoriality-galois-normal} 均可应用，证明完成。
\end{proof}




\section{特化映射}
\label{pione-section-specialization-map}

\noindent
本节构造特化映射。设 $f : X \to S$ 是具有几何连通纤维的概形固有态射。
设 $s' \leadsto s$ 是 $S$ 中点的特化。设 $\overline{s}$ 和 $\overline{s}'$
分别是位于 $s$ 和 $s'$ 上方的几何点。则有特化映射
$$
sp : \pi_1(X_{\overline{s}'}) \longrightarrow \pi_1(X_{\overline{s}})
$$
此映射的构造如下。令 $A$ 为 $\mathcal{O}_{S, s}$ 关于
$\kappa(s) \subset \kappa(s)^{sep} \subset \kappa(\overline{s})$ 的严格
Hensel 化，见 Algebra, Definition \ref{algebra-definition-henselization}。
因为 $s' \leadsto s$，点 $s'$ 对应于 $\Spec(\mathcal{O}_{S, s})$ 的一个点，
故 $\Spec(A)$ 中至少有一个（也可能有许多个）位于 $s'$ 上方的点，
其剩余域是 $\kappa(s')$ 的可分代数扩张。由于
$\kappa(\overline{s}')$ 代数闭，可以选取一个态射
$\varphi : \overline{s}' \to \Spec(A)$，从而得到交换图
$$
\xymatrix{
\overline{s}' \ar[r]_-\varphi \ar[rd] &
\Spec(A) \ar[d] &
\overline{s} \ar[l] \ar[ld] \\
& S
}
$$
特化映射是复合
$$
\pi_1(X_{\overline{s}'}) \longrightarrow
\pi_1(X_A) =
\pi_1(X_{\kappa(s)^{sep}}) =
\pi_1(X_{\overline{s}})
$$
其中第一个等号由 Lemma \ref{pione-lemma-finite-etale-on-proper-over-henselian} 给出，
第二个等号由 Lemmas \ref{pione-lemma-perfection} 和
\ref{pione-lemma-finite-etale-invariant-over-proper} 得出。按构造，特化映射属于交换图
$$
\xymatrix{
\pi_1(X_{\overline{s}'}) \ar[rr]_{sp} \ar[rd] & &
\pi_1(X_{\overline{s}}) \ar[ld] \\
& \pi_1(X)
}
$$
只要 $X$ 连通。特化映射依赖于上面对
$\varphi : \overline{s}' \to \Spec(A)$ 的选择；若要表明这一点，便记作 $sp_\varphi$。

\begin{lemma}
\label{pione-lemma-specialization-map-base-change}
考虑概形的交换图
$$
\xymatrix{
Y \ar[d]_g \ar[r] & X \ar[d]^f \\
T \ar[r] & S
}
$$
其中 $f$ 和 $g$ 都是具有几何连通纤维的固有态射。设 $t' \leadsto t$ 是 $T$
中点的特化，并考虑上述特化映射
$sp : \pi_1(Y_{\overline{t}'}) \to \pi_1(Y_{\overline{t}})$。
则有特化映射的交换图
$$
\xymatrix{
\pi_1(Y_{\overline{t}'}) \ar[r]_{sp} \ar[d] & \pi_1(Y_{\overline{t}}) \ar[d] \\
\pi_1(X_{\overline{s}'}) \ar[r]^{sp} & \pi_1(X_{\overline{s}})
}
$$
其中 $\overline{s}$ 和 $\overline{s}'$ 分别是 $\overline{t}$ 和
$\overline{t}'$ 的像。
\end{lemma}

\begin{proof}
令 $B$ 为 $\mathcal{O}_{T, t}$ 关于
$\kappa(t) \subset \kappa(t)^{sep} \subset \kappa(\overline{t})$ 的严格 Hensel 化。
像构造特化映射时那样，选取
$\psi : \overline{t}' \to \Spec(B)$，它提升 $\overline{t}' \to T$。
以 $s$ 和 $s'$ 表示 $t$ 和 $t'$ 在 $S$ 中的像。
令 $A$ 为 $\mathcal{O}_{S, s}$ 关于
$\kappa(s) \subset \kappa(s)^{sep} \subset \kappa(\overline{s})$ 的严格 Hensel 化。
由于 $\kappa(\overline{s}) = \kappa(\overline{t})$，由严格 Hensel 化的函子性
（Algebra, Lemma \ref{algebra-lemma-strictly-henselian-functorial}），
得到属于下述交换图的环同态 $A \to B$
$$
\xymatrix{
\overline{t}' \ar[r]_-\psi \ar[d] & \Spec(B) \ar[d] \ar[r] & T \ar[d] \\
\overline{s}' \ar[r]^-\varphi & \Spec(A) \ar[r] & S
}
$$
这里直接取态射 $\varphi : \overline{s}' \to \Spec(A)$ 为复合
$\overline{t}' \to \Spec(B) \to \Spec(A)$。施行基变换得到交换图
$$
\xymatrix{
Y_{\overline{t}'} \ar[r] \ar[d] & Y_B \ar[d] \\
X_{\overline{s}'} \ar[r] & X_A
}
$$
由特化映射的构造，此图的交换性蕴含引理中交换图的交换性。
\end{proof}

\begin{lemma}
\label{pione-lemma-specialization-map-composition}
设 $f : X \to S$ 是具有几何连通纤维的固有态射。设
$s'' \leadsto s' \leadsto s$ 是 $S$ 中点的特化。特化映射的复合
$\pi_1(X_{\overline{s}''}) \to \pi_1(X_{\overline{s}'}) \to
\pi_1(X_{\overline{s}})$ 是特化映射
$\pi_1(X_{\overline{s}''}) \to \pi_1(X_{\overline{s}})$.
\end{lemma}

\begin{proof}
设 $\mathcal{O}_{S, s} \to A$ 是利用
$\kappa(s) \to \kappa(\overline{s})$ 构造的严格 Hensel 化。设
$A \to \kappa(\overline{s}')$ 是构造第一个特化映射时所用的映射。设
$\mathcal{O}_{S, s'} \to A'$ 是利用
$\kappa(s') \subset \kappa(\overline{s}')$ 构造的严格 Hensel 化。
由严格 Hensel 化的函子性，存在映射 $A \to A'$，使它同
$A' \to \kappa(\overline{s}')$ 的复合就是给定映射
（Algebra, Lemma \ref{algebra-lemma-map-into-henselian-colimit}）。
接着，设 $A' \to \kappa(\overline{s}'')$ 是构造第二个特化映射时所用的映射。
显然，第一与第二个特化映射的复合就是特化映射
$\pi_1(X_{\overline{s}''}) \to \pi_1(X_{\overline{s}})$
，它利用 $A \to A' \to \kappa(\overline{s}'')$ 构造。
\end{proof}

\noindent
设 $X \to S$ 是具有几何连通纤维的固有态射。设 $R$ 是分式域代数闭的严格
Hensel 赋值环，并设 $\Spec(R) \to S$ 是态射。设
$\eta, s \in \Spec(R)$ 分别为一般点和闭点。于是可以考虑特化映射
$$
sp_R : \pi_1(X_\eta) \to \pi_1(X_s)
$$
它属于基变换 $X_R/\Spec(R)$。注意，这是有意义的，因为 $\eta$ 和 $s$
的剩余域均代数闭。

\begin{lemma}
\label{pione-lemma-specialization-map-valuation-ring}
设 $f : X \to S$ 是具有几何连通纤维的固有态射。设 $s' \leadsto s$ 是 $S$
中点的特化，并设
$sp : \pi_1(X_{\overline{s}'}) \to \pi_1(X_{\overline{s}})$
是特化映射。则存在分式域代数闭的严格 Hensel 赋值环 $R$，它位于 $S$ 上，
使 $sp$ 同构于上面定义的 $sp_R$。
\end{lemma}

\begin{proof}
设 $\mathcal{O}_{S, s} \to A$ 是利用
$\kappa(s) \to \kappa(\overline{s})$ 构造的严格 Hensel 化。设
$A \to \kappa(\overline{s}')$ 是构造 $sp$ 时所用的映射。设
$R \subset \kappa(\overline{s}')$ 是分式域为 $\kappa(\overline{s}')$、
且支配 $A$ 之像的赋值环。见 Algebra, Lemma \ref{algebra-lemma-dominate}。
例如，由 Lemma \ref{pione-lemma-normal-local-domain-separablly-closed-fraction-field}
和 Algebra, Lemma \ref{algebra-lemma-valuation-ring-normal} 可知 $R$ 是严格
Hensel 环。引理随即成立。
\end{proof}

\noindent
设 $X \to S$ 是具有几何连通纤维的固有态射。设 $R$ 是严格 Hensel 离散赋值环，
并设 $\Spec(R) \to S$ 是态射。设 $\eta, s \in \Spec(R)$ 分别为一般点和闭点。
于是可以考虑特化映射
$$
sp_R : \pi_1(X_{\overline{\eta}}) \to \pi_1(X_s)
$$
它属于基变换 $X_R/\Spec(R)$。注意，这是有意义的，因为 $s$ 的剩余域代数闭。

\begin{lemma}
\label{pione-lemma-specialization-map-discrete-valuation-ring}
设 $f : X \to S$ 是具有几何连通纤维的固有态射。设 $s' \leadsto s$ 是 $S$
中点的特化，并设
$sp : \pi_1(X_{\overline{s}'}) \to \pi_1(X_{\overline{s}})$
是特化映射。若 $S$ 是 Noether 概形，则存在严格 Hensel 离散赋值环 $R$，
它位于 $S$ 上，使得 $sp$ 同构于上面定义的 $sp_R$。
\end{lemma}

\begin{proof}
设 $\mathcal{O}_{S, s} \to A$ 是利用
$\kappa(s) \to \kappa(\overline{s})$ 构造的严格 Hensel 化。设
$A \to \kappa(\overline{s}')$ 是构造 $sp$ 时所用的映射。设
$R \subset \kappa(\overline{s}')$ 是支配 $A$ 之像的离散赋值环，见 Algebra, Lemma
\ref{algebra-lemma-exists-dvr}。选取域的交换图
$$
\xymatrix{
\kappa(\overline{s}) \ar[r] & k \\
A/\mathfrak m_A \ar[r] \ar[u] & R/\mathfrak m_R \ar[u]
}
$$
其中 $k$ 代数闭。设 $R^{sh}$ 是 $R$ 的严格 Hensel 化，它利用 $R \to k$ 构造。
由 Algebra, Lemma \ref{algebra-lemma-strictly-henselian-functorial}，
得到映射 $A \to R^{sh}$。由 More on Algebra, Lemma
\ref{more-algebra-lemma-henselization-dvr}，环 $R^{sh}$ 是离散赋值环。
以 $\eta, o$ 表示 $\Spec(R^{sh})$ 的一般点和闭点。因为概形图
$$
\xymatrix{
\overline{\eta} \ar[d] \ar[r] & \Spec(R^{sh}) \ar[d] &
\Spec(k) \ar[d] \ar[l] \\
\overline{s}' \ar[r] & \Spec(A) & \overline{s} \ar[l]
}
$$
交换，所以得到交换图
$$
\xymatrix{
\pi_1(X_{\overline{\eta}}) \ar[d] \ar[r]_{sp_{R^{sh}}} & \pi_1(X_o) \ar[d] \\
\pi_1(X_{\overline{s}'}) \ar[r]^{sp} & X_{\overline{s}}
}
$$
按这些映射的构造，这是特化映射的交换图。由于竖直箭头均为同构
（Lemma \ref{pione-lemma-finite-etale-invariant-over-proper}），引理得证。
\end{proof}








\section{限制到闭子概形}
\label{pione-section-lefschetz}

\noindent
本节证明关于限制函子
$$
\textit{F\'Et}_X \longrightarrow \textit{F\'Et}_Y,\quad
U \longmapsto V = U \times_X Y
$$
的一些结果，其中 $X$ 是概形，而 $Y$ 是闭子概形。利用基本群的拓扑不变性，
可以把对此函子的研究同有限局部自由模上的完备化函子联系起来。

\medskip\noindent
以下引理使用 Cohomology of Schemes, Section
\ref{coherent-section-coherent-formal} 中定义的凝聚形式模概念。给定 Noether 概形和拟凝聚理想层
$\mathcal{I} \subset \mathcal{O}_X$，若对象 $(\mathcal{F}_n)$ 属于
$\textit{Coh}(X, \mathcal{I})$，且每个 $\mathcal{F}_n$ 都是有限局部自由
$\mathcal{O}_X/\mathcal{I}^n$-模，就称它{\it 有限局部自由}。

\begin{lemma}
\label{pione-lemma-restriction-fully-faithful}
设 $X$ 是 Noether 概形，$Y \subset X$ 是理想层为
$\mathcal{I} \subset \mathcal{O}_X$ 的闭子概形。假设完备化函子
$$
\textit{Coh}(\mathcal{O}_X)
\longrightarrow 
\textit{Coh}(X, \mathcal{I}),\quad
\mathcal{F} \longmapsto \mathcal{F}^\wedge
$$
在有限局部自由对象的全子范畴上全忠实（见上文）。则限制函子
$\textit{F\'Et}_X \to \textit{F\'Et}_Y$ 全忠实。
\end{lemma}

\begin{proof}
由于有限 \'etale 覆盖的范畴具有内部 Hom（Lemma
\ref{pione-lemma-internal-hom-finite-etale}），只需证明如下断言：给定 $U$，它是 $X$
上的有限 \'etale 概形，并给定态射 $t : Y \to U$，它在 $X$ 上，存在唯一的截面
$s : X \to U$ 使得 $t = s|_Y$。图示为
$$
\xymatrix{
& U \ar[d]^f \\
Y \ar[r] \ar[ru] & X \ar@{..>}@/^1em/[u]
}
$$
寻找虚线箭头 $s$ 等价于寻找 $\mathcal{O}_X$-代数映射
$$
s^\sharp : f_*\mathcal{O}_U \longrightarrow \mathcal{O}_X
$$
它模掉 $Y$ 的理想层后化为给定的代数映射
$t^\sharp : f_*\mathcal{O}_U \to \mathcal{O}_Y$。由 Lemma
\ref{pione-lemma-thickening}，可以将 $t$ 唯一提升为相容映射系统
$t_n : Y_n \to U$，从而得到映射
$$
\lim t_n^\sharp : f_*\mathcal{O}_U \longrightarrow \lim \mathcal{O}_{Y_n}
$$
它是 $X$ 上代数层的映射。由于 $f_*\mathcal{O}_U$ 是有限局部自由
$\mathcal{O}_X$-模，可知存在唯一的 $\mathcal{O}_X$-模映射
$\sigma : f_*\mathcal{O}_U \to \mathcal{O}_X$，其完备化为
$\lim t_n^\sharp$。为证明 $\sigma$ 是代数同态，需要检验交换图
$$
\xymatrix{
f_*\mathcal{O}_U \otimes_{\mathcal{O}_X} f_*\mathcal{O}_U
\ar[r] \ar[d]_{\sigma \otimes \sigma} &
f_*\mathcal{O}_U \ar[d]^\sigma \\
\mathcal{O}_X \otimes_{\mathcal{O}_X} \mathcal{O}_X \ar[r] &
\mathcal{O}_X
}
$$
成立。对每个 $n$，已知此图限制到 $Y_n$ 后交换；也就是说，施行完备化函子后该图交换。
因此由完备化函子的忠实性得出结论。
\end{proof}

\begin{lemma}
\label{pione-lemma-restriction-equivalence}
设 $X$ 是 Noether 概形，$Y \subset X$ 是理想层为
$\mathcal{I} \subset \mathcal{O}_X$ 的闭子概形。假设完备化函子
$$
\textit{Coh}(\mathcal{O}_X)
\longrightarrow 
\textit{Coh}(X, \mathcal{I}),\quad
\mathcal{F} \longmapsto \mathcal{F}^\wedge
$$
在有限局部自由对象的全子范畴上给出等价（见上文）。则限制函子
$\textit{F\'Et}_X \to \textit{F\'Et}_Y$ 是等价。
\end{lemma}

\begin{proof}
由 Lemma \ref{pione-lemma-restriction-fully-faithful}，限制函子全忠实。

\medskip\noindent
设 $U_1 \to Y$ 是有限 \'etale 态射。为完成证明，我们将说明 $U_1$
属于限制函子的本质像。

\medskip\noindent
对 $n \geq 1$，设 $Y_n$ 为第 $n$ 个 $Y$ 的无穷小邻域。由 Lemma
\ref{pione-lemma-thickening}，存在唯一的有限 \'etale 态射
$\pi_n : U_n \to Y_n$，其到 $Y = Y_1$ 的基变换恢复 $U_1 \to Y_1$。
考虑层 $\mathcal{F}_n = \pi_{n, *}\mathcal{O}_{U_n}$。可以并且将
$\mathcal{F}_n$ 视为一个 $\mathcal{O}_X$-模，它位于 $X$ 上；该模局部同构于
$(\mathcal{O}_X/f^{n + 1}\mathcal{O}_X)^{\oplus r}$。这个
$(\mathcal{F}_n)$ 是 $\textit{Coh}(X, \mathcal{I})$ 中的有限局部自由对象。
按假设，存在有限局部自由 $\mathcal{O}_X$-模 $\mathcal{F}$ 和相容同构系统
$$
\mathcal{F}/\mathcal{I}^n\mathcal{F} \to \mathcal{F}_n
$$
，它们是 $\mathcal{O}_X$-模同构。

\medskip\noindent
为在 $\mathcal{F}$ 上构造代数结构，考虑乘法映射
$\mathcal{F}_n \otimes_{\mathcal{O}_X} \mathcal{F}_n \to \mathcal{F}_n$
；它们来自 $\mathcal{F}_n = \pi_{n, *}\mathcal{O}_{U_n}$ 是代数层这一事实。
这些映射定义
$$
(\mathcal{F}\otimes_{\mathcal{O}_X} \mathcal{F})^\wedge
\longrightarrow
\mathcal{F}^\wedge
$$
，它是范畴 $\textit{Coh}(X, \mathcal{I})$ 中的映射。因此按假设，可以设存在映射
$\mu : \mathcal{F}\otimes_{\mathcal{O}_X} \mathcal{F} \to \mathcal{F}$
，它限制到 $Y_n$ 后给出上述乘法映射。由引理所述函子的忠实性，可知
$\mu$ 定义了交换 $\mathcal{O}_X$-代数结构，其底层模为 $\mathcal{F}$，并与
$\mathcal{F}_n$ 上给定的代数结构相容。令
$$
U = \underline{\Spec}_X((\mathcal{F}, \mu))
$$
得到有限局部自由概形 $\pi : U \to X$，其到 $Y$ 的限制同构于 $U_1$。
$\pi$ 的判别式是截面
$$
\det(Q_\pi) :
\mathcal{O}_X
\longrightarrow
\wedge^{top}(\pi_*\mathcal{O}_U)^{\otimes -2}
$$
的零点集；该截面构造于 Discriminants, Section
\ref{discriminant-section-discriminant}。由 Discriminants, Lemma
\ref{discriminant-lemma-discriminant}，它到 $Y_n$ 的限制对所有 $n$
都是同构，故它自身也是同构。再由 Discriminants, Lemma
\ref{discriminant-lemma-discriminant}，$\pi$ 是 \'etale 的。
\end{proof}

\begin{lemma}
\label{pione-lemma-restriction-fully-faithful-general}
设 $X$ 是 Noether 概形，$Y \subset X$ 是理想层为
$\mathcal{I} \subset \mathcal{O}_X$ 的闭子概形。令 $\mathcal{V}$ 为开子概形
$V \subset X$ 所成的集合，其中每个都包含 $Y$；以反包含关系排序。假设完备化函子
$$
\colim_\mathcal{V} \textit{Coh}(\mathcal{O}_V)
\longrightarrow
\textit{Coh}(X, \mathcal{I}),
\quad
\mathcal{F} \longmapsto \mathcal{F}^\wedge
$$
在有限局部自由对象所成的全子范畴上全忠实（见上文）。则限制函子
$\colim_\mathcal{V} \textit{F\'Et}_V \to \textit{F\'Et}_Y$
全忠实。
\end{lemma}

\begin{proof}
注意 $\mathcal{V}$ 是有向集，故这里的余极限具有 Categories, Section
\ref{categories-section-directed-colimits} 中所述的形式。其余论证几乎与 Lemma
\ref{pione-lemma-restriction-fully-faithful} 的证明完全相同；建议读者略去。

\medskip\noindent
由于有限 \'etale 覆盖的范畴具有内部 Hom（Lemma
\ref{pione-lemma-internal-hom-finite-etale}），只需证明如下断言：给定 $U$，它在
$V \in \mathcal{V}$ 上有限 \'etale，并给定态射 $t : Y \to U$，它在 $V$ 上；
存在 $V' \geq V$ 以及态射 $s : V' \to U$，后者在 $V$ 上，并且 $t = s|_Y$。图示为
$$
\xymatrix{
& & U \ar[d]^f \\
Y \ar[r] \ar[rru] & V' \ar@{..>}[ru] \ar[r] & V
}
$$
寻找虚线箭头 $s$ 等价于寻找 $\mathcal{O}_{V'}$-代数映射
$$
s^\sharp : f_*\mathcal{O}_U|_{V'} \longrightarrow \mathcal{O}_{V'}
$$
，它模掉 $Y$ 的理想层后化为给定的代数映射
$t^\sharp : f_*\mathcal{O}_U \to \mathcal{O}_Y$。由 Lemma
\ref{pione-lemma-thickening}，可以将 $t$ 唯一提升为相容映射系统
$t_n : Y_n \to U$，从而得到映射
$$
\lim t_n^\sharp : f_*\mathcal{O}_U \longrightarrow \lim \mathcal{O}_{Y_n}
$$
，它是 $V$ 上代数层的映射。注意 $f_*\mathcal{O}_U$ 是有限局部自由
$\mathcal{O}_V$-模。因此可取 $V' \geq V$ 以及映射
$\sigma : f_*\mathcal{O}_U|_{V'} \to \mathcal{O}_{V'}$
，其完备化为 $\lim t_n^\sharp$。为证明 $\sigma$ 是代数同态，需要检验交换图
$$
\xymatrix{
(f_*\mathcal{O}_U \otimes_{\mathcal{O}_V} f_*\mathcal{O}_U)|_{V'}
\ar[r] \ar[d]_{\sigma \otimes \sigma} &
f_*\mathcal{O}_U|_{V'} \ar[d]^\sigma \\
\mathcal{O}_{V'} \otimes_{\mathcal{O}_{V'}} \mathcal{O}_{V'} \ar[r] &
\mathcal{O}_{V'}
}
$$
成立。对每个 $n$，已知此图限制到 $Y_n$ 后交换；也就是说，施行完备化函子后该图交换。
因此由完备化函子的忠实性可知，存在 $V'' \geq V'$，使得
$\sigma|_{V''}$ 如所需是代数同态。
\end{proof}

\begin{lemma}
\label{pione-lemma-restriction-equivalence-general}
设 $X$ 是 Noether 概形，$Y \subset X$ 是理想层为
$\mathcal{I} \subset \mathcal{O}_X$ 的闭子概形。令 $\mathcal{V}$ 为开子概形
$V \subset X$ 所成的集合，其中每个都包含 $Y$；以反包含关系排序。假设完备化函子
$$
\colim_\mathcal{V} \textit{Coh}(\mathcal{O}_V)
\longrightarrow
\textit{Coh}(X, \mathcal{I}),
\quad
\mathcal{F} \longmapsto \mathcal{F}^\wedge
$$
在有限局部自由对象所成的全子范畴之间给出等价（见上文说明）。则限制函子
$$
\colim_\mathcal{V} \textit{F\'Et}_V \to \textit{F\'Et}_Y
$$
是等价。
\end{lemma}

\begin{proof}
注意 $\mathcal{V}$ 是有向集，故这里的余极限具有 Categories, Section
\ref{categories-section-directed-colimits} 中所述的形式。其余论证几乎与 Lemma
\ref{pione-lemma-restriction-equivalence} 的证明完全相同；建议读者略去。

\medskip\noindent
由 Lemma \ref{pione-lemma-restriction-fully-faithful-general}，限制函子全忠实。

\medskip\noindent
设 $U_1 \to Y$ 是有限 \'etale 态射。为完成证明，我们将说明 $U_1$
属于限制函子的本质像。

\medskip\noindent
对 $n \geq 1$，设 $Y_n$ 为第 $n$ 个 $Y$ 的无穷小邻域。由 Lemma
\ref{pione-lemma-thickening}，存在唯一的有限 \'etale 态射
$\pi_n : U_n \to Y_n$，其到 $Y = Y_1$ 的基变换恢复 $U_1 \to Y_1$。
考虑层 $\mathcal{F}_n = \pi_{n, *}\mathcal{O}_{U_n}$。可以并且将
$\mathcal{F}_n$ 视为一个 $\mathcal{O}_X$-模，它位于 $X$ 上；该模局部同构于
$(\mathcal{O}_X/f^{n + 1}\mathcal{O}_X)^{\oplus r}$.
这个 $(\mathcal{F}_n)$ 是 $\textit{Coh}(X, \mathcal{I})$ 中的有限局部自由对象。
按假设，存在 $V \in \mathcal{V}$、有限局部自由 $\mathcal{O}_V$-模
$\mathcal{F}$，以及相容同构系统
$$
\mathcal{F}/\mathcal{I}^n\mathcal{F} \to \mathcal{F}_n
$$
，它们是 $\mathcal{O}_V$-模同构。

\medskip\noindent
为在 $\mathcal{F}$ 上构造代数结构，考虑乘法映射
$\mathcal{F}_n \otimes_{\mathcal{O}_V} \mathcal{F}_n \to \mathcal{F}_n$
；它们来自 $\mathcal{F}_n = \pi_{n, *}\mathcal{O}_{U_n}$ 是代数层这一事实。
这些映射定义
$$
(\mathcal{F}\otimes_{\mathcal{O}_V} \mathcal{F})^\wedge
\longrightarrow
\mathcal{F}^\wedge
$$
，它是范畴 $\textit{Coh}(X, \mathcal{I})$ 中的映射。因此按假设，缩小
$V$ 后可假设存在映射
$\mu : \mathcal{F}\otimes_{\mathcal{O}_V} \mathcal{F} \to \mathcal{F}$
，它限制到 $Y_n$ 后给出上述乘法映射。必要时进一步缩小后，可假设 $\mu$
定义了交换 $\mathcal{O}_V$-代数结构，其底层模为 $\mathcal{F}$，并与
$\mathcal{F}_n$ 上给定的代数结构相容。令
$$
U = \underline{\Spec}_V((\mathcal{F}, \mu))
$$
，得到 $V$ 上的有限局部自由概形，其到 $Y$ 的限制同构于 $U_1$。由此
$U \to V$ 在位于 $Y$ 上方的所有点处都是 \'etale 的；参见 More on Morphisms,
Lemma \ref{more-morphisms-lemma-check-smoothness-on-infinitesimal-nbhds}。
因此再次缩小 $V$ 后，可假设 $U \to V$ 有限 \'etale。证明完成。
\end{proof}

\begin{lemma}
\label{pione-lemma-restriction-faithful}
设 $X$ 是概形，$Y \subset X$ 是闭子概形。若 $X$ 的每个连通分支都与
$Y$ 相交，则限制函子 $\textit{F\'Et}_X \to \textit{F\'Et}_Y$ 忠实。
\end{lemma}

\begin{proof}
设 $a, b : U \to U'$ 是两个 $X$ 上有限 \'etale 概形之间的态射，并且它们到
$Y$ 的限制相同。$U$ 的一个连通分支的像是 $X$ 的一个连通分支；这是将
Topology, Lemma \ref{topology-lemma-finite-fibre-connected-components}
应用于 $U \to X$ 到 $X$ 的一个连通分支上的限制所得。因此按假设，$U$
每个连通分支的像都与 $Y$ 相交。由 \'Etale Morphisms, Proposition
\ref{etale-proposition-equality}，可知 $a = b$ 在 $U$ 的每个连通分支上的限制成立。
由于 $a$ 与 $b$ 的等化子是 $U$ 的开子概形（因为 $U'$ 在 $X$ 上的对角是开的），
结论随即成立。
\end{proof}

\begin{lemma}
\label{pione-lemma-restriction-fully-faithful-special}
设 $X$ 是 Noether 概形，$Y \subset X$ 是闭子概形。令 $Y_n \subset X$ 为第
$n$ 个无穷小邻域，即 $Y$ 在 $X$ 中的相应邻域。假设以下条件之一成立：
\begin{enumerate}
\item $X$ 拟仿射，且
$\Gamma(X, \mathcal{O}_X) \to \lim \Gamma(Y_n, \mathcal{O}_{Y_n})$
是同构；或
\item $X$ 有丰沛可逆模 $\mathcal{L}$，且
$\Gamma(X, \mathcal{L}^{\otimes m}) \to
\lim \Gamma(Y_n, \mathcal{L}^{\otimes m}|_{Y_n})$
对所有 $m \gg 0$ 都是同构；或
\item 对每个有限局部自由 $\mathcal{O}_X$-模 $\mathcal{E}$，映射
$\Gamma(X, \mathcal{E}) \to \lim \Gamma(Y_n, \mathcal{E}|_{Y_n})$
是同构。
\end{enumerate}
则限制函子 $\textit{F\'Et}_X \to \textit{F\'Et}_Y$ 全忠实。
\end{lemma}

\begin{proof}
本引理形式地由 Lemma \ref{pione-lemma-restriction-fully-faithful} 和
Algebraic and Formal Geometry, Lemma
\ref{algebraization-lemma-completion-fully-faithful} 推出。
\end{proof}

\begin{lemma}
\label{pione-lemma-restriction-fully-faithful-general-special}
设 $X$ 是 Noether 概形，$Y \subset X$ 是闭子概形。令 $Y_n \subset X$ 为第
$n$ 个无穷小邻域，即 $Y$ 在 $X$ 中的相应邻域。令 $\mathcal{V}$ 为开子概形
$V \subset X$ 所成的集合，其中每个都包含 $Y$；以反包含关系排序。假设以下条件之一成立：
\begin{enumerate}
\item $X$ 拟仿射，且
$$
\colim_\mathcal{V} \Gamma(V, \mathcal{O}_V)
\longrightarrow
\lim \Gamma(Y_n, \mathcal{O}_{Y_n})
$$
是同构；或
\item $X$ 有丰沛可逆模 $\mathcal{L}$，且
$$
\colim_\mathcal{V} \Gamma(V, \mathcal{L}^{\otimes m})
\longrightarrow
\lim \Gamma(Y_n, \mathcal{L}^{\otimes m}|_{Y_n})
$$
对所有 $m \gg 0$ 都是同构；或
\item 对每个 $V \in \mathcal{V}$ 以及每个有限局部自由
$\mathcal{O}_V$-模 $\mathcal{E}$，映射
$$
\colim_{V' \geq V} \Gamma(V', \mathcal{E}|_{V'})
\longrightarrow
\lim \Gamma(Y_n, \mathcal{E}|_{Y_n})
$$
是同构。
\end{enumerate}
则函子
$$
\colim_\mathcal{V} \textit{F\'Et}_V \to \textit{F\'Et}_Y
$$
全忠实。
\end{lemma}

\begin{proof}
本引理形式地由 Lemma \ref{pione-lemma-restriction-fully-faithful-general} 和
Algebraic and Formal Geometry, Lemma
\ref{algebraization-lemma-completion-fully-faithful-general} 推出。
\end{proof}






\section{推出与基本群}
\label{pione-section-pushouts}

\noindent
主要结果如下。

\begin{lemma}
\label{pione-lemma-pushout-along-closed-immersion-and-integral}
在 More on Morphisms, Situation
\ref{more-morphisms-situation-pushout-along-closed-immersion-and-integral} 中，
例如若 $Z \to Y$ 与 $Z \to X$ 都是概形的闭浸入，则有范畴等价
$$
\textit{F\'Et}_{Y \amalg_Z X}
\longrightarrow
\textit{F\'Et}_Y
\times_{\textit{F\'Et}_Z}
\textit{F\'Et}_X
$$
\end{lemma}

\begin{proof}
由 More on Morphisms, Proposition
\ref{more-morphisms-proposition-pushout-along-closed-immersion-and-integral}，
该推出存在。此函子把推出上的有限 \'etale 概形 $U$ 送到基变换
$Y' = U \times_{Y \amalg_Z X} Y$ 与 $X' = U \times_{Y \amalg_Z X} X$，
以及自然同构 $Y' \times_Y Z \to X' \times_X Z$；此同构位于 $Z$ 上。
为证明此函子是等价，使用 More on Morphisms, Lemma
\ref{more-morphisms-lemma-pushout-functor-equivalence-flat} 构造拟逆函子。
唯一剩下要证明的是：给定分离、拟有限且 \'etale 的态射
$U \to Y \amalg_Z X$，并且 $X' \to X$ 与 $Y' \to Y$ 都有限，则
$U \to Y \amalg_Z X$ 有限。这既可由相应的代数事实（More on Algebra,
Lemma \ref{more-algebra-lemma-finite-module-over-fibre-product}）推出，也可由下述事实看出：
$$
X' \amalg Y' \to U
$$
满射，而 $X'$ 与 $Y'$ 在 $Y \amalg_Z X$ 上都是固有的（这里使用 More on Morphisms,
Proposition \ref{more-morphisms-proposition-pushout-along-closed-immersion-and-integral}
对推出的描述）。于是可应用 Morphisms, Lemma
\ref{morphisms-lemma-scheme-theoretic-image-is-proper}，得出 $U$ 在
$Y \amalg_Z X$ 上固有。由于拟有限且固有的态射是有限的（More on Morphisms,
Lemma \ref{more-morphisms-lemma-characterize-finite}），结论成立。
\end{proof}





\section{穿孔谱的有限 \'etale 覆盖，I}
\label{pione-section-pi1-punctured-spec}

\noindent
我们先证明一些 Lefschetz 型结果。

\begin{situation}
\label{pione-situation-local-lefschetz}
设 $(A, \mathfrak m)$ 是 Noether 局部环，且 $f \in \mathfrak m$。
置 $X = \Spec(A)$ 与 $X_0 = \Spec(A/fA)$，并令
$U = X \setminus \{\mathfrak m\}$ 与
$U_0 = X_0 \setminus \{\mathfrak m\}$ 分别为
$A$ 与 $A/fA$ 的穿孔谱。
\end{situation}

\noindent
回顾：对概形 $X$，在 $X$ 上有限
\'etale 的概形所成的范畴记为 $\textit{F\'Et}_X$；见
第 \ref{pione-section-finite-etale} 节。
在情形 \ref{pione-situation-local-lefschetz} 中，
我们将研究基变换函子
$$
\xymatrix{
\textit{F\'Et}_X \ar[d] \ar[r] & \textit{F\'Et}_U \ar[d] \\
\textit{F\'Et}_{X_0} \ar[r] & \textit{F\'Et}_{U_0}
}
$$
在许多情形下，右侧竖直箭头是忠实的。

\begin{lemma}
\label{pione-lemma-faithful}
在情形 \ref{pione-situation-local-lefschetz} 中，假设下列条件之一成立：
\begin{enumerate}
\item 对每个极小素理想都有 $\dim(A/\mathfrak p) \geq 2$，其中
$\mathfrak p \subset A$ 且 $f \not \in \mathfrak p$；或
\item $U$ 的每个连通分支都与 $U_0$ 相交。
\end{enumerate}
则
$$
\textit{F\'Et}_U \longrightarrow \textit{F\'Et}_{U_0},\quad
V \longmapsto V_0 = V \times_U U_0
$$
是忠实函子。
\end{lemma}

\begin{proof}
情形 (2) 立即由引理 \ref{pione-lemma-restriction-faithful} 得到。
假设 (1) 蕴含 $U$ 的每个不可约分支都与 $U_0$ 相交；见
Algebra, Lemma \ref{algebra-lemma-one-equation}。
因此 (1) 由 (2) 推出。
\end{proof}

\noindent
在证明更有内容的结果之前，需要先给出几个引理。

\begin{lemma}
\label{pione-lemma-fill-in-missing}
在情形 \ref{pione-situation-local-lefschetz} 中，设 $V \to U$ 是有限
态射。令 $A^\wedge$ 为 $\mathfrak m$-进完备化（对环 $A$ 而言），
令 $X' = \Spec(A^\wedge)$，并令 $U'$ 与 $V'$ 分别为
$U$ 与 $V$ 到 $X'$ 的基变换。若 $Y' \to X'$ 是满足
$V' = Y' \times_{X'} U'$ 的有限态射，则存在有限态射 $Y \to X$，
使得 $V = Y \times_X U$ 且 $Y' = Y \times_X X'$。
\end{lemma}

\begin{proof}
这是 More on Algebra, Proposition
\ref{more-algebra-proposition-equivalence} 的直接应用。
具体地，选取生成元 $f_1, \ldots, f_t$，它们生成 $\mathfrak m$。
对每个 $i$，记 $V \times_U D(f_i) = \Spec(B_i)$。
对 $1 \leq i, j \leq n$，得到同构
$\alpha_{ij} : (B_i)_{f_j} \to (B_j)_{f_i}$；它是
$A_{f_if_j}$-代数同构，因为二者的谱都表示
$V \times_U D(f_if_j)$。
记 $Y' = \Spec(B')$。由于 $V \times_U U' = Y \times_{X'} U'$，
得到同构 $\alpha_i : B'_{f_i} \to B_i \otimes_A A^\wedge$。
直接验证可知，$(B', B_i, \alpha_i, \alpha_{ij})$ 是
$\text{Glue}(A \to A^\wedge, f_1, \ldots, f_t)$ 的对象；见
More on Algebra, Remark \ref{more-algebra-remark-glueing-data}。
应用上引命题（并使用 More on Algebra, Remark
\ref{more-algebra-remark-formal-glueing-algebras} 得到代数结构），
可找到 $A$-代数 $B$，使得 $\text{Can}(B)$ 同构于
$(B', B_i, \alpha_i, \alpha_{ij})$。置 $Y = \Spec(B)$，
可见 $Y \to X$ 是配备有相容同构
$V \cong Y \times_X U$ 与 $Y' = Y \times_X X'$ 的态射，正合所需。
\end{proof}

\begin{lemma}
\label{pione-lemma-fully-faithful-henselian-completion}
在情形 \ref{pione-situation-local-lefschetz} 中，假设 $A$ 是 Hensel 环，
或更一般地，$(A, (f))$ 是 Hensel 对。
令 $A^\wedge$ 为 $\mathfrak m$-进完备化（对环 $A$ 而言），
令 $X' = \Spec(A^\wedge)$，并令 $U'$ 与 $U'_0$ 分别为
$U$ 与 $U_0$ 到 $X'$ 的基变换。若
$\textit{F\'Et}_{U'} \to \textit{F\'Et}_{U'_0}$ 全忠实，则
$\textit{F\'Et}_U \to \textit{F\'Et}_{U_0}$ 全忠实。
\end{lemma}

\begin{proof}
假设 $\textit{F\'Et}_{U'} \longrightarrow \textit{F\'Et}_{U'_0}$
全忠实。由于 $X' \to X$ 忠实平坦，立即可知函子
$V \to V_0 = V \times_U U_0$ 忠实。
由于有限 \'etale 覆盖的范畴具有内部 Hom
（引理 \ref{pione-lemma-internal-hom-finite-etale}），
只需证明以下断言：给定 $V$，它在 $U$ 上有限且 \'etale，则有
$$
\Mor_U(U, V) = \Mor_{U_0}(U_0, V_0)
$$
于是，假设给定态射 $s_0 : U_0 \to V_0$（在 $U_0$ 上），
我们将构造态射 $s : U \to V$（在 $U$ 上）。

\medskip\noindent
由假设，确实存在态射 $s' : U' \to V'$，其在 $V'_0$ 上的限制
为 $s'_0$，后者是 $s_0$ 的基变换。由于 $V' \to U'$ 有限且 \'etale，
这意味着 $V' = s'(U') \amalg W'$，其中 $W' \to U'$ 有限且 \'etale。
选取有限态射 $Z' \to X'$，使得
$W' = Z' \times_{X'} U'$。这可由下述形式的 Zariski 主定理得到：
More on Morphisms, Lemma
\ref{more-morphisms-lemma-quasi-finite-separated-pass-through-finite}
（略去一个小细节）。于是
$$
V' = s'(U') \amalg W' \longrightarrow X' \amalg Z' = Y'
$$
是开浸入，并且 $V' = Y' \times_{X'} U'$。
由引理 \ref{pione-lemma-fill-in-missing}，可找到有限态射 $Y \to X$，
使得 $V = Y \times_X U$ 且 $Y' = Y \times_X X'$。
记 $Y = \Spec(B)$，于是 $Y' = \Spec(B \otimes_A A^\wedge)$。
此时 $B \otimes_A A^\wedge$ 有幂等元 $e'$，它对应于开闭子概形
$X'$；后者包含于 $Y' = X' \amalg Z'$。

\medskip\noindent
先考虑 $A$ 为 Hensel 环的情形（稍为容易）。记 $\overline{e}$
为 $e'$ 在 $B \otimes_A \kappa(\mathfrak m) = B/\mathfrak mB$ 中的像；
它可提升为幂等元 $e$（属于 $B$），因为 $A$ 是 Hensel 环
（由 Algebra, Lemma \ref{algebra-lemma-characterize-henselian}，
$B$ 是局部环的乘积）。再由幂等元提升的唯一性可知，$e$ 映到 $e'$
（这里使用 $B \otimes_A A^\wedge$ 是局部环的乘积）。
令 $Y_1 \subset Y$ 为对应于 $e$ 的开闭子概形。于是
$Y_1 \times_X X' = s'(X')$，这蕴含 $Y_1 \to X$ 是同构
（由忠实平坦下降），并给出所需截面。

\medskip\noindent
再考虑 $(A, (f))$ 为 Hensel 对的情形。这里使用 $s'$ 是
$s'_0$ 的提升。具体地，令 $Y_{0, 1} \subset Y_0 = Y \times_X X_0$
为 $s_0(U_0) \subset V_0 = Y_0 \times_{X_0} U_0$ 的闭包。
由于 $X' \to X$ 平坦，基变换 $Y'_{0, 1} \subset Y'_0$
是 $s'_0(U'_0)$ 的闭包，因而等于 $X'_0 \subset Y'_0$
（见 Morphisms, Lemma
\ref{morphisms-lemma-flat-base-change-scheme-theoretic-image}）。
由于 $Y'_0 \to Y_0$ 是下沉的
（Morphisms, Lemma \ref{morphisms-lemma-fpqc-quotient-topology}），
可知 $Y_{0, 1}$ 在 $Y_0$ 中是开闭的。令 $e_0 \in B/fB$
为相应的幂等元。由 More on Algebra, Lemma
\ref{more-algebra-lemma-characterize-henselian-pair}，
可将 $e_0$ 提升为幂等元 $e \in B$。随后同前得出结论。
\end{proof}

\noindent
在情形 \ref{pione-situation-local-lefschetz} 中，限制函子
$\textit{F\'Et}_U \longrightarrow \textit{F\'Et}_{U_0}$
的全忠实性在相当温和的假设下即成立。特别地，这些假设通常并不蕴含
$U$ 是连通概形，但结论保证 $U$ 与 $U_0$ 的连通分支数相同。

\begin{lemma}
\label{pione-lemma-fully-faithful-simple}
在情形 \ref{pione-situation-local-lefschetz} 中，假设
\begin{enumerate}
\item[(a)] $A$ 有对偶化复形；
\item[(b)] $(A, (f))$ 是 Hensel 对；
\item[(c)] 下列条件之一成立：
\begin{enumerate}
\item[(i)] $A_f$ 满足 $(S_2)$，且 $X$ 的每个不包含于
$X_0$ 的不可约分支的维数都 $\geq 3$；或
\item[(ii)] 对每个素理想 $\mathfrak p \subset A$，若
$f \not \in \mathfrak p$，则有
$\text{depth}(A_\mathfrak p) + \dim(A/\mathfrak p) > 2$。
\end{enumerate}
\end{enumerate}
则限制函子
$\textit{F\'Et}_U \longrightarrow \textit{F\'Et}_{U_0}$
全忠实。
\end{lemma}

\begin{proof}
令 $A'$ 为关于 $\mathfrak m$ 的环 $A$ 的进完备化。我们将证明这些假设对
$A'$ 仍成立。对条件 (a) 与 (b)，这是显然的。条件 (c)(ii) 由
Local Cohomology, Lemma \ref{local-cohomology-lemma-change-completion}
在完备化下保持。其次，假设 (c)(i) 成立。由于 $A$ 普遍链状
（Dualizing Complexes, Lemma \ref{dualizing-lemma-universally-catenary}），
可知 $\Spec(A')$ 的每个不包含于 $V(f)$ 的不可约分支的维数都
$\geq 3$；见
More on Algebra, Proposition \ref{more-algebra-proposition-ratliff}。
由于 $A \to A'$ 平坦且具有 Gorenstein 纤维，
$A_f$ 满足 $(S_2)$ 蕴含 $A'_f$ 满足 $(S_2)$。
这里使用的参考结果为：
Dualizing Complexes, Section \ref{dualizing-section-formal-fibres},
More on Algebra, Section \ref{more-algebra-section-properties-formal-fibres},
以及 Algebra, Lemma \ref{algebra-lemma-Sk-goes-up}。
因此，由引理 \ref{pione-lemma-fully-faithful-henselian-completion}，
可设 $A$ 是完备 Noether 局部环。

\medskip\noindent
除其他假设外，再假设 $A$ 是完备局部环。由引理
\ref{pione-lemma-restriction-fully-faithful}，结论可由
Algebraic and Formal Geometry, Lemma
\ref{algebraization-lemma-fully-faithful-simple-two} 得到。
\end{proof}

\begin{lemma}
\label{pione-lemma-fully-faithful-minimal}
\begin{reference}
\cite[Corollary 1.11]{Bhatt-local}
\end{reference}
在情形 \ref{pione-situation-local-lefschetz} 中，假设
\begin{enumerate}
\item $H^1_\mathfrak m(A)$ 与 $H^2_\mathfrak m(A)$ 均被 $f$ 的某个幂零化；并且
\item $A$ 是 Hensel 环，或更一般地，$(A, (f))$ 是 Hensel 对。
\end{enumerate}
则限制函子
$\textit{F\'Et}_U \longrightarrow \textit{F\'Et}_{U_0}$
全忠实。
\end{lemma}

\begin{proof}
由引理 \ref{pione-lemma-fully-faithful-henselian-completion}，
可设 $A$ 是完备 Noether 局部环。
（这些假设会保持；使用 Dualizing Complexes, Lemma
\ref{dualizing-lemma-torsion-change-rings}。）
由引理 \ref{pione-lemma-restriction-fully-faithful}，结论可由
Algebraic and Formal Geometry, Lemma
\ref{algebraization-lemma-fully-faithful-alternative} 得到。
\end{proof}

\begin{lemma}
\label{pione-lemma-fully-faithful}
在情形 \ref{pione-situation-local-lefschetz} 中，假设 $A$ 的深度 $\geq 3$，
并且 $A$ 是 Hensel 环，或更一般地，$(A, (f))$ 是 Hensel 对。则限制函子
$\textit{F\'Et}_U \to \textit{F\'Et}_{U_0}$
全忠实。
\end{lemma}

\begin{proof}
深度假设蕴含
$H^1_\mathfrak m(A) = H^2_\mathfrak m(A) = 0$；见
Dualizing Complexes, Lemma \ref{dualizing-lemma-depth}。
因此可应用引理 \ref{pione-lemma-fully-faithful-minimal}。
\end{proof}



























\section{局部情形中的纯性，I}
\label{pione-section-local-purity}

\noindent
设 $(A, \mathfrak m)$ 是 Noether 局部环。置 $X = \Spec(A)$，
并令 $U = X \setminus \{\mathfrak m\}$ 为穿孔谱。若对
$(A, \mathfrak m)$，限制函子
$$
\textit{F\'Et}_X \longrightarrow \textit{F\'Et}_U
$$
本质满，则称{\it 纯性成立}。本节试图理解：从 $X$ 过渡到
超曲面 $X_0$（位于 $X$ 中）时，这个问题会如何变化。
换言之，我们研究穿孔谱基本群的一种局部 Lefschetz 性质。
这些结果将用于在证明局部纯性的主要结果时对维数作归纳；这些结果是
引理 \ref{pione-lemma-local-purity}、
命题 \ref{pione-proposition-purity-complete-intersection} 以及
命题 \ref{pione-proposition-purity-smooth-over-depth2}。

\begin{lemma}
\label{pione-lemma-sections-over-punctured-spec}
设 $(A, \mathfrak m)$ 是 Noether 局部环。置 $X = \Spec(A)$，
并令 $U = X \setminus \{\mathfrak m\}$。设 $\pi : Y \to X$
是有限态射，并且对所有闭点都有
$\text{depth}(\mathcal{O}_{Y, y}) \geq 2$，其中 $y \in Y$。
则 $Y$ 是 $B = \mathcal{O}_Y(\pi^{-1}(U))$ 的谱。
\end{lemma}

\begin{proof}
置 $V = \pi^{-1}(U)$，并以 $\pi' : V \to U$ 表示 $\pi$ 的限制。
考虑 $\mathcal{O}_X$-模映射
$$
\pi_*\mathcal{O}_Y \longrightarrow j_*\pi'_*\mathcal{O}_V
$$
其中 $j : U \to X$ 是包含态射。我们断言 Divisors, Lemma
\ref{divisors-lemma-depth-2-hartog} 可应用于此映射。若然，则
$B = \Gamma(Y, \mathcal{O}_Y)$，从而引理成立。令 $x \in X$
为闭点。只需证明
$\text{depth}((\pi_*\mathcal{O}_Y)_x) \geq 2$。
令 $y_1, \ldots, y_n \in Y$ 为映到 $x$ 的各点。
由 Algebra, Lemma \ref{algebra-lemma-depth-goes-down-finite}，
只需证明 $\text{depth}(\mathcal{O}_{Y, y_i}) \geq 2$
对 $i = 1, \ldots, n$ 成立。
这正是引理的假设，证明完毕。
\end{proof}

\begin{lemma}
\label{pione-lemma-reformulate-purity}
设 $(A, \mathfrak m)$ 是 Noether 局部环。置 $X = \Spec(A)$，
并令 $U = X \setminus \{\mathfrak m\}$。
设 $V$ 在 $U$ 上有限且 \'etale。假设 $A$ 的深度 $\geq 2$。
则下列条件等价：
\begin{enumerate}
\item 有 $V = Y \times_X U$，其中 $Y \to X$ 有限且 \'etale；
\item $B = \Gamma(V, \mathcal{O}_V)$ 在 $A$ 上有限且 \'etale。
\end{enumerate}
\end{lemma}

\begin{proof}
以 $\pi : V \to U$ 表示给定的有限 \'etale 态射。
假设 (1) 中的 $Y$ 存在。令 $x \in X$ 为对应于 $\mathfrak m$ 的点。
令 $y \in Y$ 为映到 $x$ 的点。我们断言
$\text{depth}(\mathcal{O}_{Y, y}) \geq 2$。
这是因为 $Y \to X$ 是 \'etale 的，因而
$A = \mathcal{O}_{X, x}$ 与 $\mathcal{O}_{Y, y}$ 具有相同深度
（Algebra, Lemma \ref{algebra-lemma-apply-grothendieck}）。
因此可应用引理 \ref{pione-lemma-sections-over-punctured-spec}，并有
$Y = \Spec(B)$。

\medskip\noindent
蕴含 (2) $\Rightarrow$ (1) 较为容易，细节从略。
\end{proof}

\begin{lemma}
\label{pione-lemma-reformulate-purity-normal}
设 $(A, \mathfrak m)$ 是 Noether 局部环。置 $X = \Spec(A)$，
并令 $U = X \setminus \{\mathfrak m\}$。假设 $A$ 正规且维数
$\geq 2$。函子
$$
\textit{F\'Et}_U \longrightarrow
\left\{
\begin{matrix}
\text{finite normal }A\text{-algebras }B\text{ such} \\
\text{that }\Spec(B) \to X\text{ is \'etale over }U
\end{matrix}
\right\},
\quad
V \longmapsto \Gamma(V, \mathcal{O}_V)
$$
是等价。此外，有 $V = Y \times_X U$，其中 $Y \to X$
有限且 \'etale，当且仅当 $B = \Gamma(V, \mathcal{O}_V)$
在 $A$ 上有限且 \'etale。
\end{lemma}

\begin{proof}
注意 $\text{depth}(A) \geq 2$，因为 $A$ 正规
（Serre 正规性判据，Algebra, Lemma
\ref{algebra-lemma-criterion-normal}）。
因此最后的陈述由引理 \ref{pione-lemma-reformulate-purity} 得到。
给定有限 \'etale 态射 $\pi : V \to U$，置
$B = \Gamma(V, \mathcal{O}_V)$。若能证明 $B$ 正规且在 $A$ 上有限，
便得到所显示的函子。由于有显然的拟逆函子，这也就是全部所需证明的内容。

\medskip\noindent
由于 $A$ 正规，概形 $V$ 正规
（Descent, Lemma \ref{descent-lemma-normal-local-smooth}）。
因此 $V$ 是有限多个整概形的不交并
（Properties, Lemma \ref{properties-lemma-normal-Noetherian}）。
所以可设 $V$ 整。在这种情形下，函数域 $L$（属于 $V$）
（Morphisms, Section \ref{morphisms-section-rational-maps}）
是 $A$ 的分式域的有限可分扩张
（这是通过考察 $V \to U$ 的一般纤维并使用 Morphisms, Lemma
\ref{morphisms-lemma-etale-over-field} 得到的）。
由 Algebra, Lemma
\ref{algebra-lemma-Noetherian-normal-domain-finite-separable-extension}，
整闭包 $B' \subset L$（即 $A$ 在 $L$ 中的整闭包）在 $A$ 上有限。
由 More on Algebra, Lemma
\ref{more-algebra-lemma-integral-closure-reflexive}，$B'$ 是自反的
$A$-模；进而由 More on Algebra, Lemma
\ref{more-algebra-lemma-reflexive-over-normal} 可知
$\text{depth}_A(B') \geq 2$。

\medskip\noindent
设 $f \in \mathfrak m$。则
$B_f = \Gamma(V \times_U D(f), \mathcal{O}_V)$
（Properties, Lemma \ref{properties-lemma-invert-f-sections}）。
所以 $B'_f = B_f$，因为 $B_f$ 正规（见上文）、在 $A_f$ 上有限，
且分式域为 $L$。于是 $V = \Spec(B') \times_X U$。
由引理 \ref{pione-lemma-sections-over-punctured-spec} 可得 $B = B'$；
这里将该引理应用于 $\Spec(B') \to X$。该引理适用，是因为局部化
$B'_{\mathfrak m'}$（它是 $B'$ 在极大理想
$\mathfrak m' \subset B'$ 处的局部化，而这些理想位于
$\mathfrak m$ 上方）的深度
$\geq 2$；这由 Algebra, Lemma
\ref{algebra-lemma-depth-goes-down-finite} 以及上一段关于深度的说明得到。
\end{proof}

\begin{lemma}
\label{pione-lemma-purity-and-completion}
设 $(A, \mathfrak m)$ 是 Noether 局部环。置 $X = \Spec(A)$，
并令 $U = X \setminus \{\mathfrak m\}$。设 $V$ 在 $U$ 上有限且
\'etale。令 $A^\wedge$ 为 $\mathfrak m$-进完备化（对环 $A$ 而言），
令 $X' = \Spec(A^\wedge)$，并令 $U'$ 与 $V'$ 分别为
$U$ 与 $V$ 到 $X'$ 的基变换。下列条件等价：
\begin{enumerate}
\item 有 $V = Y \times_X U$，其中 $Y \to X$ 有限且 \'etale；并且
\item 有 $V' = Y' \times_{X'} U'$，其中 $Y' \to X'$
有限且 \'etale。
\end{enumerate}
\end{lemma}

\begin{proof}
蕴含 (1) $\Rightarrow$ (2) 可由对一个解 $Y \to X$ 作基变换得到。
设 $Y' \to X'$ 如 (2) 所述。由引理 \ref{pione-lemma-fill-in-missing}，
可找到有限态射 $Y \to X$，使得 $V = Y \times_X U$ 且
$Y' = Y \times_X X'$。由下降可知 $Y \to X$ 有限且 \'etale
（Algebra, Lemmas \ref{algebra-lemma-descend-properties-modules} 以及
\ref{algebra-lemma-etale}）。证明完毕。
\end{proof}

\noindent
以下两个引理的要点在于，这些假设并不强制 $A$ 的深度
$\geq 3$。例如，若 $A$ 是维数 $\geq 3$ 的完备正规局部整环，
且 $f \in \mathfrak m$ 非零，则这些假设成立。

\begin{lemma}
\label{pione-lemma-lift-simple}
在情形 \ref{pione-situation-local-lefschetz} 中，设 $V$ 在 $U$ 上有限且
\'etale。假设
\begin{enumerate}
\item[(a)] $A$ 有对偶化复形；
\item[(b)] $(A, (f))$ 是 Hensel 对；
\item[(c)] 下列条件之一成立：
\begin{enumerate}
\item[(i)] $A_f$ 满足 $(S_2)$，且 $X$ 的每个不包含于
$X_0$ 的不可约分支的维数都 $\geq 3$；或
\item[(ii)] 对每个素理想 $\mathfrak p \subset A$，若
$f \not \in \mathfrak p$，则
$\text{depth}(A_\mathfrak p) + \dim(A/\mathfrak p) > 2$。
\end{enumerate}
\item[(d)] $V_0 = V \times_U U_0$ 等于 $Y_0 \times_{X_0} U_0$，
其中 $Y_0 \to X_0$ 有限且 \'etale。
\end{enumerate}
则有 $V = Y \times_X U$，其中 $Y \to X$ 有限且 \'etale。
\end{lemma}

\begin{proof}
使用引理 \ref{pione-lemma-purity-and-completion}，归约到完备情形。
（这些假设会保持；见引理 \ref{pione-lemma-fully-faithful-simple} 的证明。）

\medskip\noindent
在完备情形下，由 More on Algebra, Lemma
\ref{more-algebra-lemma-finite-etale-equivalence}，可将
$Y_0 \to X_0$ 提升为有限 \'etale 态射 $Y \to X$；注意，由
More on Algebra, Lemma \ref{more-algebra-lemma-complete-henselian}，
$(A, fA)$ 是 Hensel 对。随后可使用引理
\ref{pione-lemma-fully-faithful-simple} 看出 $V$ 同构于
$Y \times_X U$，证明完毕。
\end{proof}

\begin{lemma}
\label{pione-lemma-lift-purity-general}
在情形 \ref{pione-situation-local-lefschetz} 中，设 $V$ 在 $U$ 上有限且
\'etale。假设
\begin{enumerate}
\item $H^1_\mathfrak m(A)$ 与 $H^2_\mathfrak m(A)$ 均被 $f$ 的某个幂零化；
\item $V_0 = V \times_U U_0$ 等于 $Y_0 \times_{X_0} U_0$，
其中 $Y_0 \to X_0$ 有限且 \'etale。
\end{enumerate}
则有 $V = Y \times_X U$，其中 $Y \to X$ 有限且 \'etale。
\end{lemma}

\begin{proof}
使用引理 \ref{pione-lemma-purity-and-completion}，归约到完备情形。
（这些假设会保持；使用 Dualizing Complexes, Lemma
\ref{dualizing-lemma-torsion-change-rings}。）

\medskip\noindent
在完备情形下，由 More on Algebra, Lemma
\ref{more-algebra-lemma-finite-etale-equivalence}，可将
$Y_0 \to X_0$ 提升为有限 \'etale 态射 $Y \to X$；注意，由
More on Algebra, Lemma \ref{more-algebra-lemma-complete-henselian}，
$(A, fA)$ 是 Hensel 对。随后可使用引理
\ref{pione-lemma-fully-faithful-minimal} 看出 $V$ 同构于
$Y \times_X U$，证明完毕。
\end{proof}

\begin{lemma}
\label{pione-lemma-lift-purity}
在情形 \ref{pione-situation-local-lefschetz} 中，设 $V$ 在 $U$ 上有限且
\'etale。假设
\begin{enumerate}
\item $A$ 的深度 $\geq 3$；
\item $V_0 = V \times_U U_0$ 等于 $Y_0 \times_{X_0} U_0$，
其中 $Y_0 \to X_0$ 有限且 \'etale。
\end{enumerate}
则有 $V = Y \times_X U$，其中 $Y \to X$ 有限且 \'etale。
\end{lemma}

\begin{proof}
深度假设蕴含
$H^1_\mathfrak m(A) = H^2_\mathfrak m(A) = 0$；见
Dualizing Complexes, Lemma \ref{dualizing-lemma-depth}。
因此可应用引理 \ref{pione-lemma-lift-purity-general}。
\end{proof}










\section{分歧轨迹的纯性}
\label{pione-section-purity}

\noindent
我们将使用有限局部自由态射的判别式。见《判别式》Section
\ref{discriminant-section-discriminant}。

\begin{lemma}
\label{pione-lemma-find-point-codim-1}
设 $(A, \mathfrak m)$ 是满足 $\dim(A) \geq 1$ 的 Noether 局部环。
设 $f \in \mathfrak m$。则存在 $\mathfrak p \in V(f)$，使得
$\dim(A_\mathfrak p) = 1$。
\end{lemma}

\begin{proof}
对 $\dim(A)$ 作归纳。若 $\dim(A) = 1$，则取 $\mathfrak p = \mathfrak m$
即可。若 $\dim(A) > 1$，则设 $Z \subset \Spec(A)$ 是维数 $> 1$ 的一个
不可约分支。于是 $V(f) \cap Z$ 的维数 $> 0$
（《代数》引理 \ref{algebra-lemma-one-equation}）。选取素理想
$\mathfrak q \in V(f) \cap Z$，$\mathfrak q \not = \mathfrak m$，
它对应于 $A$ 的穿孔谱的一个闭点；由《性质》引理
\ref{properties-lemma-complement-closed-point-Jacobson}，这是可以做到的。
于是 $\mathfrak q$ 不是 $Z$ 的一般点。因此
$0 < \dim(A_\mathfrak q) < \dim(A)$ 且 $f \in \mathfrak q A_\mathfrak q$。
由关于维数的归纳，可以找到
$f \in \mathfrak p \subset A_\mathfrak q$，使得
$\dim((A_\mathfrak q)_\mathfrak p) = 1$。
于是 $\mathfrak p \cap A$ 满足要求。
\end{proof}

\begin{lemma}
\label{pione-lemma-ramification-quasi-finite-flat}
设 $f : X \to Y$ 是局部 Noether 概形的态射。设 $x \in X$。假设
\begin{enumerate}
\item $f$ 平坦，
\item $f$ 在 $x$ 处拟有限，
\item $x$ 不是 $X$ 的某个不可约分支的一般点，
\item 对特化 $x' \leadsto x$，若
$\dim(\mathcal{O}_{X, x'}) = 1$，则态射 $f$ 在 $x'$ 处非分歧。
\end{enumerate}
则 $f$ 在 $x$ 处是 \'etale 的。
\end{lemma}

\begin{proof}
注意，$f$ 非分歧的点所成集合与 $f$ 为 \'etale 的点所成集合相同，并且
这个集合是开的。见《态射》定义 \ref{morphisms-definition-unramified}
和 \ref{morphisms-definition-etale} 以及引理
\ref{morphisms-lemma-flat-unramified-etale}。
为了检验 $f$ 在 $x$ 处是 \'etale 的，我们可以在基底和目标上 \'etale
局部地工作（《下降》引理 \ref{descent-lemma-descending-property-etale} 和
\ref{descent-lemma-etale-etale-local-source}）。
因此，可以应用《态射补篇》引理
\ref{more-morphisms-lemma-etale-makes-quasi-finite-finite-at-point}
并假设 $f : X \to Y$ 有限，且 $x$ 是 $X$ 中位于 $y = f(x)$ 上方的唯一一点。
于是 $f$ 有限局部自由（《态射》引理
\ref{morphisms-lemma-finite-flat}）。

\medskip\noindent
假设 $f$ 有限局部自由，且 $x$ 是 $X$ 中位于 $y = f(x)$ 上方的唯一一点。
由《判别式》引理 \ref{discriminant-lemma-discriminant}，可得局部主闭子概形
$D_\pi \subset Y$，使得 $y' \in D_\pi$ 当且仅当存在 $x' \in X$，满足
$f(x') = y'$ 且 $f$ 在 $x'$ 处分歧。因此必须证明 $y \not \in D_\pi$。
假设 $y \in D_\pi$，以导出矛盾。

\medskip\noindent
由条件 (3)，有 $\dim(\mathcal{O}_{X, x}) \geq 1$。由《代数》引理
\ref{algebra-lemma-dimension-base-fibre-equals-total}，有
$\dim(\mathcal{O}_{X, x}) = \dim(\mathcal{O}_{Y, y})$。
由引理 \ref{pione-lemma-find-point-codim-1}，可以找到 $y' \in D_\pi$，它特化到 $y$，
且满足 $\dim(\mathcal{O}_{Y, y'}) = 1$。
选取 $x' \in X$，满足 $f(x') = y'$ 且 $f$ 在该点分歧。由于 $f$
有限，故它是闭态射，从而 $x' \leadsto x$。同前可得
$\dim(\mathcal{O}_{X, x'}) = \dim(\mathcal{O}_{Y, y'}) = 1$。
这与性质 (4) 矛盾。
\end{proof}

\begin{lemma}
\label{pione-lemma-local-purity}
设 $(A, \mathfrak m)$ 是维数为 $d \geq 2$ 的正则局部环。置
$X = \Spec(A)$ 和 $U = X \setminus \{\mathfrak m\}$。则
\begin{enumerate}
\item 函子 $\textit{F\'Et}_X \to \textit{F\'Et}_U$ 本质满，即纯性对 $A$
成立；
\item 任意有限同态 $A \to B$，若 $B$ 正规且在穿孔谱上诱导有限 \'etale
态射，则它是 \'etale 的。
\end{enumerate}
\end{lemma}

\begin{proof}
回顾由《代数》引理 \ref{algebra-lemma-regular-normal}，正则局部环是正规环。
因此，由引理 \ref{pione-lemma-reformulate-purity-normal}，(1) 与 (2) 等价。
我们对 $d$ 作归纳来证明该引理。

\medskip\noindent
考虑 $d = 2$ 的情形。此时 $A \to B$ 平坦。事实上，由《代数》命题
\ref{algebra-proposition-going-down-normal-integral}，$A \to B$ 满足降链性质。
再由《代数》引理 \ref{algebra-lemma-dimension-base-fibre-equals-total}，有
$\dim(B_{\mathfrak m'}) = 2$ 对所有极大理想 $\mathfrak m' \subset B$ 成立。
由《代数》引理 \ref{algebra-lemma-criterion-normal}，$B_{\mathfrak m'}$ 是
Cohen--Macaulay 环。因此，关键的一步是应用《代数》引理
\ref{algebra-lemma-CM-over-regular-flat}，得到 $A \to B_{\mathfrak m'}$ 平坦。
再由《代数》引理 \ref{algebra-lemma-flat-localization}，$A \to B$ 平坦。
于是可以应用引理 \ref{pione-lemma-ramification-quasi-finite-flat}（也可以直接应用较简单的
《判别式》引理 \ref{discriminant-lemma-discriminant} 进行论证），看出
$A \to B$ 是 \'etale 的。

\medskip\noindent
考虑 $d \geq 3$ 的情形。设 $V \to U$ 有限 \'etale。取
$f \in \mathfrak m_A$，$f \not \in \mathfrak m_A^2$。于是 $A/fA$ 是维数为
$d - 1 \geq 2$ 的正则局部环，见《代数》引理
\ref{algebra-lemma-regular-ring-CM}。设 $U_0$ 是 $A/fA$ 的穿孔谱，并置
$V_0 = V \times_U U_0$。由引理 \ref{pione-lemma-lift-purity}，只需证明
$V_0$ 位于 $\textit{F\'Et}_{\Spec(A/fA)} \to \textit{F\'Et}_{U_0}$ 的本质像中。
这由归纳假设得出。
\end{proof}

\begin{lemma}[分歧轨迹的纯性]
\label{pione-lemma-purity}
\begin{reference}
\cite{Nagata-Purity} 和 \cite[Exp. X, Thm. 3.1]{SGA1}
\end{reference}
\begin{history}
Zariski 最早在 \cite{Zariski-Purity} 中以几何形式陈述并证明了这一结果。
Nagata 对非完美基域的推广，作为同卷下一篇文章发表于
《美国国家科学院院刊》，见 \cite{Nagata-Remarks-Purity}。次年，Nagata 在
\cite{Nagata-Purity} 中证明了 Noether 局部环的结果。他的证明使用了 Chow
的一个结果，即完备局部整环的 Bertini 定理，见 \cite{Chow-Bertini}；
Bertini 定理的历史见 Kleiman 的历史性文章 \cite{Kleiman-Bertini}。
几年后，Auslander 找到了一个完全不同的证明，见 \cite{Auslander-Purity}。
\end{history}
设 $f : X \to Y$ 是局部 Noether 概形的态射。设 $x \in X$，并置
$y = f(x)$。假设
\begin{enumerate}
\item $\mathcal{O}_{X, x}$ 正规；
\item $\mathcal{O}_{Y, y}$ 正则；
\item $f$ 在 $x$ 处拟有限；
\item $\dim(\mathcal{O}_{X, x}) = \dim(\mathcal{O}_{Y, y}) \geq 1$；
\item 对特化 $x' \leadsto x$，若
$\dim(\mathcal{O}_{X, x'}) = 1$，则态射 $f$ 在 $x'$ 处非分歧。
\end{enumerate}
则 $f$ 在 $x$ 处是 \'etale 的。
\end{lemma}

\begin{proof}
我们对 $d = \dim(\mathcal{O}_{X, x}) = \dim(\mathcal{O}_{Y, y})$ 作归纳来证明该引理。

\medskip\noindent
先考虑较平凡的情形 $d = 1$。此时已假设 $f$ 在 $x$ 处非分歧，且
$\mathcal{O}_{Y, y}$ 是离散赋值环（《代数》引理
\ref{algebra-lemma-characterize-dvr}）。于是 $\mathcal{O}_{X, x}$ 在
$\mathcal{O}_{Y, y}$ 上平坦（否则该态射在 $x$ 处不会拟有限），从而看出
$f$ 在 $x$ 处平坦。由于平坦 $+$ 非分歧即为 \'etale，结论成立（略去一些细节）。

\medskip\noindent
考虑 $d \geq 2$ 的情形。我们将对 $d$ 作归纳，把问题约化到引理
\ref{pione-lemma-local-purity} 中讨论的情形。为了检验 $f$ 在 $x$ 处是 \'etale 的，
可以在基底和目标上 \'etale 局部地工作（《下降》引理
\ref{descent-lemma-descending-property-etale} 和
\ref{descent-lemma-etale-etale-local-source}）。因此，可以应用《态射补篇》引理
\ref{more-morphisms-lemma-etale-makes-quasi-finite-finite-at-point}
并假设 $f : X \to Y$ 有限，且 $x$ 是 $X$ 中位于 $y$ 上方的唯一一点。这里使用了
局部环的 \'etale 扩张不改变维数、正规性和正则性；见《代数补篇》Section
\ref{more-algebra-section-permanence-etale} 和《\'Etale 态射》Section
\ref{etale-section-properties-permanence}。

\medskip\noindent
接着，可以由 $\Spec(\mathcal{O}_{Y, y})$ 作基变换，并假设 $Y$ 是一个正则局部环的谱。
于是 $X = \Spec(\mathcal{O}_{X, x})$，因为 $X$ 的每一点都必定特化到 $x$。

\medskip\noindent
环同态 $\mathcal{O}_{Y, y} \to \mathcal{O}_{X, x}$ 有限，并且必为单射
（由维数相等）。由《代数》命题
\ref{algebra-proposition-going-down-normal-integral}，可知
$\mathcal{O}_{Y, y} \to \mathcal{O}_{X, x}$ 满足降链性质（并使用了由《代数》引理
\ref{algebra-lemma-regular-normal} 得出的正则环是正规环这一事实）。选取
$x' \in X$，$x' \not = x$，并令其像为 $y' = f(x')$。由于正规整环的局部化仍
正规，$\mathcal{O}_{X, x'}$ 正规。类似地，$\mathcal{O}_{Y, y'}$ 正则（见《代数》引理
\ref{algebra-lemma-localization-of-regular-local-is-regular}）。由《代数》引理
\ref{algebra-lemma-dimension-base-fibre-equals-total}（上面已经验证降链性质），有
$\dim(\mathcal{O}_{X, x'}) = \dim(\mathcal{O}_{Y, y'})$。这些维数当然严格小于
$d$，因为 $x' \not = x$；再由关于 $d$ 的归纳，得出 $f$ 在 $x'$ 处是 \'etale 的。

\medskip\noindent
于是归结到如下情形：存在 Noether 局部环之间的有限局部同态 $A \to B$，
这些环的维数为 $d \geq 2$，其中 $A$ 正则，$B$ 正规，并且它在穿孔谱上诱导有限 \'etale 态射
$V \to U$。我们的目标是证明 $A \to B$ 为 \'etale。这由引理
\ref{pione-lemma-local-purity} 得出，证明完毕。
\end{proof}

\noindent
下面的引理有时可用于找出有限 \'etale 态射能够延拓到的最大开子集。

\begin{lemma}
\label{pione-lemma-extend-S2}
设 $j : U \to X$ 是局部 Noether 概形的开浸入，且
$\text{depth}(\mathcal{O}_{X, x}) \geq 2$ 对 $x \not \in U$ 成立。
设 $\pi : V \to U$ 有限 \'etale。则
\begin{enumerate}
\item $\mathcal{B} = j_*\pi_*\mathcal{O}_V$ 是自反凝聚 $\mathcal{O}_X$-代数；置
$Y = \underline{\Spec}_X(\mathcal{B})$；
\item $Y \to X$ 是满足如下条件的唯一有限态射：$V = Y \times_X U$，且
$\text{depth}(\mathcal{O}_{Y, y}) \geq 2$ 对 $y \in Y \setminus V$ 成立；
\item $Y \to X$ 在 $y$ 处为 \'etale 当且仅当 $Y \to X$ 在 $y$ 处平坦；
\item $Y \to X$ 为 \'etale 当且仅当 $\mathcal{B}$ 作为 $\mathcal{O}_X$-模有限局部自由。
\end{enumerate}
此外，(a) $\mathcal{B}$ 和 $Y \to X$ 的构造与局部 Noether 概形的平坦态射
$X' \to X$ 所给的基变换可交换；(b) 若 $V' \to U'$ 是有限 \'etale 态射，
$U \subset U' \subset X$ 为开子概形，并且它限制为 $V \to U$（在 $U$ 上），
则存在唯一同构 $Y' \times_X U' = V'$（在 $U'$ 上）。
\end{lemma}

\begin{proof}
注意，$\pi_*\mathcal{O}_V$ 是有限局部自由 $\mathcal{O}_U$-模，特别地是自反模。
由《除子》引理 \ref{divisors-lemma-reflexive-S2-extend}，模
$j_*\pi_*\mathcal{O}_V$ 是 $X$ 上以 $\pi_*\mathcal{O}_V$ 为其在 $U$ 上限制的
唯一自反凝聚模。这证明了 (1)。

\medskip\noindent
由构造，$Y \times_X U = V$。由于 $\mathcal{B}$ 凝聚，可知 $Y \to X$ 有限。
由《除子》引理 \ref{divisors-lemma-reflexive-S2}，
$\text{depth}(\mathcal{B}_x) \geq 2$ 对 $x \in X \setminus U$ 成立。因此，由
《代数》引理 \ref{algebra-lemma-depth-goes-down-finite}，
$\text{depth}(\mathcal{O}_{Y, y}) \geq 2$ 对 $y \in Y \setminus V$ 成立。反之，假设
$\pi' : Y' \to X$ 是有限态射，满足 $V = Y' \times_X U$，且
$\text{depth}(\mathcal{O}_{Y', y'}) \geq 2$ 对 $y' \in Y' \setminus V$ 成立。
则 $\pi'_*\mathcal{O}_{Y'}$ 限制为 $\pi_*\mathcal{O}_V$（在 $U$ 上），并且
$\text{depth}((\pi'_*\mathcal{O}_{Y'})_x) \geq 2$ 对
$x \in X \setminus U$ 成立；这由《代数》引理
\ref{algebra-lemma-depth-goes-down-finite} 得出。
例如，由《除子》引理 \ref{divisors-lemma-depth-2-hartog}，
$\pi'_*\mathcal{O}_{Y'}$ 与 $j_*\pi_*\mathcal{O}_V$ 典范同构。这证明了 (2)。

\medskip\noindent
若 $Y \to X$ 在 $y$ 处为 \'etale，则 $Y \to X$ 在 $y$ 处平坦。反之，假设
$Y \to X$ 在 $y$ 处平坦。若 $y \in V$，则 $Y \to X$ 在 $y$ 处为 \'etale。
若 $y \not \in V$，则检验引理 \ref{pione-lemma-ramification-quasi-finite-flat} 的
(1)、(2)、(3)、(4) 成立，即可看出 $Y \to X$ 在 $y$ 处为 \'etale。
(1) 和 (2) 显然成立；由于 $\text{depth}(\mathcal{O}_{Y, y}) \geq 2$，(3) 也显然成立。
若 $y' \leadsto y$ 是特化且 $\dim(\mathcal{O}_{Y, y'}) = 1$，则
$y' \in V$，因为否则由《代数》引理 \ref{algebra-lemma-bound-depth}，
这个局部环的深度将为 $2$，矛盾。因此，$Y \to X$ 在 $y'$ 处为 \'etale，
从而引理 \ref{pione-lemma-ramification-quasi-finite-flat} 的 (4) 也成立。
这就完成了 (3) 的证明。

\medskip\noindent
(4) 由 (3) 以及如下事实得出：$((Y \to X)_*\mathcal{O}_Y)_x$ 是平坦
$\mathcal{O}_{X, x}$-模，当且仅当 $\mathcal{O}_{Y, y}$ 是平坦
$\mathcal{O}_{X, x}$-模；这里取遍所有 $y \in Y$，且它们映到 $x$。见《代数》引理
\ref{algebra-lemma-flat-localization}。这里还使用了 Noether 环上的有限平坦模是
有限局部自由模这一事实，见《代数》引理 \ref{algebra-lemma-finite-projective}
（以及《代数》引理
\ref{algebra-lemma-Noetherian-finite-type-is-finite-presentation}）。

\medskip\noindent
关于引理最后的断言，(a) 由平坦基变换得出，见《概形的上同调》引理
\ref{coherent-lemma-flat-base-change-cohomology}；把 (2) 中的唯一性应用于限制
$Y \times_X U'$，便得到 (b)。
\end{proof}

\begin{lemma}
\label{pione-lemma-extend-pure}
设 $j : U \to X$ 是 Noether 概形的开浸入，并且纯性对
$\mathcal{O}_{X, x}$ 成立；这里取遍所有 $x \not \in U$。则
$$
\textit{F\'Et}_X \longrightarrow \textit{F\'Et}_U
$$
本质满。
\end{lemma}

\begin{proof}
设 $V \to U$ 是有限 \'etale 态射。由 Noether 归纳，只需把 $V \to U$ 延拓为
到 $X$ 的一个严格更大开子集的有限 \'etale 态射。设 $x \in X \setminus U$
是 $X \setminus U$ 的一个不可约分支的一般点。则逆像 $U_x$，即 $U$ 在
$\Spec(\mathcal{O}_{X, x})$ 中的逆像，是 $\mathcal{O}_{X, x}$ 的穿孔谱。
由假设，$V_x = V \times_U U_x$ 是有限 \'etale 态射
$Y_x \to \Spec(\mathcal{O}_{X, x})$ 到 $U_x$ 的限制。由《极限》引理
\ref{limits-lemma-glueing-near-point}，可以找到开子概形 $U \subset U' \subset X$，
它包含 $x$，以及有限表示态射 $V' \to U'$，使得它在 $U$ 上的限制
恢复 $V \to U$，并且它在 $\Spec(\mathcal{O}_{X, x})$ 上的限制恢复 $Y_x$。
最后，态射 $V' \to U'$ 在可能把 $U'$ 缩小为更小的开子概形后是有限 \'etale 的，
见《极限》引理 \ref{limits-lemma-glueing-near-point-properties}。
\end{proof}
















\section{穿孔谱的有限 \'etale 覆盖，II}
\label{pione-section-pi1-punctured-spec-II}

\noindent
本节证明 Section \ref{pione-section-pi1-punctured-spec} 中所讨论内容的一些变式。
设有 Noether 局部环 $(A, \mathfrak m)$ 和 $f \in \mathfrak m$。置
$X = \Spec(A)$ 和 $X_0 = \Spec(A/fA)$，并令
$U = X \setminus \{\mathfrak m\}$ 和 $U_0 = X_0 \setminus \{\mathfrak m\}$
分别为 $A$ 和 $A/fA$ 的穿孔谱。所有这些恰如情形
\ref{pione-situation-local-lefschetz}。不同之处在于，我们将考虑限制函子
$$
\colim_{U_0 \subset U' \subset U\text{ 开}} \textit{F\'Et}_{U'}
\longrightarrow
\textit{F\'Et}_{U_0}
$$
换言之，我们不会试图把 $U_0$ 的有限 \'etale 覆盖提升到整个 $U$，而只提升到
某个开邻域 $U'$，即 $U_0$ 在 $U$ 中的一个开邻域。

\begin{lemma}
\label{pione-lemma-faithful-general}
在情形 \ref{pione-situation-local-lefschetz} 下，设 $U' \subset U$ 是包含 $U_0$ 的开子集。
假设对素理想 $\mathfrak p \subset A$，若点 $\mathfrak p \in U'$ 是极小点，且
$\mathfrak p \not \in U_0$，则
$\dim(A/\mathfrak p) \geq 2$。那么
$$
\textit{F\'Et}_{U'} \longrightarrow \textit{F\'Et}_{U_0},\quad
V' \longmapsto V_0 = V' \times_{U'} U_0
$$
是忠实函子。此外，存在满足该假设的 $U'$；而任意更小开子集
$U'' \subset U'$，若包含 $U_0$，也满足该假设。特别地，限制函子
$$
\colim_{U_0 \subset U' \subset U\text{ 开}} \textit{F\'Et}_{U'}
\longrightarrow
\textit{F\'Et}_{U_0}
$$
是忠实的。
\end{lemma}

\begin{proof}
由《代数》引理 \ref{algebra-lemma-one-equation} 可知，$V(\mathfrak p)$ 与
$U_0$ 相交，其中 $\mathfrak p$ 是 $A$ 的任意满足
$\dim(A/\mathfrak p) \geq 2$ 的素理想。因此，由引理
\ref{pione-lemma-restriction-faithful}，陈述中的这种 $U$
所对应的展示函子是忠实的。为了看出这样的 $U'$ 存在，注意若
$\mathfrak p \subset A$ 满足 $\mathfrak p \in U$、$\mathfrak p \not \in U_0$ 和
$\dim(A/\mathfrak p) = 1$，则 $\mathfrak p$ 对应于 $U$ 的一个闭点，因而
$V(\mathfrak p) \cap U_0 = \emptyset$。因此，可以取 $U'$ 为 $X$ 的那些不与
$U_0$ 相交且维数为 $1$ 的不可约分支之补集。
\end{proof}

\begin{lemma}
\label{pione-lemma-fully-faithful-general-better}
在情形 \ref{pione-situation-local-lefschetz} 下，假设
\begin{enumerate}
\item $A$ 有对偶化复形，且关于 $f$-进拓扑完备；
\item $X$ 的每个不包含于 $X_0$ 的不可约分支的维数都 $\geq 3$。
\end{enumerate}
则限制函子
$$
\colim_{U_0 \subset U' \subset U\text{ 开}} \textit{F\'Et}_{U'}
\longrightarrow
\textit{F\'Et}_{U_0}
$$
全忠实。
\end{lemma}

\begin{proof}
由基本群的拓扑不变性，可以在证明中把 $A$ 替换为其约化；见引理
\ref{pione-lemma-thickening}。于是，结论由引理
\ref{pione-lemma-restriction-fully-faithful-general} 和《代数化与形式几何》引理
\ref{algebraization-lemma-fully-faithful-general} 得出。
\end{proof}

\begin{lemma}
\label{pione-lemma-fully-faithful-general}
在情形 \ref{pione-situation-local-lefschetz} 下，假设
\begin{enumerate}
\item $A$ 关于 $f$-进拓扑完备；
\item $f$ 是非零因子；
\item $H^1_\mathfrak m(A/fA)$ 是有限 $A$-模。
\end{enumerate}
则限制函子
$$
\colim_{U_0 \subset U' \subset U\text{ 开}} \textit{F\'Et}_{U'}
\longrightarrow
\textit{F\'Et}_{U_0}
$$
全忠实。
\end{lemma}

\begin{proof}
这由引理 \ref{pione-lemma-restriction-fully-faithful-general} 和《代数化与形式几何》引理
\ref{algebraization-lemma-fully-faithful-general-alternative} 得出。
\end{proof}








\section{穿孔谱的有限 \'etale 覆盖，III}
\label{pione-section-pi1-punctured-spec-III}

\noindent
本节研究在情形 \ref{pione-situation-local-lefschetz} 下，限制函子
$$
\colim_{U_0 \subset U' \subset U\text{ 开}} \textit{F\'Et}_{U'}
\longrightarrow
\textit{F\'Et}_{U_0}
$$
何时为范畴等价。

\begin{lemma}
\label{pione-lemma-essentially-surjective-general-better}
在情形 \ref{pione-situation-local-lefschetz} 下，假设
\begin{enumerate}
\item $A$ 有对偶化复形，且关于 $f$-进拓扑完备；
\item 下列条件之一成立：
\begin{enumerate}
\item $A_f$ 满足 $(S_2)$，且 $X$ 的每个不包含于 $X_0$ 的不可约分支的维数都
$\geq 4$；或
\item 若 $\mathfrak p \not \in V(f)$ 且
$V(\mathfrak p) \cap V(f) \not = \{\mathfrak m\}$，则
$\text{depth}(A_\mathfrak p) + \dim(A/\mathfrak p) > 3$。
\end{enumerate}
\end{enumerate}
则限制函子
$$
\colim_{U_0 \subset U' \subset U\text{ 开}} \textit{F\'Et}_{U'}
\longrightarrow
\textit{F\'Et}_{U_0}
$$
为等价。
\end{lemma}

\begin{proof}
这由引理 \ref{pione-lemma-restriction-equivalence-general} 和《代数化与形式几何》引理
\ref{algebraization-lemma-equivalence-better} 得出。
\end{proof}

\begin{lemma}
\label{pione-lemma-essentially-surjective-general}
在情形 \ref{pione-situation-local-lefschetz} 下，假设
\begin{enumerate}
\item $A$ 关于 $f$-进拓扑完备；
\item $f$ 是非零因子；
\item $H^1_\mathfrak m(A/fA)$ 和 $H^2_\mathfrak m(A/fA)$ 都是有限 $A$-模。
\end{enumerate}
则限制函子
$$
\colim_{U_0 \subset U' \subset U\text{ 开}} \textit{F\'Et}_{U'}
\longrightarrow
\textit{F\'Et}_{U_0}
$$
为等价。
\end{lemma}

\begin{proof}
这由引理 \ref{pione-lemma-restriction-equivalence-general} 和《代数化与形式几何》引理
\ref{algebraization-lemma-equivalence} 得出。
\end{proof}

\begin{remark}
\label{pione-remark-combine}
设 $(A, \mathfrak m)$ 是完备 Noether 局部环，并设 $f \in \mathfrak m$ 非零。
假设 $A_f$ 满足 $(S_2)$，且 $\Spec(A)$ 的每个不可约分支的维数都 $\geq 4$。
那么，引理 \ref{pione-lemma-essentially-surjective-general-better} 告诉我们，范畴
$$
\colim\nolimits_{U' \subset U\text{ open, }U_0 \subset U}
\text{ category of schemes finite \'etale over }U'
$$
等价于在 $U_0$ 上有限 \'etale 的概形范畴。例如，当 $A$ 是维数 $\geq 4$
的正规整环时，这一结论成立！
\end{remark}




\section{穿孔谱的有限 \'etale 覆盖，IV}
\label{pione-section-pi1-punctured-spec-IV}

\noindent
设 $X, X_0, U, U_0$ 如情形 \ref{pione-situation-local-lefschetz}。本节考察限制函子
$$
\textit{F\'Et}_U
\longrightarrow
\textit{F\'Et}_{U_0}
$$
何时本质满。我们将先取用 Section \ref{pione-section-pi1-punctured-spec-III} 的结果，
再用纯性填补其间的缺口。回顾：若对 Noether 局部环 $(A, \mathfrak m)$，
限制函子 $\textit{F\'Et}_X \to \textit{F\'Et}_U$ 本质满，其中
$X = \Spec(A)$ 且 $U = X \setminus \{\mathfrak m\}$，则称该环{\it 满足纯性}。

\begin{lemma}
\label{pione-lemma-equivalence-better}
在情形 \ref{pione-situation-local-lefschetz} 下，假设
\begin{enumerate}
\item $A$ 有对偶化复形，且关于 $f$-进拓扑完备；
\item 下列条件之一成立：
\begin{enumerate}
\item $A_f$ 满足 $(S_2)$，且 $X$ 的每个不包含于 $X_0$ 的不可约分支的维数都
$\geq 4$；或
\item 若 $\mathfrak p \not \in V(f)$ 且
$V(\mathfrak p) \cap V(f) \not = \{\mathfrak m\}$，则
$\text{depth}(A_\mathfrak p) + \dim(A/\mathfrak p) > 3$。
\end{enumerate}
\item 对每个极大理想 $\mathfrak p \subset A_f$，纯性对
$(A_f)_\mathfrak p$ 成立。
\end{enumerate}
则限制函子 $\textit{F\'Et}_U \to \textit{F\'Et}_{U_0}$ 本质满。
\end{lemma}

\begin{proof}
设 $V_0 \to U_0$ 是有限 \'etale 态射。由引理
\ref{pione-lemma-essentially-surjective-general-better}，存在开子集 $U' \subset U$，
它包含 $U_0$，以及有限 \'etale 态射 $V' \to U$，其到 $U_0$ 的基变换同构于
$V_0 \to U_0$。由于 $U' \supset U_0$，可见 $U \setminus U'$ 由与 (3) 中的素理想
$\mathfrak p_1, \ldots, \mathfrak p_n$ 对应的点组成。由假设，可以找到有限
\'etale 态射 $V'_i \to \Spec(A_{\mathfrak p_i})$，它与 $V' \to U'$ 一致
（在 $U' \times_U \Spec(A_{\mathfrak p_i})$ 上）。把《极限》引理
\ref{limits-lemma-glueing-near-closed-point} 应用 $n$ 次，可知 $V' \to U'$ 延拓为
有限 \'etale 态射 $V \to U$。
\end{proof}

\begin{lemma}
\label{pione-lemma-equivalence}
设 $(A, \mathfrak m)$ 是 Noether 局部环。设 $f \in \mathfrak m$。假设
\begin{enumerate}
\item $A$ 关于 $f$-进拓扑完备；
\item $f$ 是非零因子；
\item $H^1_\mathfrak m(A/fA)$ 和 $H^2_\mathfrak m(A/fA)$ 都是有限 $A$-模；
\item 对每个极大理想 $\mathfrak p \subset A_f$，纯性对
$(A_f)_\mathfrak p$ 成立。
\end{enumerate}
则限制函子 $\textit{F\'Et}_U \to \textit{F\'Et}_{U_0}$ 本质满。
\end{lemma}

\begin{proof}
证明与引理 \ref{pione-lemma-equivalence-better} 的证明完全相同，只需使用引理
\ref{pione-lemma-essentially-surjective-general} 代替引理
\ref{pione-lemma-essentially-surjective-general-better}。
\end{proof}



\section{局部情形中的纯性，II}
\label{pione-section-local-purity-II}

\noindent
本节是 Section \ref{pione-section-local-purity} 的续篇。回顾：若对 Noether 局部环
$(A, \mathfrak m)$，限制函子 $\textit{F\'Et}_X \to \textit{F\'Et}_U$
本质满，其中 $X = \Spec(A)$ 且 $U = X \setminus \{\mathfrak m\}$，
则称该环{\it 满足纯性}。

\begin{lemma}
\label{pione-lemma-purity-inherited-by-hypersurface-better}
设 $(A, \mathfrak m)$ 是 Noether 局部环。设 $f \in \mathfrak m$。假设
\begin{enumerate}
\item $A$ 有对偶化复形，且关于 $f$-进拓扑完备；
\item 下列条件之一成立：
\begin{enumerate}
\item $A_f$ 满足 $(S_2)$，且 $X$ 的每个不包含于 $X_0$ 的不可约分支的维数都
$\geq 4$；或
\item 若 $\mathfrak p \not \in V(f)$ 且
$V(\mathfrak p) \cap V(f) \not = \{\mathfrak m\}$，则
$\text{depth}(A_\mathfrak p) + \dim(A/\mathfrak p) > 3$。
\end{enumerate}
\item 对每个极大理想 $\mathfrak p \subset A_f$，纯性对
$(A_f)_\mathfrak p$ 成立；并且
\item 纯性对 $A$ 成立。
\end{enumerate}
则纯性对 $A/fA$ 成立。
\end{lemma}

\begin{proof}
置 $X = \Spec(A)$，并令 $U = X \setminus \{\mathfrak m\}$ 为穿孔谱。类似地，
有 $X_0 = \Spec(A/fA)$ 和 $U_0 = X_0 \setminus \{\mathfrak m\}$。设
$V_0 \to U_0$ 是有限 \'etale 态射。由引理 \ref{pione-lemma-equivalence-better}，
可得有限 \'etale 态射 $V \to U$，其到 $U_0$ 的基变换同构于
$V_0 \to U_0$。由假设 (5)，$V \to U$ 延拓为有限 \'etale 态射
$Y \to X$。于是，$Y$ 到 $X_0$ 的限制就是 $V_0 \to U_0$ 所需的延拓。
\end{proof}

\begin{lemma}
\label{pione-lemma-purity-inherited-by-hypersurface}
设 $(A, \mathfrak m)$ 是 Noether 局部环。设 $f \in \mathfrak m$。假设
\begin{enumerate}
\item $A$ 关于 $f$-进拓扑完备；
\item $f$ 是非零因子；
\item $H^1_\mathfrak m(A/fA)$ 和 $H^2_\mathfrak m(A/fA)$ 都是有限 $A$-模；
\item 对每个极大理想 $\mathfrak p \subset A_f$，纯性对
$(A_f)_\mathfrak p$ 成立；
\item 纯性对 $A$ 成立。
\end{enumerate}
则纯性对 $A/fA$ 成立。
\end{lemma}

\begin{proof}
证明与引理 \ref{pione-lemma-purity-inherited-by-hypersurface-better} 的证明完全相同，
只需使用引理 \ref{pione-lemma-equivalence} 代替引理 \ref{pione-lemma-equivalence-better}。
\end{proof}

\noindent
现在可以利用前面的结果，证明纯性对维数 $\geq 3$ 的完全交成立。回顾：若一个
Noether 局部环的完备化，是某个正则局部环模一条正则序列所生成理想的商，
则称该环为完全交。见《除幂代数》Section \ref{dpa-section-lci} 中的讨论。

\begin{proposition}
\label{pione-proposition-purity-complete-intersection}
设 $(A, \mathfrak m)$ 是 Noether 局部环。若 $A$ 是维数 $\geq 3$ 的完全交，
则纯性对 $A$ 成立；也就是说，穿孔谱的任意有限 \'etale 覆盖都可延拓。
\end{proposition}

\begin{proof}
由引理 \ref{pione-lemma-purity-and-completion}，可设 $A$ 是完备局部环。由假设，可以写成
$A = B/(f_1, \ldots, f_r)$，其中 $B$ 是完备正则局部环，且
$f_1, \ldots, f_r$ 是正则序列。我们对 $r$ 作归纳以完成证明。基本情形为
$r = 0$；这由引理 \ref{pione-lemma-local-purity} 得出，该引理适用于维数 $\geq 2$
的正则环。

\medskip\noindent
假设 $A = B/(f_1, \ldots, f_r)$，且命题对 $r - 1$ 成立。置
$A' = B/(f_1, \ldots, f_{r - 1})$，并把引理
\ref{pione-lemma-purity-inherited-by-hypersurface} 应用于 $f_r \in A'$。这是允许的：
条件 (1) 成立，因为 $f_1, \ldots, f_r$ 是正则序列；
条件 (2) 成立，因为 $B$ 完备，从而 $A'$ 完备；
条件 (3) 成立，因为 $A = A'/f_r A'$ 是维数为 $\dim(A) \geq 3$ 的
Cohen--Macaulay 环，见《对偶化复形》引理 \ref{dualizing-lemma-depth}；
条件 (4) 由归纳假设成立，因为 $\dim((A'_{f_r})_\mathfrak p) \geq 3$，其中
$\mathfrak p$ 是 $A'_{f_r}$ 的一个极大素理想，并且
$(A'_{f_r})_\mathfrak p = B_\mathfrak q/(f_1, \ldots, f_{r - 1})$ 对某个
$\mathfrak q \subset B$ 成立；
条件 (5) 由归纳假设成立。
\end{proof}




\section{局部情形中的纯性，III}
\label{pione-section-local-purity-III}

\noindent
本节继续 Sections \ref{pione-section-local-purity} 和
\ref{pione-section-local-purity-II} 中的讨论。

\begin{lemma}
\label{pione-lemma-fully-faithful-power-series-over-depth2}
设 $(A, \mathfrak m)$ 是深度 $\geq 2$ 的 Noether 局部环。设
$B = A[[x_1, \ldots, x_d]]$，其中 $d \geq 1$。置
$Y = \Spec(B)$ 和 $Y_0 = V(x_1, \ldots, x_d)$。对任意开子概形
$V \subset Y$，若 $V_0 = V \cap Y_0$ 等于
$Y_0 \setminus \{\mathfrak m_B\}$，则限制函子
$$
\textit{F\'Et}_V \longrightarrow \textit{F\'Et}_{V_0}
$$
全忠实。
\end{lemma}

\begin{proof}
置 $I = (x_1, \ldots, x_d)$。置 $X = \Spec(A)$。利用态射 $Y \to X$
把 $Y_0$ 与 $X$ 等同后，$V_0$ 就与穿孔谱 $U$ 等同，即 $A$ 的穿孔谱。
沿这个仿射态射
前推模，得到
\begin{align*}
\lim_n \Gamma(V_0, \mathcal{O}_V/I^n\mathcal{O}_V)
& =
\lim_n \Gamma(V_0, \mathcal{O}_Y/I^n\mathcal{O}_Y) \\
& =
\lim_n \Gamma(U, \mathcal{O}_U[x_1, \ldots, x_d]/(x_1, \ldots, x_d)^n) \\
& =
\lim_n A[x_1, \ldots, x_d]/(x_1, \ldots, x_d)^n \\
& =
B
\end{align*}
事实上，由于 $A$ 的深度 $\geq 2$，有 $\Gamma(U, \mathcal{O}_U) = A$；
见《局部上同调》引理
\ref{local-cohomology-lemma-finiteness-pushforwards-and-H1-local}。
因此，对引理中任意开子集 $V \subset Y$，有
$$
B = \Gamma(Y, \mathcal{O}_Y) \to \Gamma(V, \mathcal{O}_V) \to
\lim_n \Gamma(V_0, \mathcal{O}_Y/I^n\mathcal{O}_Y) = B
$$
这蕴含两个箭头都是同构（略去一个小细节）。由《代数化与形式几何》引理
\ref{algebraization-lemma-completion-fully-faithful}，可知
$\textit{Coh}(\mathcal{O}_V) \to \textit{Coh}(V, I\mathcal{O}_V)$
在有限局部自由对象所成的满子范畴上全忠实。因此，由引理
\ref{pione-lemma-restriction-fully-faithful} 即得结论。
\end{proof}

\begin{lemma}
\label{pione-lemma-purity-power-series-over-depth2}
\begin{slogan}
有限 \'etale 覆盖的 Ramanujam--Samuel 定理
\end{slogan}
设 $(A, \mathfrak m)$ 是深度 $\geq 2$ 的 Noether 局部环。设
$B = A[[x_1, \ldots, x_d]]$，其中 $d \geq 1$。对任意开子集
$V \subset Y = \Spec(B)$，若它包含
\begin{enumerate}
\item 任意素理想 $\mathfrak q \subset B$，只要
$\mathfrak q \cap A \not = \mathfrak m$；
\item 素理想 $\mathfrak m B$，
\end{enumerate}
则函子
$
\textit{F\'Et}_Y
\to
\textit{F\'Et}_V
$
为等价。特别地，纯性对 $B$ 成立。
\end{lemma}

\begin{proof}
素理想 $\mathfrak q \subset B$ 若不属于 $V$，则位于 $\mathfrak m$ 上方。在这种情形下，
$A \to B_\mathfrak q$ 是平坦局部同态，因而
$\text{depth}(B_\mathfrak q) \geq 2$（《代数》引理
\ref{algebra-lemma-apply-grothendieck}）。因此，结合引理
\ref{pione-lemma-quasi-compact-dense-open-connected-at-infinity-Noetherian} 与《局部上同调》引理
\ref{local-cohomology-lemma-depth-2-connected-punctured-spectrum}，该函子全忠实。

\medskip\noindent
设 $W \to V$ 是有限 \'etale 态射。设 $B \to C$ 是唯一的有限环同态，使得
$\Spec(C) \to Y$ 是引理 \ref{pione-lemma-extend-S2} 中构造的、延拓 $W \to V$
的有限态射。注意，$C = \Gamma(W, \mathcal{O}_W)$。

\medskip\noindent
置 $Y_0 = V(x_1, \ldots, x_d)$ 和 $V_0 = V \cap Y_0$。置 $X = \Spec(A)$。
利用态射 $Y \to X$ 把 $Y_0$ 与 $X$ 等同后，$V_0$ 就与穿孔谱 $U$ 等同，
即 $A$ 的穿孔谱。因此，可把 $W_0 = W \times_Y Y_0$ 看作 $U$ 上的有限 \'etale 概形。于是
$$
W_0 \times_U (U \times_X Y)
\quad\text{且}\quad
W \times_V (U \times_X Y)
$$
是 $U \times_X Y$ 上的有限 \'etale 概形，它们在 $V_0$ 上的限制同构。把引理
\ref{pione-lemma-fully-faithful-power-series-over-depth2} 应用于开子集 $U \times_X Y$，
得到同构
$$
W_0 \times_U (U \times_X Y) \longrightarrow W \times_V (U \times_X Y)
$$
（在 $U \times_X Y$ 上）。

\medskip\noindent
注意，$C_0 = \Gamma(W_0, \mathcal{O}_{W_0})$ 是有限 $A$-代数；这是把引理
\ref{pione-lemma-extend-S2} 应用于 $W_0 \to U \subset X$ 所得（与上面对
$B \to C$ 所做的完全相同）。由于引理
\ref{pione-lemma-extend-S2} 中的构造与平坦基变换及开集的改变相容，上面的同构诱导同构
$$
\Psi : C \longrightarrow C_0 \otimes_A B
$$
（作为有限 $B$-代数）。然而，我们知道，$\Spec(C) \to Y$ 在位于 $Y$
中至少一个点上方的所有点处是 \'etale 的，而该点位于 $\mathfrak m \in X$
上方。由于 $\Psi$ 是同构，可知 $\Spec(C_0) \to X$ 在 $\mathfrak m$
上方是 \'etale 的（略去一个小细节）。当然，这意味着 $A \to C_0$ 有限
\'etale，因而 $B \to C$ 有限 \'etale。
\end{proof}

\begin{lemma}
\label{pione-lemma-purity-smooth-over-depth2}
设 $f : X \to S$ 是概形的态射。设 $U \subset X$ 是开子概形。假设
\begin{enumerate}
\item $f$ 光滑；
\item $S$ 是 Noether 概形；
\item 对 $s \in S$，若 $\text{depth}(\mathcal{O}_{S, s}) \leq 1$，则有
$X_s = U_s$；
\item $U_s \subset X_s$ 对所有 $s \in S$ 都稠密。
\end{enumerate}
则 $\textit{F\'Et}_X \to \textit{F\'Et}_U$ 为等价。
\end{lemma}

\begin{proof}
结合引理 \ref{pione-lemma-quasi-compact-dense-open-connected-at-infinity-Noetherian} 与
《局部上同调》引理 \ref{local-cohomology-lemma-depth-2-connected-punctured-spectrum}，
该函子全忠实（此外应用《代数》引理
\ref{algebra-lemma-apply-grothendieck} 来检验深度条件）。

\medskip\noindent
设 $\pi : V \to U$ 是有限 \'etale 态射。设 $Y \to X$ 是引理
\ref{pione-lemma-extend-S2} 中构造的有限态射。必须证明 $Y \to X$ 有限 \'etale。
为了证明这一点对所有 $x \in X$ 成立，固定其像 $s \in S$；可以用 Noether
概形的平坦态射 $S' \to S$ 作基变换，只要 $s$ 位于其像中。这由引理
\ref{pione-lemma-extend-S2} 中的构造与平坦基变换相容得出。

\medskip\noindent
扩大 $U$ 后，可设 $U \subset X$ 是使 $Y \to X$ 有限 \'etale 的最大开集。
令 $Z \subset X$ 为 $U$ 的补集。为导出矛盾，假设 $Z \not = \emptyset$。
选取一点 $s \in S$，它位于 $Z \to S$ 的像中，并且 $s$ 的任何严格一般化都不在该像中。
对 $\Spec(\mathcal{O}_{S, s})$ 作基变换后，可知 $S = \Spec(A)$，其中
$(A, \mathfrak m, \kappa)$ 是深度 $\geq 2$ 的 Noether 局部环，且 $Z$
包含于闭纤维 $X_s$ 中并在 $X_s$ 中无处稠密。选取闭点 $z \in Z$。于是
$\kappa(z)/\kappa$ 有限（由 Hilbert 零点定理；见《代数》定理
\ref{algebra-theorem-nullstellensatz}）。选取局部概形的有限平坦态射
$(S', s') \to (S, s)$ 来实现剩余域扩张 $\kappa(z)/\kappa$；见《代数》引理
\ref{algebra-lemma-finite-free-given-residue-field-extension}。由 $S' \to S$
作基变换后，约化到 $\kappa(z) = \kappa$ 的情形。

\medskip\noindent
由《态射补篇》引理 \ref{more-morphisms-lemma-slice-smooth}，存在局部闭子概形
$S' \subset X$，它经过 $z$，使得 $S' \to S$ 在 $z$ 处为 \'etale。由
$S' \to S$ 作基变换后，可设存在截面 $\sigma : S \to X$，满足
$\sigma(s) = z$。选取仿射邻域 $\Spec(B) \subset X$，它包含 $s$。于是
$A \to B$ 是光滑环同态，并有截面 $\sigma : B \to A$。记
$I = \Ker(\sigma)$，并以 $B^\wedge$ 表示关于 $I$ 的 $B$ 之完备化。于是有
$B^\wedge \cong A[[x_1, \ldots, x_d]]$，其中 $d \geq 0$；见《代数》引理
\ref{algebra-lemma-section-smooth}。注意 $d > 0$，否则 $X \to S$ 在 $z$
处为 \'etale，这将蕴含 $z$ 是 $X_s$ 的一般点，继而由假设 (4) 有
$z \in U$。类似地，若 $d > 0$，则 $\mathfrak m B^\wedge$ 映入 $U$，这里所用态射为
$\Spec(B^\wedge) \to X$。只需证明 $Y \to X$ 在基变换到
$\Spec(B^\wedge)$ 后有限 \'etale。由于 $B \to B^\wedge$ 平坦（《代数》引理
\ref{algebra-lemma-completion-flat}），这由引理
\ref{pione-lemma-purity-power-series-over-depth2} 和 $Y \to X$ 构造中的唯一性得出。
\end{proof}

\begin{proposition}
\label{pione-proposition-purity-smooth-over-depth2}
设 $A \to B$ 是 Noether 局部环之间的局部同态。假设 $A$ 的深度 $\geq 2$，
$A \to B$ 关于 $\mathfrak m_B$-进拓扑形式光滑，且 $\dim(B) > \dim(A)$。
对任意开子集 $V \subset Y = \Spec(B)$，若它包含
\begin{enumerate}
\item 任意素理想 $\mathfrak q \subset B$，只要
$\mathfrak q \cap A \not = \mathfrak m_A$；
\item 素理想 $\mathfrak m_A B$，
\end{enumerate}
则函子 $\textit{F\'Et}_Y \to \textit{F\'Et}_V$ 为等价。特别地，纯性对 $B$ 成立。
\end{proposition}

\begin{proof}
素理想 $\mathfrak q \subset B$ 若不属于 $V$，则位于 $\mathfrak m_A$ 上方。
在这种情形下，$A \to B_\mathfrak q$ 是平坦局部同态，因而
$\text{depth}(B_\mathfrak q) \geq 2$（《代数》引理
\ref{algebra-lemma-apply-grothendieck}）。因此，结合引理
\ref{pione-lemma-quasi-compact-dense-open-connected-at-infinity-Noetherian} 与《局部上同调》引理
\ref{local-cohomology-lemma-depth-2-connected-punctured-spectrum}，该函子全忠实。

\medskip\noindent
以 $A^\wedge$ 和 $B^\wedge$ 表示 $A$ 和 $B$ 关于各自极大理想的完备化。注意，
命题的假设对 $A^\wedge \to B^\wedge$ 成立；见《代数补篇》引理
\ref{more-algebra-lemma-completion-dimension}、
\ref{more-algebra-lemma-completion-depth} 和
\ref{more-algebra-lemma-formally-smooth-completion}。由引理
\ref{pione-lemma-extend-S2} 中构造的唯一性及其与平坦基变换的相容性，只需对
$A^\wedge \to B^\wedge$ 和 $V$ 的逆像证明本质满性（略去细节；当 $V$
为穿孔谱时，可与引理 \ref{pione-lemma-purity-and-completion} 比较）。由《代数补篇》命题
\ref{more-algebra-proposition-fs-regular}，这意味着可设 $A \to B$ 正则。

\medskip\noindent
设 $W \to V$ 是有限 \'etale 态射。由 Popescu 定理（《环同态的光滑化》定理
\ref{smoothing-theorem-popescu}），可把 $B = \colim B_i$ 写成光滑 $A$-代数的
滤过余极限。可以选取某个 $i$ 和开子集 $V_i \subset \Spec(B_i)$，使其逆像为
$V$（《极限》引理 \ref{limits-lemma-descend-opens}）。增大 $i$ 后，可设存在有限
\'etale 态射 $W_i \to V_i$，其到 $V$ 的基变换为 $W \to V$；见《极限》引理
\ref{limits-lemma-descend-finite-presentation}、
\ref{limits-lemma-descend-finite-finite-presentation} 和
\ref{limits-lemma-descend-etale}。可以假设 $V_i$ 的补集包含于
$\Spec(B_i) \to \Spec(A)$ 的闭纤维，因为这一点对 $V$ 成立（或者如此选取
$V_i$，或者使用上面的引理证明它对充分大的 $i$ 成立）。设 $\eta$ 是
$\Spec(B) \to \Spec(A)$ 之闭纤维的一般点。由于 $\eta \in V$，$\eta$ 的像位于
$V_i$。因此，把 $V_i$ 替换为 $\Spec(B)$ 之闭点的像的一个仿射开邻域后，
可设 $\Spec(B_i) \to \Spec(A)$ 的闭纤维不可约，且其一般点包含于 $V_i$
（略去细节；使用域上光滑的概形是不可约概形的不交并这一事实）。此时可以应用引理
\ref{pione-lemma-purity-smooth-over-depth2}，看出 $W_i \to V_i$ 延拓为有限 \'etale 态射
$\Spec(C_i) \to \Spec(B_i)$。把它拉回到 $\Spec(B)$，便可得出 $W$ 位于函子
$\textit{F\'Et}_Y \to \textit{F\'Et}_V$ 的本质像中，正如所需。
\end{proof}





\section{基本群的 Lefschetz 定理}
\label{pione-section-global-lefschetz}

\noindent
当然，我们已经在局部情形中证明了许多这种类型的结果。本节讨论射影情形。

\begin{proposition}
\label{pione-proposition-lefschetz-fully-faithful}
设 $k$ 是域。设 $X$ 是 $k$ 上的固有概形。设 $\mathcal{L}$ 是丰沛可逆
$\mathcal{O}_X$-模。设 $s \in \Gamma(X, \mathcal{L})$。设 $Y = Z(s)$ 是
$s$ 的零点概形。假设对所有 $x \in X \setminus Y$，都有
$$
\text{depth}(\mathcal{O}_{X, x}) + \dim(\overline{\{x\}}) > 1
$$
则限制函子 $\textit{F\'Et}_X \to \textit{F\'Et}_Y$ 全忠实。事实上，
对任意开子概形 $V \subset X$，若它包含 $Y$，则限制函子
$\textit{F\'Et}_V \to \textit{F\'Et}_Y$ 全忠实。
\end{proposition}

\begin{proof}
第一个断言是引理 \ref{pione-lemma-restriction-fully-faithful-special} 和《代数化与形式几何》命题
\ref{algebraization-proposition-lefschetz} 的形式推论。第二个断言由引理
\ref{pione-lemma-restriction-fully-faithful-special} 和《代数化与形式几何》引理
\ref{algebraization-lemma-lefschetz-addendum} 得出。
\end{proof}

\begin{proposition}
\label{pione-proposition-lefschetz-equivalence-general}
设 $k$ 是域。设 $X$ 是 $k$ 上的固有概形。设 $\mathcal{L}$ 是丰沛可逆
$\mathcal{O}_X$-模。设 $s \in \Gamma(X, \mathcal{L})$。设 $Y = Z(s)$ 是
$s$ 的零点概形。令 $\mathcal{V}$ 为 $X$ 中包含 $Y$ 的开子概形所成集合，
按逆包含关系排序。假设对所有 $x \in X \setminus Y$，都有
$$
\text{depth}(\mathcal{O}_{X, x}) + \dim(\overline{\{x\}}) > 2
$$
则限制函子
$$
\colim_\mathcal{V} \textit{F\'Et}_V \to \textit{F\'Et}_Y
$$
为等价。
\end{proposition}

\begin{proof}
这是引理 \ref{pione-lemma-restriction-equivalence-general} 和《代数化与形式几何》命题
\ref{algebraization-proposition-lefschetz-equivalence} 的形式推论。
\end{proof}

\begin{proposition}
\label{pione-proposition-lefschetz-equivalence}
设 $k$ 是域。设 $X$ 是 $k$ 上的固有概形。设 $\mathcal{L}$ 是丰沛可逆
$\mathcal{O}_X$-模。设 $s \in \Gamma(X, \mathcal{L})$。设 $Y = Z(s)$ 是
$s$ 的零点概形。假设对所有 $x \in X \setminus Y$，都有
$$
\text{depth}(\mathcal{O}_{X, x}) + \dim(\overline{\{x\}}) > 2
$$
并且对闭点 $x \in X \setminus Y$，纯性对 $\mathcal{O}_{X, x}$ 成立。则限制函子
$\textit{F\'Et}_X \to \textit{F\'Et}_Y$ 为等价。若 $X$（等价地，$Y$）连通，则
$$
\pi_1(Y, \overline{y}) \to \pi_1(X, \overline{y})
$$
该态射对任意几何点 $\overline{y}$ 均为同构；这里该点取自 $Y$。
\end{proposition}

\begin{proof}
由命题 \ref{pione-proposition-lefschetz-fully-faithful}，全忠实性成立。由命题
\ref{pione-proposition-lefschetz-equivalence-general}，$\textit{F\'Et}_Y$ 的任意对象都同构于
纤维积 $U \times_V Y$，其中 $U \to V$ 是某个有限 \'etale 态射，且
$V \subset X$ 是包含 $Y$ 的开子概形。补集 $T = X \setminus V$
是\footnote{事实上，$T$ 在 $k$ 上固有（因为它在 $X$ 中闭），并且仿射
（因为它在仿射概形 $X \setminus Y$ 中闭；见《态射》引理
\ref{morphisms-lemma-proper-ample-delete-affine}），因此在 $k$ 上有限
（《态射》引理 \ref{morphisms-lemma-finite-proper}）。所以 $T$ 是闭点的有限集合。}
$X \setminus Y$ 中闭点的有限集合。写成 $T = \{x_1, \ldots, x_n\}$。
由假设，可以找到有限 \'etale 态射
$V'_i \to \Spec(\mathcal{O}_{X, x_i})$，它与 $U \to V$ 一致（在
$V \times_X \Spec(\mathcal{O}_{X, x_i})$ 上）。把《极限》引理
\ref{limits-lemma-glueing-near-closed-point} 应用 $n$ 次，可见 $U \to V$
按要求延拓为有限 \'etale 态射 $U' \to X$。最后一个断言见引理
\ref{pione-lemma-what-equivalence-gives}。
\end{proof}










\section{分歧轨迹的纯性}
\label{pione-section-purity-ramification}

\noindent
本节讨论泛有限态射情形下与分支轨迹纯性相对应的结果。该结果似乎归功于
Gabber。其一个特殊情形是 van der Waerden 的纯性定理：从正规簇到光滑簇的
双有理态射不为同构的轨迹具有纯性。

\begin{lemma}
\label{pione-lemma-characterize-rational-singularity}
设 $A$ 是维数为 $2$ 的 Noether 正规局部整环。假设 $A$ 是 Nagata 环，
具有对偶化模 $\omega_A$，并且存在奇点解消
$f : X \to \Spec(A)$。按《曲面奇点的消解》注记
\ref{resolve-remark-dualizing-setup} 取 $\omega_X$。若
$\omega_X \cong \mathcal{O}_X(E)$，其中 $E \subset X$ 是某个支撑在例外纤维上的
有效 Cartier 除子，则 $A$ 定义一个有理奇点。
若 $f$ 是极小解消，则 $E = 0$。
\end{lemma}

\begin{proof}
存在迹映射 $Rf_*\omega_X \to \omega_A$，见《概形的对偶性》节
\ref{duality-section-trace}。由 Grauert--Riemenschneider 定理
（《曲面奇点的消解》命题
\ref{resolve-proposition-Grauert-Riemenschneider}），有
$R^1f_*\omega_X = 0$。因此迹映射就是映射
$f_*\omega_X \to \omega_A$。于是可以考虑
$$
\mathcal{O}_{\Spec(A)} = f_*\mathcal{O}_X \to f_*\omega_X \to \omega_A
$$
其中第一个映射来自引理陈述中假设存在的映射
$\mathcal{O}_X \to \mathcal{O}_X(E) = \omega_X$。由《除子》引理
\ref{divisors-lemma-check-isomorphism-via-depth-and-ass}，该复合映射是同构：
它在 $A$ 的穿孔谱上是同构（由引理中的假设以及 $f$ 在穿孔谱上为同构）；
而 $A$ 与 $\omega_A$ 都是 $A$-模，且深度为 $2$（由《代数》引理
\ref{algebra-lemma-criterion-normal} 与《对偶化复形》引理
\ref{dualizing-lemma-depth-dualizing-module}）。因此
$f_*\omega_X \to \omega_A$ 满射，进而为同构。于是
$Rf_*\omega_X = \omega_A$；由对偶性，这蕴含
$Rf_*\mathcal{O}_X = \mathcal{O}_{\Spec(A)}$。从而
$H^1(X, \mathcal{O}_X) = 0$，这说明 $A$ 定义一个有理奇点（见《曲面奇点的消解》
节 \ref{resolve-section-bounded} 中的讨论，特别是引理
\ref{resolve-lemma-regular-rational} 与
\ref{resolve-lemma-exact-sequence}）。若 $f$ 极小，则 $E = 0$：反复应用
《曲面奇点的消解》引理
\ref{resolve-lemma-dualizing-blow-up-rational} 可知映射
$f^*\omega_A \to \omega_X$ 满射，而如上所见 $\omega_A \cong A$。
\end{proof}

\begin{lemma}
\label{pione-lemma-key-purity-ramification}
设 $f : X \to \Spec(A)$ 是有限型态射。设 $x \in X$ 是一点。假设
\begin{enumerate}
\item $A$ 是优良正则局部环；
\item $\mathcal{O}_{X, x}$ 正规且维数为 $2$；
\item $f$ 在 $\overline{\{x\}}$ 之外为 \'etale。
\end{enumerate}
则 $f$ 在 $x$ 处为 \'etale。
\end{lemma}

\begin{proof}
首先把 $X$ 换成 $x$ 的一个仿射开邻域。注意，
$\mathcal{O}_{X, x}$ 是优良局部环（《代数补篇》引理
\ref{more-algebra-lemma-finite-type-over-excellent}）。因此可以选取一个极小奇点解消
$W \to \Spec(\mathcal{O}_{X, x})$，见《曲面奇点的消解》定理
\ref{resolve-theorem-resolve}。必要时把 $X$ 换成 $x$ 的仿射开邻域，可以找到固有态射
$b : X' \to X$，使得
$X' \times_X \Spec(\mathcal{O}_{X, x}) = W$，见《极限》引理
\ref{limits-lemma-glueing-near-closed-point}。进一步缩小 $X$ 后，可设 $X'$ 正则。
事实上，已知 $W$ 正则且 $X'$ 优良，而优良环的谱的正则轨迹是开的。由于
$W \to \Spec(\mathcal{O}_{X, x})$ 是射影态射（它是一列正规化爆破），
缩小 $X$ 后可设 $b$ 射影（略去细节）。令
$U = X \setminus \overline{\{x\}}$。由于
$W \to \Spec(\mathcal{O}_{X, x})$ 在穿孔谱上为同构，可设
$b : X' \to X$ 在 $U$ 上为同构。因此也把 $U$ 看作 $X'$ 的开子概形。
置 $f' = f \circ b : X' \to \Spec(A)$。

\medskip\noindent
由于 $A$ 正则，可见 $\mathcal{O}_Y$ 是 $Y$ 的对偶化复形。因此
$f^!\mathcal{O}_Y$ 是 $X$ 上的对偶化复形（《概形的对偶性》引理
\ref{duality-lemma-shriek-dualizing}）。由《概形的对偶性》引理
\ref{duality-lemma-dualizing-module-CM-scheme}，$X$ 的 Cohen--Macaulay
轨迹是开的（也可以利用优良性证明这一点）。由于
$\mathcal{O}_{X, x}$ 是 Cohen--Macaulay 的，缩小 $X$ 后可设
$X$ 是 Cohen--Macaulay 的。注意，\'etale 态射是局部完全交。因此可以把
《概形的对偶性》引理
\ref{duality-lemma-fundamental-class-almost-lci} 用于 $r = 0$，得到映射
$$
\mathcal{O}_X \longrightarrow \omega_{X/Y} = H^0(f^!\mathcal{O}_Y)
$$
它在 $X \setminus \overline{\{x\}}$ 上为同构。由《概形的对偶性》引理
\ref{duality-lemma-shriek-over-CM}，$\omega_{X/Y}$ 满足 $(S_2)$，故由
《除子》引理 \ref{divisors-lemma-check-isomorphism-via-depth-and-ass}
可知该映射为同构。这已经说明 $X$ 是 Gorenstein 的，特别地，
$\mathcal{O}_{X, x}$ 是 Gorenstein 环。

\medskip\noindent
置 $\omega_{X'/Y} = H^0((f')^!\mathcal{O}_Y)$。与上面完全相同的论证表明
$(f')^!\mathcal{O}_Y = \omega_{X'/Y}[0]$ 是 $X'$ 的对偶化复形。由于
$X'$ 正则，态射 $X' \to Y$ 是局部完全交态射，见《态射补篇》引理
\ref{more-morphisms-lemma-morphism-regular-schemes-is-lci}。由《概形的对偶性》引理
\ref{duality-lemma-fundamental-class-lci}，存在映射
$$
\mathcal{O}_{X'} \longrightarrow \omega_{X'/Y}
$$
它在 $U$ 上为同构。于是
$\omega_{X'/Y} = \mathcal{O}_{X'}(E)$，其中 $E \subset X'$ 是某个与
$U$ 不交的有效 Cartier 除子。

\medskip\noindent
由于 $\omega_{X/Y} = \mathcal{O}_Y$，可见
$\omega_{X'/Y} = b^! f^!\mathcal{O}_Y = b^!\mathcal{O}_X$。回到
$W \to \Spec(\mathcal{O}_{X, x})$，可得
$\omega_W = \mathcal{O}_W(E|_W)$。由引理
\ref{pione-lemma-characterize-rational-singularity}，有 $E|_W = 0$。
于是由上面已经用过的《概形的对偶性》引理
\ref{duality-lemma-fundamental-class-lci}，$f' : X' \to Y$ 为 \'etale。
这立即完成证明，因为 $f'$ 为 \'etale 迫使 $b$ 是同构。
\end{proof}

\begin{lemma}[分歧轨迹的纯性]
\label{pione-lemma-purity-ramification}
\begin{reference}
复空间情形下的这一结果见 \cite{Fischer} 第 170 页。一般情形是
\cite[Theorem 2.4]{Zong} 中归于 Gabber 的结果。
\end{reference}
设 $f : X \to Y$ 是局部 Noether 概形之间的态射。设 $x \in X$，并置
$y = f(x)$。假设
\begin{enumerate}
\item $\mathcal{O}_{X, x}$ 正规且维数 $\geq 1$；
\item $\mathcal{O}_{Y, y}$ 正则；
\item $f$ 是局部有限型的；并且
\item 对特化 $x' \leadsto x$，若
$\dim(\mathcal{O}_{X, x'}) = 1$，则态射 $f$ 在 $x'$ 处为 \'etale。
\end{enumerate}
则 $f$ 在 $x$ 处为 \'etale。
\end{lemma}

\begin{proof}
对 $d = \dim(\mathcal{O}_{X, x})$ 作归纳来证明该引理。

\medskip\noindent
先看平凡情形 $d = 1$；此时由假设，态射 $f$ 在 $x$ 处为 \'etale。
以下假设 $d \geq 2$。

\medskip\noindent
沿 $\Spec(\mathcal{O}_{Y, y}) \to Y$ 作基变换不会影响引理的结论，见
《态射》引理 \ref{morphisms-lemma-set-points-where-fibres-etale}。因此可设
$Y = \Spec(A)$，其中 $A$ 是正则局部环，而 $y$ 所对应的是极大理想
$\mathfrak m$（在 $A$ 中）。

\medskip\noindent
设 $x' \leadsto x$ 是满足 $x' \not = x$ 的特化。
$\mathcal{O}_{X, x'}$ 正规，因为它是 $\mathcal{O}_{X, x}$ 的局部化。若 $x'$ 不是
$X$ 的一般点，则 $1 \leq \dim(\mathcal{O}_{X, x'}) < d$，由归纳假设可知
$f$ 在 $x'$ 处为 \'etale。因此可设 $f$ 在所有特化到 $x$ 的点处均为
\'etale。由于 $f$ 为 \'etale 的点所成集合在 $X$ 中是开的（依定义），
把 $X$ 换成 $x$ 的一个开邻域后，可设 $f$ 在
$\overline{\{x\}}$ 之外为 \'etale。特别地，可见 $f$ 仅可能在位于闭点
$y \in Y = \Spec(A)$ 上方的点处不是 \'etale。

\medskip\noindent
令 $X' = X \times_{\Spec(A)} \Spec(A^\wedge)$。设 $x' \in X'$ 是位于
$x$ 上方的唯一一点。由上可见，$X'$ 在闭纤维之外于
$\Spec(A^\wedge)$ 上为 \'etale，因而 $X'$ 在闭纤维之外正规。由于
$X$ 正规，由《曲面奇点的消解》引理
\ref{resolve-lemma-normalization-completion} 可知 $X'$ 正规。于是，只要能证明
$X' \to \Spec(A^\wedge)$ 在 $x'$ 处为 \'etale，就可推出 $f$ 在 $x$
处为 \'etale（由前述《态射》引理
\ref{morphisms-lemma-set-points-where-fibres-etale}）。故可以并且确实假设
$A$ 是正则完备局部环。

\medskip\noindent
此时 $d = 2$ 的情形由引理 \ref{pione-lemma-key-purity-ramification} 得出。

\medskip\noindent
假设 $d > 2$。取 $t \in \mathfrak m$，满足
$t \not \in \mathfrak m^2$。置 $Y_0 = \Spec(A/tA)$ 以及
$X_0 = X \times_Y Y_0$。则 $X_0 \to Y_0$ 在闭点上方的纤维之外为
\'etale。由于 $d > 2$，有
$\dim(\mathcal{O}_{X_0, x}) = d - 1$ 且该数 $\geq 2$。正规化
$X_0' \to X_0$ 是满射且有限（因为我们在完备局部环上工作，而这种环是
Nagata 环）。设 $x' \in X_0'$ 是映到 $x$ 的一点。由归纳假设，态射
$X'_0 \to Y$ 在 $x'$ 处为 \'etale。由包含关系
$\kappa(y) \subset \kappa(x) \subset \kappa(x')$
可知 $\kappa(x)$ 是 $\kappa(y)$ 的有限扩张。因此 $x$ 是
$X \to Y$ 在 $y$ 上方纤维的闭点。但 $x$ 也是该纤维的一般点，所以
$f$ 在 $x$ 处拟有限，从而归约到分支轨迹纯性的情形，见引理
\ref{pione-lemma-purity}。
\end{proof}




\section{分歧轨迹补集的仿射性}
\label{pione-section-stronger-purity}

\noindent
设 $f : X \to Y$ 是 Noether 概形之间的有限型态射，其中 $X$ 正规而
$Y$ 正则。设 $V \subset X$ 是 $f$ 为 \'etale 的最大开子概形。
\cite[Chapter IV, Section 21.12]{EGA} 中的讨论表明，
$V \to X$ 可能是仿射态射。注意，若 $V \to X$ 仿射，则可推出分歧轨迹的纯性
（引理 \ref{pione-lemma-purity-ramification}）；这里使用了《除子》引理
\ref{divisors-lemma-complement-affine-open-immersion}。因此，$V \to X$ 的仿射性是
分歧轨迹纯性的一种“强”形式。本节证明 $V \to X$ 仿射，条件是
$X$ 与 $Y$ 等特征且优良，见定理 \ref{pione-theorem-global}。有理由猜测，当 $X$ 与 $Y$
具有混合特征（但仍优良）时，该结果仍然成立。

\begin{lemma}
\label{pione-lemma-structure-cohomology}
设 $(A, \mathfrak m)$ 是包含一个域的正则局部环。设
$f : V \to \Spec(A)$ 为 \'etale 且准紧。假设
$\mathfrak m \not \in f(V)$，并且
$g : V \to \Spec(A) \setminus \{\mathfrak m\}$ 仿射。则对
$H^i(V, \mathcal{O}_V)$（$i > 0$）同构于 $A$ 的剩余域的单射包的若干副本的直和。
\end{lemma}

\begin{proof}
记 $U = \Spec(A) \setminus \{\mathfrak m\}$ 为穿孔谱。于是
$g : V \to U$ 仿射。由《概形的上同调》引理
\ref{coherent-lemma-relative-affine-cohomology}，有
$H^i(V, \mathcal{O}_V) = H^i(U, g_*\mathcal{O}_V)$。由《概形》引理
\ref{schemes-lemma-push-forward-quasi-coherent}，
$\mathcal{O}_U$-模 $g_*\mathcal{O}_V$ 是拟凝聚的。对任意拟凝聚
$\mathcal{O}_U$-模 $\mathcal{F}$，上同调 $H^i(U, \mathcal{F})$（$i > 0$）是
$\mathfrak m$-幂挠的；例如见《局部上同调》引理
\ref{local-cohomology-lemma-local-cohomology}。特别地，$A$-模
$H^i(V, \mathcal{O}_V)$（$i > 0$）是 $\mathfrak m$-幂挠的。

对任意平坦环映射 $A \to A'$，平坦基变换给出
$H^i(V, \mathcal{O}_V) \otimes_A A' = H^i(V', \mathcal{O}_{V'})$，
其中 $V' = V \times_{\Spec(A)} \Spec(A')$；见《概形的上同调》引理
\ref{coherent-lemma-flat-base-change-cohomology}。取 $A'$ 为 $A$ 的完备化
（由《代数补篇》节 \ref{more-algebra-section-permanence-completion}，它是平坦的），
则
$$
H^i(V, \mathcal{O}_V) = H^i(V, \mathcal{O}_V) \otimes_A A' =
H^i(V', \mathcal{O}_{V'}),\quad\text{对 } i > 0
$$
第一个等号由刚证明的挠性和《代数补篇》引理
\ref{more-algebra-lemma-neighbourhood-equivalence} 得出。此外，剩余域
$k$ 的单射包对 $A$ 与 $A'$ 相同，见《对偶化复形》引理
\ref{dualizing-lemma-compare}。这样便归约到
$A = k[[x_1, \ldots, x_d]]$ 的情形；见《代数》节
\ref{algebra-section-cohen-structure-theorem}。

\medskip\noindent
假设 $k$ 的特征为 $p > 0$。映射 $F : A \to A$，
$a \mapsto a^p$ 是平坦的（《局部上同调》引理
\ref{local-cohomology-lemma-frobenius-flat-regular}）；而由《\'Etale 态射》引理
\ref{etale-lemma-relative-frobenius-etale}，有
$V \times_{\Spec(A), \Spec(F)} \Spec(A) \cong V$，这是 $\Spec(A)$ 上概形的同构。
因此上面的结论给出
$H^i(V, \mathcal{O}_V) \otimes_{A, F} A \cong H^i(V, \mathcal{O}_V)$。
于是所求结论由《局部上同调》引理
\ref{local-cohomology-lemma-structure-torsion-Frobenius-regular} 得出。

\medskip\noindent
假设 $k$ 的特征为 $0$。由《局部上同调》引理
\ref{local-cohomology-lemma-etale-derivation}，存在加性算子
$D_j$，$j = 1, \ldots, d$，作用在 $H^i(V, \mathcal{O}_V)$ 上，
它们满足关于 $\partial_j = \partial/\partial x_j$ 的 Leibniz 法则。
于是所求结论由《局部上同调》引理
\ref{local-cohomology-lemma-structure-torsion-D-module-regular} 得出。
\end{proof}

\begin{lemma}
\label{pione-lemma-conclude}
在引理 \ref{pione-lemma-structure-cohomology} 的情形下，假设
$H^i(V, \mathcal{O}_V) = 0$ 对 $i \geq \dim(A) - 1$ 成立。则 $V$ 仿射。
\end{lemma}

\begin{proof}
令 $k = A/\mathfrak m$。由于
$V \times_{\Spec(A)} \Spec(k) = \emptyset$，由上同调与基变换可得
$$
R\Gamma(V, \mathcal{O}_V) \otimes_A^\mathbf{L} k = 0
$$
见《概形的导出范畴》引理 \ref{perfect-lemma-compare-base-change}。因此存在谱序列
（《代数补篇》例 \ref{more-algebra-example-tor}）
$$
E_2^{p, q} = \text{Tor}_{-p}(k, H^q(V, \mathcal{O}_V)),\quad
d_2^{p, q} : E_2^{p, q} \to E_2^{p + 2, q - 1}
$$
以及 $d_r^{p, q} : E_r^{p, q} \to E_r^{p + r, q - r + 1}$，
并且它收敛到零。由引理 \ref{pione-lemma-structure-cohomology}、《对偶化复形》引理
\ref{dualizing-lemma-tor-injective-hull} 以及假设
$H^i(V, \mathcal{O}_V) = 0$（$i \geq \dim(A) - 1$），可知没有非零微分进入或离开
$(p, q) = (0, 0)$ 项。因此
$H^0(V, \mathcal{O}_V) \otimes_A k = 0$。这意味着，若
$\mathfrak m = (x_1, \ldots, x_d)$，则有开覆盖
$V = \bigcup V \times_{\Spec(A)} \Spec(A_{x_i})$，它由仿射开子概形
$V \times_{\Spec(A)} \Spec(A_{x_i})$ 组成（因为 $V$ 在 $A$ 的穿孔谱上仿射），
并且 $x_1, \ldots, x_d$ 在 $\Gamma(V, \mathcal{O}_V)$ 中生成单位理想。
由《性质》引理 \ref{properties-lemma-characterize-affine}，这蕴含 $V$ 仿射。
\end{proof}

\begin{theorem}
\label{pione-theorem-global}
设 $Y$ 是域上的优良正则概形。设 $f : X \to Y$ 是概形的有限型态射，
且 $X$ 正规。设 $V \subset X$ 是 $f$ 为 \'etale 的最大开子概形。则含入态射
$V \to X$ 仿射。
\end{theorem}

\begin{proof}
设 $x \in X$，其像为 $y \in Y$。只须证明 $V \cap W$ 仿射，其中 $W$ 是
$x$ 的某个仿射开邻域。由于 $\Spec(\mathcal{O}_{X, x})$ 是这些概形 $W$ 的极限，
这当且仅当
$$
V_x = V \times_X \Spec(\mathcal{O}_{X, x})
$$
仿射（《极限》引理 \ref{limits-lemma-limit-affine}）。因此，只要该定理对态射
$X \times_Y \Spec(\mathcal{O}_{Y, y}) \to \Spec(\mathcal{O}_{Y, y})$ 成立，
原定理便成立。特别地，可以设 $Y$ 正则且维数有限，从而可对维数
$d = \dim(Y)$ 作归纳。再次结合相同论证，可以设 $Y$ 是以 $y$ 为闭点的局部概形，
并且 $V \cap (X \setminus f^{-1}(\{y\}) \to X \setminus f^{-1}(\{y\})$
仿射。

\medskip\noindent
设 $x \in X$ 是位于 $y$ 上方的一点。若 $x \in V$，则无须证明。注意，
$f^{-1}(\{y\}) \cap V$ 是由有限个闭点组成的集合（\'etale 态射的纤维是离散的）。
因此把 $X$ 换成 $x$ 的一个仿射开邻域后，可设
$y \not \in f(V)$。我们要证明 $V$ 仿射。

\medskip\noindent
令 $e(V)$ 为最大整数 $i$，使 $H^i(V, \mathcal{O}_V) \not = 0$。
由于 $X$ 仿射，整数 $e(V)$ 等于诸数 $e(V_x)$ 的最大值，其中
$x \in X \setminus V$；见《局部上同调》引理
\ref{local-cohomology-lemma-cd-local} 以及《局部上同调》引理
\ref{local-cohomology-lemma-cd} 中对上同调维数的刻画。由《局部上同调》引理
\ref{local-cohomology-lemma-cd-dimension}，有
$e(V_x) \leq \dim(\mathcal{O}_{X, x}) - 1$。若
$\dim(\mathcal{O}_{X, x}) \geq 2$，则分歧轨迹的纯性（引理
\ref{pione-lemma-purity-ramification}）表明 $V_x$ 是
$\mathcal{O}_{X, x}$ 的穿孔谱的真开子概形。由于 $\mathcal{O}_{X, x}$ 正规且优良，
Hartshorne--Lichtenbaum 消失定理（《局部上同调》引理
\ref{local-cohomology-lemma-affine-complement}）蕴含
$e(V_x) \leq \dim(\mathcal{O}_{X, x}) - 2$。另一方面，由于
$X \to Y$ 是有限型的，并且 $V \subset X$ 稠密（必要时把 $X$ 换成
$V$ 的闭包），维数公式（《态射》引理
\ref{morphisms-lemma-dimension-formula}）表明
$\dim(\mathcal{O}_{X, x}) \leq d$。因而
$e(V) \leq \max(0, d -  2)$。于是由引理
\ref{pione-lemma-conclude} 可知 $V$ 仿射，只要 $d \geq 2$。若 $d = 1$ 或 $d = 0$，则
$\mathcal{O}_{Y, y}$ 的穿孔谱仿射，故 $V$ 仿射。
\end{proof}




\section{光滑固有情形下的特化映射}
\label{pione-section-specialization-smooth-proper}

\noindent
本节讨论如下结果。设 $f : X \to S$ 是概形的固有光滑态射。设
$s \leadsto s'$ 是 $S$ 中点的一个特化。则特化映射
$$
sp : \pi_1(X_{\overline{s}}) \longrightarrow \pi_1(X_{\overline{s}'})
$$
（见节 \ref{pione-section-specialization-map}）是满射，并且
\begin{enumerate}
\item 若 $\kappa(s')$ 的特征为零，则它是同构；或者
\item 若 $\kappa(s')$ 的特征为 $p > 0$，则它在最大与 $p$ 互素商上诱导同构。
\end{enumerate}

\begin{lemma}
\label{pione-lemma-specialization-map-surjective}
设 $f : X \to S$ 是具有几何连通纤维的平坦固有态射。设
$s' \leadsto s$ 是一个特化。若 $X_s$ 几何既约，则特化映射
$sp : \pi_1(X_{\overline{s}'}) \to \pi_1(X_{\overline{s}})$ 是满射。
\end{lemma}

\begin{proof}
由于 $X_s$ 几何既约，必要时缩小 $S$ 后，可以假设所有纤维都几何既约；见
《态射补篇》引理 \ref{more-morphisms-lemma-geometrically-reduced-open}。令
$\mathcal{O}_{S, s} \to A \to \kappa(\overline{s}')$ 如特化映射的构造中那样；
见节 \ref{pione-section-specialization-map}。因此只须证明
$$
\pi_1(X_{\overline{s}'}) \to \pi_1(X_A)
$$
是满射。这由命题 \ref{pione-proposition-first-homotopy-sequence} 和
$\pi_1(\Spec(A)) = \{1\}$ 得出。
\end{proof}

\begin{proposition}
\label{pione-proposition-specialization-map-isomorphism}
设 $f : X \to S$ 是具有几何连通纤维的光滑固有态射。设
$s' \leadsto s$ 是一个特化。若 $\kappa(s)$ 的特征为零，则特化映射
$$
sp : \pi_1(X_{\overline{s}'}) \to \pi_1(X_{\overline{s}})
$$
是同构。
\end{proposition}

\begin{proof}
由引理 \ref{pione-lemma-specialization-map-surjective}，该映射是满射。因此只须证明它是单射。

\medskip\noindent
可设 $S$ 仿射。于是 $S$ 是一族在 $\mathbf{Z}$ 上有限型的仿射概形的余滤过极限。
因此可以假设 $X \to S$ 是 $X_0 \to S_0$ 的基变换，其中 $S_0$ 是一个有限型
$\mathbf{Z}$-代数的谱，而 $X_0 \to S_0$ 光滑且固有。见《极限》引理
\ref{limits-lemma-descend-finite-presentation}、
\ref{limits-lemma-descend-smooth} 和
\ref{limits-lemma-eventually-proper}。由引理
\ref{pione-lemma-specialization-map-base-change}，归约到基为 Noether 的情形。

\medskip\noindent
应用引理 \ref{pione-lemma-specialization-map-discrete-valuation-ring}，归约到如下情形：
基 $S$ 是严格 Hensel 离散赋值环 $A$ 的谱，并且所考察的是 $A$ 上的特化映射。
令 $K$ 为 $A$ 的分式域。选取代数闭包 $\overline{K}$，它对应于几何一般点
$\overline{\eta}$（在 $\Spec(A)$ 上）。对有限可分的
$\overline{K}/L/K$，令 $B \subset L$ 为 $A$ 在 $L$ 中的整闭包。由《代数补篇》注
\ref{more-algebra-remark-finite-separable-extension}，这是离散赋值环。

\medskip\noindent
设 $X \to \Spec(A)$ 如上一段。为了证明特化映射的单射性，只须证明每个有限
\'etale 覆盖 $V$（在 $X_{\overline{\eta}}$ 上）都是某个有限 \'etale 覆盖
$Y \to X$ 的基变换。事实上，此时由引理
\ref{pione-lemma-functoriality-galois-injective}，
$\pi_1(X_{\overline{\eta}}) \to \pi_1(X) = \pi_1(X_s)$ 是单射。

\medskip\noindent
给定 $V$，先由引理 \ref{pione-lemma-perfection} 把 $V$ 下降为
$V' \to X_{K^{sep}}$，再由引理 \ref{pione-lemma-limit} 下降为
$V'' \to X_L$。令 $Z \to X_B$ 为 $X_B$ 在 $V''$ 中的正规化。注意，
$Z$ 正规，并且有 $Z_L = V''$（作为 $X_L$ 上的概形）。因此
$Z \to X_B$ 在一般纤维上有限 \'etale。问题在于尚不知道
$Z \to X_B$ 是否处处 \'etale。由于 $X \to \Spec(A)$ 具有几何连通的光滑纤维，
可见特殊纤维 $X_s$ 几何不可约。因此
$X_B \to \Spec(B)$ 的特殊纤维不可约；令 $\xi_B$ 为其一般点。令
$\xi_1, \ldots, \xi_r$ 为 $Z$ 中映到 $\xi_B$ 的点。现在第一个问题
（事实将表明也是唯一的问题）是离散赋值环的扩张
$$
\mathcal{O}_{X_B, \xi_B} \subset \mathcal{O}_{Z, \xi_i}
$$
可能分歧。令 $e_i$ 为该扩张的分歧指数。注意，由于
$\kappa(s)$ 的特征为零，该分歧是温和的！

\medskip\noindent
为消除分歧，选取更进一步的有限可分扩张 $K^{sep}/L'/L/K$，使分歧指数
$e$（对应诱导扩张 $B'/B$）可被 $e_i$ 整除。考虑正规化基变换
$Z'$，它由 $Z$ 沿 $\Spec(B') \to \Spec(B)$ 得到；见《态射补篇》节
\ref{more-morphisms-section-reduced-fibre-theorem} 中的讨论。令
$\xi_{i, j}$ 为 $Z'$ 中既映到 $\xi_{B'}$ 又映到 $\xi_i$（在 $Z$ 中）的点。则局部环
$$
\mathcal{O}_{Z', \xi_{i, j}}
$$
是 $\mathcal{O}_{Z, \xi_i}$ 在 $L' \otimes_L F_i$ 中的整闭包的局部化，其中
$F_i$ 是 $\mathcal{O}_{Z, \xi_i}$ 的分式域；略去细节。因此 Abhyankar 引理
（《代数补篇》引理 \ref{more-algebra-lemma-abhyankar}）表明
$$
\mathcal{O}_{X_{B'}, \xi_{B'}} \subset \mathcal{O}_{Z', \xi_{i, j}}
$$
不分歧。由此可知态射 $Z' \to X_{B'}$ 在所有余维数 $1$ 的点处为 \'etale。
于是由分支轨迹的纯性（引理 \ref{pione-lemma-purity}），可见
$Z' \to X_{B'}$ 是有限 \'etale 的！

\medskip\noindent
不过，$A \to B'$ 诱导的剩余域扩张是平凡的（因为 $A$ 的剩余域可分闭且特征为零，
所以代数闭）。两次应用引理
\ref{pione-lemma-finite-etale-on-proper-over-henselian}，可知 $Z'$ 是某个有限
\'etale 覆盖 $Y \to X$ 的基变换（第一次得到 $Y$（在 $A$ 上），第二次证明拉回到
$B$ 后与 $Z'$ 同构）。这就完成了证明。
\end{proof}

\noindent
设 $G$ 是 profinite 群。设 $p$ 是素数。{\it 最大与 $p$ 互素商}按定义为
$$
G' = \lim_{U \subset G\text{ open, normal, index prime to }p} G/U
$$
若 $X$ 是连通概形且给定了 $p$，则最大与 $p$ 互素商（即 $\pi_1(X)$ 的相应商）
记为 $\pi'_1(X)$。

\begin{theorem}
\label{pione-theorem-specialization-map-isomorphism-prime-to-p}
设 $f : X \to S$ 是具有几何连通纤维的光滑固有态射。设
$s' \leadsto s$ 是一个特化。若 $\kappa(s)$ 的特征为 $p$，则特化映射
$$
sp : \pi_1(X_{\overline{s}'}) \to \pi_1(X_{\overline{s}})
$$
是满射，并诱导同构
$$
\pi'_1(X_{\overline{s}'}) \cong \pi'_1(X_{\overline{s}})
$$
于最大与 p 互素商之间。
\end{theorem}

\begin{proof}
证明与命题 \ref{pione-proposition-specialization-map-isomorphism} 完全相同，但有以下区别：
\begin{enumerate}
\item 给定 $X/A$，不再证明函子
$\textit{F\'Et}_X \to \textit{F\'Et}_{X_{\overline{\eta}}}$
本质满。只证明 Galois 群的阶与 $p$ 互素的 Galois 对象位于本质像中。
这足以用完全相同的论证推出
$\pi'_1(X_{\overline{s}'}) \to \pi'_1(X_{\overline{s}})$ 是单射。
\item 扩张
$\mathcal{O}_{X_B, \xi_B} \subset \mathcal{O}_{Z, \xi_i}$
温和分歧，因为相应的分式域扩张是 Galois 扩张，且 Galois 群的阶与 $p$ 互素。
见《代数补篇》引理 \ref{more-algebra-lemma-galois-conclusion}。
\item 扩张 $\kappa_B/\kappa_A$ 未必仍然平凡，但它是纯不可分的。因此态射
$X_{\kappa_B} \to X_{\kappa_A}$ 是泛同胚，并由命题
\ref{pione-proposition-universal-homeomorphism} 在基本群上诱导同构。
\end{enumerate}
\end{proof}













\section{温和分歧}
\label{pione-section-tame}

\noindent
设 $X \to Y$ 是 $\mathbf{Z}$ 上有限型概形之间的有限 \'etale 态射。
关于 $f$ 在 $\infty$ 处温和分歧的含义，有多种定义。文献
\cite{Kerz-Schmidt} 讨论了这些概念在何种程度上相互等价。

\medskip\noindent
本节讨论另一个更初等的问题，它先于无穷远处温和性的概念。
请与文献 \cite{Grothendieck-Murre} 中略有不同的讨论比较。
假设给定
\begin{enumerate}
\item 局部 Noether 概形 $X$；
\item 稠密开子概形 $U \subset X$；
\item 有限 \'etale 态射 $f : Y \to U$，
\end{enumerate}
使得对每个素除子 $Z \subset X$，若 $Z \cap U = \emptyset$，则局部环
$\mathcal{O}_{X, \xi}$，即 $X$ 在一般点 $\xi$（$Z$ 的一般点）处的局部环，
都是离散赋值环。令 $K_\xi$ 为 $\mathcal{O}_{X, \xi}$ 的分式域，
便得到概形的 Cartesian 方形
$$
\xymatrix{
\Spec(K_\xi) \ar[r] \ar[d] & U \ar[d] \\
\Spec(\mathcal{O}_{X, \xi}) \ar[r] & X
}
$$
特别地，$Y \times_U \Spec(K_\xi)$ 是某个有限可分代数
$L_\xi/K_\xi$ 的谱。
此时称 {\it $Y$ 在 $X$ 上于余维数 $1$ 处不分歧}，分别地称
{\it $Y$ 在 $X$ 上于余维数 $1$ 处温和分歧}，若
$L_\xi/K_\xi$ 关于 $\mathcal{O}_{X, \xi}$ 对每个上述
$(Z, \xi)$ 不分歧，分别地温和分歧；见
《代数补篇》定义 \ref{more-algebra-definition-types-of-extensions}。
更确切地说，将 $L_\xi$ 分解为 $K_\xi$ 的有限可分域扩张之积，
并要求每个因子关于 $\mathcal{O}_{X, \xi}$ 不分歧，分别地温和分歧。

\begin{lemma}
\label{pione-lemma-pullback-tame-codim1}
设 $X' \to X$ 是局部 Noether 概形的态射，$U \subset X$ 是稠密开子概形。假设
\begin{enumerate}
\item $U' = f^{-1}(U)$ 是 $X'$ 中的稠密开子概形；
\item 对每个素除子 $Z \subset X$，若 $Z \cap U = \emptyset$，则局部环
$\mathcal{O}_{X, \xi}$，即 $X$ 在一般点 $\xi$（$Z$ 的一般点）处的局部环，是离散赋值环；
\item 对每个素除子 $Z' \subset X'$，若 $Z' \cap U' = \emptyset$，则局部环
$\mathcal{O}_{X', \xi'}$，即 $X'$ 在一般点 $\xi'$（$Z'$ 的一般点）处的局部环，是离散赋值环；
\item 若 $\xi' \in X'$ 如 (3) 所述，则 $\xi = f(\xi')$ 如 (2) 所述。
\end{enumerate}
若 $f : Y \to U$ 有限 \'etale，且 $Y$ 在 $X$ 上于余维数 $1$ 处不分歧，
分别地温和分歧，则 $Y' = Y \times_X X' \to U'$ 有限 \'etale，且
$Y'$ 在 $X'$ 上于余维数 $1$ 处不分歧，分别地温和分歧。
\end{lemma}

\begin{proof}
本引理中唯一值得说明的事实是《代数补篇》引理
\ref{more-algebra-lemma-tame-goes-up} 给出的交换代数结果。
\end{proof}

\noindent
利用刚引入的术语，可以把先前得到的纯性结果重新表述为如下简洁形式。

\begin{lemma}
\label{pione-lemma-purity-one-divisor-general}
设 $X$ 是局部 Noether 概形，$U \subset X$ 开且稠密，并设
$Y \to U$ 是有限 \'etale 态射。假设
\begin{enumerate}
\item $Y$ 在 $X$ 上于余维数 $1$ 处不分歧；
\item 局部环 $\mathcal{O}_{X, x}$ 对所有 $x \in X \setminus U$ 都正则。
\end{enumerate}
则存在有限 \'etale 态射 $Y' \to X$，其在 $U$ 上的限制为 $Y$。
\end{lemma}

\begin{proof}
令 $\xi \in X \setminus U$ 为 $X \setminus U$ 的某个余维数 $1$ 不可约分支的一般点。
则 $\mathcal{O}_{X, \xi}$ 是离散赋值环。如上文讨论，记
$Y \times_U \Spec(K_\xi) = \Spec(L_\xi)$.
以 $B_\xi$ 表示 $\mathcal{O}_{X, \xi}$ 在 $L_\xi$ 中的整闭包。
$Y$ 在 $X$ 上于余维数 $1$ 处不分歧这一假设意味着
$\mathcal{O}_{X, \xi} \to B_\xi$ 有限 \'etale。因此得到有限 \'etale 态射
$Y_\xi \to \Spec(\mathcal{O}_{X, \xi})$ 以及同构
$$
Y \times_U \Spec(K_\xi) \cong
Y_\xi \times_{\Spec(\mathcal{O}_{X, \xi})} \Spec(K_\xi)
$$
（在 $\Spec(K_\xi)$ 上）。由《极限》引理
\ref{limits-lemma-glueing-near-point}，可以找到开子概形
$U \subset U' \subset X$（它包含 $\xi$）以及有限表现态射 $Y' \to U'$，使其在 $U$ 上的限制恢复
$Y$，而在 $\Spec(\mathcal{O}_{X, \xi})$ 上的限制恢复 $Y_\xi$。
再由《极限》引理 \ref{limits-lemma-glueing-near-point-properties}，态射
$Y' \to U'$ 在必要时把 $U'$ 缩小为更小的开集后有限 \'etale。
对 $X \setminus U$ 的其他余维数 $1$ 分支的一般点重复此论证，可以假设已有
有限 \'etale 态射 $Y' \to U'$，它把 $Y \to U$ 延拓到某个开子概形
$U' \subset X$，该开子概形包含 $U$ 以及所有余维数 $1$ 的 $X \setminus U$ 中的点。
最后对 $Y' \to U'$ 应用引理 \ref{pione-lemma-extend-pure}。事实上，所有局部环
$\mathcal{O}_{X, x}$（其中 $x \in X \setminus U'$）都正则且满足
$\dim(\mathcal{O}_{X, x}) \geq 2$。故由引理 \ref{pione-lemma-local-purity}，
$\mathcal{O}_{X, x}$ 具有纯性。
\end{proof}

\begin{lemma}
\label{pione-lemma-purity-one-divisor}
设 $X$ 是局部 Noether 概形，$D \subset X$ 是有效 Cartier 除子，且
$D$ 为正则概形。设 $Y \to X \setminus D$ 是有限 \'etale 态射。
若 $Y$ 在 $X$ 上于余维数 $1$ 处不分歧，则存在有限 \'etale 态射
$Y' \to X$，其在 $X \setminus D$ 上的限制为 $Y$。
\end{lemma}

\begin{proof}
这是引理 \ref{pione-lemma-purity-one-divisor-general} 的特殊情形。首先，
$D$ 在 $X$ 中处处不稠密（见《除子》节
\ref{divisors-section-effective-Cartier-divisors} 中的讨论），故
$X \setminus D$ 在 $X$ 中稠密。其次，环 $\mathcal{O}_{X, x}$ 对所有
$x \in D$ 都是正则局部环；这是由《代数》引理
\ref{algebra-lemma-regular-mod-x} 以及 $\mathcal{O}_{D, x}$ 正则这一假设得出的。
\end{proof}

\begin{example}[标准温和分歧态射]
\label{pione-example-tamely-ramified}
设 $A$ 是 Noether 环，$f \in A$ 是非零因子，且 $A/fA$ 既约。
这意味着 $A_\mathfrak p$ 是以 $f$ 为一致化元的离散赋值环，其中
$\mathfrak p$ 是 $f$ 上方的任意极小素理想。设 $e \geq 1$ 是在 $A$ 中可逆的整数。令
$$
C = A[x]/(x^e - f)
$$
则 $\Spec(C) \to \Spec(A)$ 是有限局部自由态射，并在 $A_f$ 的谱上 \'etale。
有限 \'etale 态射
$$
\Spec(C_f) \longrightarrow \Spec(A_f)
$$
在 $\Spec(A)$ 上于余维数 $1$ 处温和分歧。温和性立即由《代数补篇》引理
\ref{more-algebra-lemma-characterize-tame} 对温和分歧扩张的刻画得出。
\end{example}

\noindent
下面给出适用于正则除子的 Abhyankar 引理版本。

\begin{lemma}[正则除子的 Abhyankar 引理]
\label{pione-lemma-abhyankar-one-divisor}
设 $X$ 是局部 Noether 概形，$D \subset X$ 是有效 Cartier 除子，且
$D$ 为正则概形。设 $Y \to X \setminus D$ 是有限 \'etale 态射。
若 $Y$ 在 $X$ 上于余维数 $1$ 处温和分歧，则在 $X$ 上 \'etale 局部，
态射 $Y \to X$ 是例 \ref{pione-example-tamely-ramified} 所述标准温和分歧态射的有限不交并。
\end{lemma}

\begin{proof}
对每个 $x \in X$，将找出一个 \'etale 邻域 $(U, u) \to (X, x)$，其中
$U = \Spec(A)$，使得基变换 $Y \times_X U \to U$ 是例
\ref{pione-example-tamely-ramified} 中标准温和分歧态射的有限不交并。
可设 $x \in D$；当 $x \not \in D$ 时，结论由《\'Etale 态射》引理
\ref{etale-lemma-finite-etale-etale-local} 以及在例
\ref{pione-example-tamely-ramified} 中取 $e = 1$、$f = 1$ 得出。

\medskip\noindent
本段把问题归约到 $X$ 为严格 Hensel 局部环之谱且 $x$ 为闭点的情形。
具体地，缩小 $X$ 后可以假设 $X = \Spec(A)$，而 $D \subset X$ 由非零因子
$f \in A$ 给出。此时 $Y$ 作为 $\Spec(A_f)$ 的有限 \'etale 覆盖是仿射的；记
$Y = \Spec(B)$。令 $A^{sh}$ 为 $\mathcal{O}_{X, x}$ 的严格 Hensel 化。
由于 $A \to A^{sh}$ 平坦且 $x \in D$，可见 $f$ 映为 $A^{sh}$ 的极大理想中的非零因子。
注意，$A^{sh}/fA^{sh} = (A/fA)^{sh}$（《代数》引理
\ref{algebra-lemma-quotient-strict-henselization}）是 $\mathcal{O}_{D, x}$
的严格 Hensel 化，故为正则局部环；见《代数补篇》引理
\ref{more-algebra-lemma-henselization-regular}。由引理
\ref{pione-lemma-pullback-tame-codim1}，$Y$ 到 $\Spec(A^{sh})$ 的基变换在余维数
$1$ 处温和分歧。假设已经对 $Y$ 到 $\Spec(A^{sh})$ 的基变换证明了断言。由于 $A^{sh}$ 严格 Hensel，
每个 \'etale 邻域都有截面，从而得到同构
$$
A^{sh} \otimes_A B \cong
\prod\nolimits_{i \in I} A^{sh}_f[x]/(x^{e_i} - f)
$$
其中 $I$ 是有限集，每个 $e_i$ 都是在 $A^{sh}$ 中可逆的整数。把
$A^{sh} = \colim A_\lambda$ 写成 \'etale $A$-代数 $A_\lambda$ 的滤过余极限。
上示同构下降为同构
$$
A_\lambda \otimes_A B \cong \prod\nolimits_{i \in I}
(A_\lambda)_f[x]/(x^{e_i} - f)
$$
（对充分大的 $\lambda$）；例如见《代数》引理
\ref{algebra-lemma-colimit-category-fp-algebras}。再稍微增大 $\lambda$，
可以假设 $e_i$ 在 $A_\lambda$ 中可逆。于是
$\Spec(A_\lambda) \to \Spec(A)$ 就是所求的 $x$ 的 \'etale 邻域。

\medskip\noindent
现在设 $X = \Spec(A)$，其中 $A$ 是严格 Hensel 局部环，而 $x \in X$
对应于极大理想 $\mathfrak m$，这是 $A$ 的极大理想。令 $f \in \mathfrak m$ 为切出 $D$ 的非零因子。
则 $A/fA$ 正则，由《代数》引理 \ref{algebra-lemma-regular-mod-x}，
$A$ 也正则。将使用正则局部环的一些性质，例如它们是正规整环；见《代数》引理
\ref{algebra-lemma-regular-domain} 和 \ref{algebra-lemma-regular-normal}。
特别地，$A$ 是整环。如上可见 $Y = \Spec(B)$ 仿射，且 $A_f \to B$ 有限
\'etale。故由《代数》引理 \ref{algebra-lemma-normal-goes-up}，$B$ 也正规；
再由《代数》引理 \ref{algebra-lemma-characterize-reduced-ring-normal}，
$Y$ 是若干正规整环之谱的有限不交并。因此也可假设 $B$ 是整环。
由于 $A$ 和 $A/fA$ 都是整环，$f$ 是 $A$ 的素元，并生成高度为 $1$ 的素理想
$\mathfrak p = fA$。注意，$A_\mathfrak p$ 是以 $f$ 为一致化元的离散赋值环。

\medskip\noindent
令 $K$ 为 $A$ 的分式域，$L$ 为 $B$ 的分式域。温和分歧假设意味着
$L$ 关于 $A_\mathfrak p$ 温和分歧。按《代数补篇》注
\ref{more-algebra-remark-finite-separable-extension}，令 $e \geq 1$ 为
$L$ 在 $A_\mathfrak p$ 上各分歧指数之积。于是 $e$ 在
$\kappa(\mathfrak p)$ 中可逆，但此时尚不知道 $e$ 在 $A/fA$ 或 $A$ 中是否可逆。

\medskip\noindent
考虑有限自由 $A$-代数
$$
A' = A[x]/(x^e - f)
$$
它是严格 Hensel 局部环的局部有限扩张，因而仍严格 Hensel。
注意，$f' = x$ 是 $A'$ 中的非零因子，而 $A'/f'A' \cong A/fA$ 是正则环。
因此如上，$A'$ 正则，进而是整环。令
$B' = B \otimes_A A' = B \otimes_{A_f} A'_{f'}$.
由 Abhyankar 引理（《代数补篇》引理 \ref{more-algebra-lemma-abhyankar}），
$\Spec(B')$ 在 $\Spec(A')$ 上于余维数 $1$ 处不分歧。事实上，由引理
\ref{pione-lemma-pullback-tame-codim1}，$\Spec(B')$ 在 $\Spec(A')$ 上于余维数
$1$ 处至少仍温和分歧，而 Abhyankar 引理说明所有分歧指数都已变为 $1$。
由引理 \ref{pione-lemma-purity-one-divisor}，$\Spec(B') \to \Spec(A'_{f'})$
延拓为有限 \'etale 态射 $\Spec(C) \to \Spec(A')$。然而 $A'$ 严格 Hensel，
所以 $\Spec(C)$ 是若干份 $\Spec(A')$ 的有限不交并。因而至少存在一个态射
$\Spec(A'_{f'}) \to \Spec(B')$，从而存在包含
$B \subset A'_{f'}$，它是 $A_f$-代数的包含。

\medskip\noindent
于是有包含 $K \subset L \subset K[f^{1/e}]$。由《代数补篇》引理
\ref{more-algebra-lemma-pull-root-uniformizer}，$L$ 等于
$K[f^{1/n}]$，其中 $n$ 是 $e$ 的某个因子。$B$ 的正规性随即推出
$B \cong A_f[y]/(y^n - f)$；略去细节。不过，环映射
$A_f \to A_f[y]/(y^n - f)$ 的分歧轨迹由 $ny^{n - 1}$ 切出，
这必然意味着 $n$ 在 $A_f$ 中为单位。由于 $f$ 是素元，这意味着
$n = u f^s$，其中 $u$ 是 $A$ 的某个单位，$s \in \mathbf{Z}$。
又因 $n$ 在 $\kappa(\mathfrak p)$ 中映为单位，故 $s = 0$，证明完成。
\end{proof}

\begin{lemma}
\label{pione-lemma-extend-tame-covering-normal}
在引理 \ref{pione-lemma-abhyankar-one-divisor} 的情形中，$X$ 在 $Y$ 中的正规化是
有限局部自由态射 $\pi : Y' \to X$，并且
\begin{enumerate}
\item $Y'$ 在 $X \setminus D$ 上的限制同构于 $Y$；
\item $D' = \pi^{-1}(D)_{red}$ 是 $Y'$ 上的有效 Cartier 除子；
\item $D'$ 是正则概形。
\end{enumerate}
此外，在 $X$ 上 \'etale 局部，态射 $Y' \to X$ 是下列态射的有限不交并：
$$
\Spec(A[x]/(x^e - f)) \to \Spec(A)
$$
其中 $A$ 是 Noether 环，$f \in A$ 是非零因子，$A/fA$ 正则，并且
$e \geq 1$ 在 $A$ 中可逆。
\end{lemma}

\begin{proof}
这只是引理 \ref{pione-lemma-abhyankar-one-divisor} 的补充；事实上，读过该引理的证明后，
本引理几乎立即成立。也可以直接从引理 \ref{pione-lemma-abhyankar-one-divisor} 的结论推出。
具体地，取 $X$ 在 $Y$ 中的正规化与 \'etale 基变换可交换；见《态射补篇》引理
\ref{more-morphisms-lemma-normalization-smooth-localization}。所以关于
$Y' \to X$ 局部结构的断言可以在 $X$ 上 \'etale 局部证明。由引理
\ref{pione-lemma-abhyankar-one-divisor}，可以假设 $X = \Spec(A)$，其中 $A$ 是
Noether 环；存在非零因子 $f\in A$ 使 $A/fA$ 正则；并且 $Y$ 是若干环
$A_f[x]/(x^e - f)$ 之谱的有限不交并，其中 $e$ 在 $A$ 中可逆。
略去如下验证：$A$ 在 $A_f[x]/(x^e - f)$ 中的整闭包等于
$A' = A[x]/(x^e - f)$。（为此，只须证明 $A'$ 在位于 $(f)$ 上方的素理想处的
局部化正则。）其余细节从略。
\end{proof}

\begin{lemma}
\label{pione-lemma-tame-covering-split}
在引理 \ref{pione-lemma-abhyankar-one-divisor} 的情形中，令 $Y' \to X$ 如引理
\ref{pione-lemma-extend-tame-covering-normal} 所述。设 $R$ 是分式域为 $K$ 的离散赋值环。令
$$
t : \Spec(R) \to X
$$
为一个态射，使得概形论逆像 $t^{-1}D$ 是 $\Spec(R)$ 的约化闭点。
\begin{enumerate}
\item 若 $t|_{\Spec(K)}$ 可提升为 $Y$ 的一点，则得到提升
$t' : \Spec(R) \to Y'$，并且 $Y' \to X$ 沿 $t'(\Spec(R))$ 为 \'etale；
\item 若 $\Spec(K) \times_X Y$ 同构于若干份 $\Spec(K)$ 的不交并，
则 $Y' \to X$ 在 $t(\Spec(R))$ 的某个开邻域上有限 \'etale。
\end{enumerate}
\end{lemma}

\begin{proof}
把固有性的赋值判据应用于有限态射 $Y' \to X$，可知
$\Spec(K)$-值点（在 $Y$ 中）若与 $t|_{\Spec(K)}$ 作为到 $X$ 的态射相符，便唯一提升为态射
$t' : \Spec(R) \to Y'$。因而陈述 (1) 有意义。

\medskip\noindent
选取 \'etale 邻域 $(U, u) \to (X, t(\mathfrak m_R))$，使
$U = \Spec(A)$，并使 $Y' \times_X U \to U$ 对某个 $f \in A$ 具有引理
\ref{pione-lemma-extend-tame-covering-normal} 中的描述。于是
$\Spec(R) \times_X U \to \Spec(R)$ \'etale 且满。若 $R'$ 表示
$\Spec(R) \times_X U$ 在 $\Spec(R)$ 闭点上方的局部环，则 $R'$ 是离散赋值环，
且 $R \subset R'$ 是离散赋值环的不分歧扩张（《代数补篇》引理
\ref{more-algebra-lemma-Dedekind-etale-extension}）。关于 $t$ 的假设意味着映射
$A \to R'$（它对应于
$$
\Spec(R') \to \Spec(R) \times_X U \to U
$$
）把 $f$ 映为一致化元 $\pi \in R'$。现设
$$
Y' \times_X U =
\coprod\nolimits_{i \in I} \Spec(A[x]/(x^{e_i} - f))
$$
其中 $e_i \geq 1$。于是
$$
\Spec(R') \times_U (Y' \times_X U) =
\coprod\nolimits_{i \in I} \Spec(R'[x]/(x^{e_i} - \pi))
$$
环 $R'[x]/(x^{e_i} - f)$ 是离散赋值环（《代数补篇》引理
\ref{more-algebra-lemma-pull-root-uniformizer}），所以它们不存在到 $R'$ 的分式域的映射，
除非 $e_i = 1$。

\medskip\noindent
(1) 的证明。在此情形中，映射 $t' : \Spec(R) \to Y'$ 经基变换给出相应映射
$t'' : \Spec(R') \to Y' \times_X U$。由上文讨论，它必映入某个对应于
$i \in I$ 且 $e_i = 1$ 的分支。因此显然 $Y' \times_X U \to U$ 沿
$t''$ 的像为 \'etale。\'etale 性可在 \'etale 基变换后检验，故 (1) 得证。

\medskip\noindent
(2) 的证明。在此情形中，假设蕴含 $e_i = 1$，对所有 $i \in I$ 均如此。
因此 $Y' \times_X U \to U$ 有限 \'etale，结论同上。
\end{proof}

\begin{lemma}
\label{pione-lemma-extend-covering}
设 $S$ 是一般点为 $\eta$ 的正规整 Noether 概形。设 $f : X \to S$ 是具有几何连通
纤维的光滑态射。设 $\sigma : S \to X$ 是 $f$ 的截面。设 $Z \to X_\eta$ 是有限
\'etale Galois 覆盖（节 \ref{pione-section-finite-etale-under-galois}），其群 $G$ 的阶在
$S$ 上可逆，并且 $Z$ 有一个 $\kappa(\eta)$-理性点映到 $\sigma(\eta)$。
则存在有限 \'etale Galois 覆盖 $Y \to X$，其群为 $G$，且其在
$X_\eta$ 上的限制为 $Z$。
\end{lemma}

\begin{proof}
先设 $S = \Spec(R)$ 是离散赋值环 $R$ 的谱，闭点为 $s \in S$。于是
$X_s$ 是 $X$ 中的有效 Cartier 除子，并且作为域上的光滑概形，$X_s$ 正则。
此外，一般纤维 $X_\eta$ 就是开子概形 $X \setminus X_s$。由《代数补篇》引理
\ref{more-algebra-lemma-galois-conclusion} 及关于 $G$ 的假设，$Z$ 在 $X$ 上于
余维数 $1$ 处温和分歧。令 $Z' \to X$ 如引理
\ref{pione-lemma-extend-tame-covering-normal} 所述。注意，$G$ 在 $Z$ 上的作用延拓为
$G$ 在 $Z'$ 上的作用。由引理 \ref{pione-lemma-tame-covering-split}，
$Z' \to X$ 在 $\sigma(y)$ 的某个开邻域上有限 \'etale。由于 $X_s$ 不可约，
这蕴含 $Z \to X_\eta$ 在 $X$ 上于余维数 $1$ 处不分歧。于是得到有限 \'etale 态射
$Y \to X$，其在 $X_\eta$ 上的限制为 $Z$；这是由引理
\ref{pione-lemma-purity-one-divisor} 得出的。当然 $Y \cong Z'$（略去细节；提示：在 \'etale
局部计算），故 $Y$ 是群为 $G$ 的 Galois 覆盖。

\medskip\noindent
一般情形。令 $U \subset S$ 为满足如下条件的极大开子概形：存在有限
\'etale Galois 覆盖 $Y \to X \times_S U$，其群为 $G$，且其在 $X_\eta$ 上的限制同构于 $Z$。
为导出矛盾，假设 $U \not = S$。令 $s \in S \setminus U$ 为
$S \setminus U$ 某个不可约分支的一般点。则逆像 $U_s$（即 $U$ 在
$\Spec(\mathcal{O}_{S, s})$ 中的逆像）是
$\mathcal{O}_{S, s}$ 的穿孔谱。我们断言
$Y \times_S U_s \to X \times_S U_s$ 是有限 \'etale Galois 覆盖
$Y'_s \to X \times_S \Spec(\mathcal{O}_{S, s})$ 的限制，该覆盖的群为 $G$。

\medskip\noindent
先说明该断言如何产生所需矛盾。由《极限》引理
\ref{limits-lemma-glueing-near-point}，可找到开子概形
$U \subset U' \subset S$（它包含 $s$）以及有限表现态射 $Y'' \to U'$，使其在 $U$ 上的限制
恢复 $Y' \to U$，而在 $\Spec(\mathcal{O}_{S, s})$ 上的限制恢复 $Y'_s$。
此外，由该引理给出的范畴等价，缩小 $U'$ 后可假设存在态射
$Y'' \to U' \times_S X$，并且 $G$ 在 $Y''$ 上、在 $U' \times_S X$ 上方作用；
这些态射与作用经基变换到 $U$ 和 $\Spec(\mathcal{O}_{S, s})$ 后与给定者相容。
必要时进一步缩小 $U'$，可假设 $Y'' \to U \times_S X$ 有限 \'etale；见
《极限》引理 \ref{limits-lemma-glueing-near-point-properties}。这说明找到了
$S$ 的一个严格更大的开集，使 $Y$ 在其上延拓为群为 $G$ 的有限 \'etale Galois
覆盖，从而得到所需矛盾。

\medskip\noindent
断言的证明。替换 $S$，所用者为 $\Spec(\mathcal{O}_{S, s})$。于是
$S = \Spec(A)$，其中 $(A, \mathfrak m)$ 是正规局部整环；$U \subset S$
为穿孔谱，并已有有限 \'etale Galois 覆盖
$Y \to X \times_S U$，其群为 $G$。若 $\dim(A) = 1$，则由证明第一段可把 $Y$ 延拓为
$X$ 的 Galois 覆盖。因此可设 $\dim(A) \geq 2$，从而有
$\text{depth}(A) \geq 2$，因为 $S$ 正规；见《代数》引理
\ref{algebra-lemma-criterion-normal}。由于 $X \to S$ 平坦，可知
$\text{depth}(\mathcal{O}_{X, x}) \geq 2$，对每个点 $x \in X$（它映到 $s$）均如此；
见《代数》引理 \ref{algebra-lemma-apply-grothendieck}。令
$$
Y' \longrightarrow X
$$
为用 $Y \to X \times_S U$ 在引理 \ref{pione-lemma-extend-S2} 中构造的有限态射。
注意，得到 $G$ 在 $Y$ 上的典范作用。因此只剩证明 $Y'$ 在 $X$ 上 \'etale。
事实上，例如由引理 \ref{pione-lemma-purity-smooth-over-depth2}，只须证明
$Y' \to X$ 在纤维 $X_s$ 的（唯一）一般点上方为 \'etale。为此，在
$\sigma(s)$ 的一个（形式）邻域中作局部计算。

\medskip\noindent
选取仿射开子概形 $\Spec(B) \subset X$，它包含 $\sigma(s)$。则
$A \to B$ 是具有截面 $\sigma : B \to A$ 的光滑环映射。记
$I = \Ker(\sigma)$，并以 $B^\wedge$ 表示 $I$-进的 $B$ 的完备化。由《代数》引理
\ref{algebra-lemma-section-smooth}，有
$B^\wedge \cong A[[x_1, \ldots, x_d]]$，其中 $d \geq 0$。当然 $B \to B^\wedge$ 平坦
（《代数》引理 \ref{algebra-lemma-completion-flat}），且
$\Spec(B^\wedge) \to X$ 的像包含 $X_s$ 的一般点。令
$V \subset \Spec(B^\wedge)$ 为 $U$ 的逆像。考虑有限 \'etale 态射
$$
W = Y \times_{(X \times_S U)} V \longrightarrow V
$$
引理 \ref{pione-lemma-extend-S2} 中 $Y'$ 的构造与平坦基变换相容，因此基变换
$Y' \times_X \Spec(B^\wedge) \to \Spec(B^\wedge)$ 是从 $W \to V$
在 $\Spec(B^\wedge)$ 上按引理 \ref{pione-lemma-extend-S2} 的步骤构造的。
令 $V_0 = V \cap V(x_1, \ldots, x_d) \subset V$，并令
$W_0 = W \times_V V_0$。这是正规整概形，它映入 $\sigma(S)$，所用态射为
$\Spec(B^\wedge) \to X$，且事实上与 $\sigma(U)$ 等同。
因此 $W_0 \to V_0 = U$ 完全分裂：其一般纤维具有此性质，因为假设
$Z \to X_\eta$ 有一个 $\kappa(\eta)$-理性点位于 $\sigma(\eta)$ 上方
（而 $G$ 的作用当然随即蕴含整个纤维 $Z_{\sigma(\eta)}$ 是若干份概形
$\eta = \Spec(\kappa(\eta))$ 的不交并）。最后，由引理
\ref{pione-lemma-fully-faithful-power-series-over-depth2} 有
$$
W_0 \times_U V \cong W
$$
这说明 $W$ 是若干份 $V$ 的不交并，因而
$Y' \times_X \Spec(B^\wedge)$ 是若干份 $\Spec(B^\wedge)$ 的不交并，证明完成。
\end{proof}

\begin{lemma}
\label{pione-lemma-extend-covering-general}
设 $S$ 是一般点为 $\eta$ 的拟紧拟分离正规整概形。设 $f : X \to S$ 是具有
几何连通纤维的拟紧拟分离光滑态射。设 $\sigma : S \to X$ 是 $f$ 的截面。
设 $Z \to X_\eta$ 是有限 \'etale Galois 覆盖（节
\ref{pione-section-finite-etale-under-galois}），其群 $G$ 的阶在 $S$ 上可逆，
并且 $Z$ 有一个 $\kappa(\eta)$-理性点映到 $\sigma(\eta)$。
则存在有限 \'etale Galois 覆盖 $Y \to X$，其群为 $G$，且其在 $X_\eta$ 上的限制为 $Z$。
\end{lemma}

\begin{proof}
若 $S$ 为 Noether 概形，这就是引理 \ref{pione-lemma-extend-covering} 的结论。
一般情形由标准极限论证推出。我们强烈建议读者略过本证明。

\medskip\noindent
可把 $S = \lim S_i$ 写成一族概形的有向极限，其中过渡态射仿射，而每个
$S_i$ 都在 $\mathbf{Z}$ 上有限型；见《极限》命题
\ref{limits-proposition-approximate}。对每个 $i$，令
$S \to S'_i \to S_i$ 为 $S_i$ 在 $S$ 中的正规化；见《态射》节
\ref{morphisms-section-normalization-X-in-Y}。结合《代数》命题
\ref{algebra-proposition-ubiquity-nagata} 以及《态射》引理
\ref{morphisms-lemma-nagata-normalization-finite} 和
\ref{morphisms-lemma-normal-normalization}，可知 $S'_i$ 在 $\mathbf{Z}$ 上有限型，
在 $S_i$ 上有限，并且 $S'_i$ 是正规整概形，而 $S \to S'_i$ 占优。
由《态射》引理 \ref{morphisms-lemma-functoriality-normalization}，得到过渡态射
$S'_{i'} \to S'_i$；它们与过渡态射 $S_{i'} \to S_i$ 及源为 $S$ 的态射相容。
断言 $S = \lim S'_i$。略去证明（提示：在 $S_i$ 的一个选定仿射开集上方，
对某个 $i$，考察仿射开集，
把问题转化为直接的代数问题）。因此可把
$S = \lim S_i$ 写成一族正规整概形 $S_i$ 的有向极限，其中过渡态射仿射，且
$S_i$ 在 $\mathbf{Z}$ 上有限型。

\medskip\noindent
对某个 $i$，可找到有限表现光滑态射 $X_i \to S_i$，其到 $S$ 的基变换为
$X \to S$；见《极限》引理 \ref{limits-lemma-descend-finite-presentation} 和
\ref{limits-lemma-descend-smooth}。增大 $i$ 后，可假设截面 $\sigma$ 提升为截面
$\sigma_i : S_i \to X_i$（由《极限》引理
\ref{limits-lemma-descend-finite-presentation} 中的范畴等价）。可把 $X_i$ 替换为
《态射补篇》节 \ref{more-morphisms-section-connected-components} 所研究的开子概形
$X_i^0$，因为 $X \to X_i$ 的像显然映入其中（其开放性见《态射补篇》引理
\ref{more-morphisms-lemma-connected-along-section-open}）。于是可假设
$X_i \to S_i$ 的纤维几何连通。再增大 $i$，可假设 $|G|$ 在 $S_i$ 上可逆。
令 $\eta_i \in S_i$ 为一般点。由于 $X_\eta$ 是概形 $X_{i, \eta_i}$ 的极限，
可用完全相同的论证，把 $Z \to X_\eta$ 下降为某个有限 \'etale Galois 覆盖
$Z_i \to X_{i, \eta_i}$，必要时增大 $i$；见引理 \ref{pione-lemma-limit}。
必要时再次增大 $i$，可假设 $Z_i$ 有一个 $\kappa(\eta_i)$-理性点，
它映到 $\sigma_i(\eta_i)$。此时应用 Noether 情形的引理，再拉回到 $X$，即得结论。
\end{proof}











\section{正特征中的技巧}
\label{pione-section-achinger}

\noindent
Piotr Achinger 的论文 \cite{Achinger} 证明了正特征中的仿射概形总是
$K(\pi, 1)$。本节说明 \cite{Achinger} 中较初等的部分。具体地，证明：
对正特征域 $k$，一个在 $\mathbf{A}^n_k$ 上 \'etale 的仿射概形实际上通过另一个态射
在 $\mathbf{A}^n_k$ 上有限 \'etale。还将证明，正特征中连通仿射概形的闭浸入在
\'etale 基本群上诱导单射。

\medskip\noindent
设 $k$ 是特征为 $p > 0$ 的域。设
$$
k[x_1, \ldots, x_n] \longrightarrow A
$$
是有限型 $k$-代数的满射，其源是关于 $x_1, \ldots, x_n$ 的多项式代数。
以 $I \subset k[x_1, \ldots, x_n]$ 表示核，于是
$A = k[x_1, \ldots, x_n]/I$。不假设 $A$ 非零（换言之，允许 $A$ 为零环且
$I = k[x_1, \ldots, x_n]$）。最后，假设给定有限 \'etale 环映射
$\pi : A \to B$。

\medskip\noindent
假设给定 $k, n, k[x_1, \ldots, x_n] \to A, I, \pi : A \to B$。
设 $C$ 是 $k$-代数。考虑交换图
$$
\xymatrix{
& B \\
C \ar[r] & C/\varphi(I)C \ar[u]^\tau \\
k[x_1, \ldots, x_n] \ar[u]^\varphi \ar[r] &
A \ar[u] \ar@/_3em/[uu]_\pi
}
$$
其中 $\varphi$ 是 \'etale $k$-代数映射，$\tau$ 是满的 $k$-代数映射。
给定 $C, \varphi, \tau$。对任意 $r \geq 0$ 以及元素
$y_1, \ldots, y_r \in C$，若它们生成代数 $C$（在 $\Im(\varphi)$ 上），令
$s = s(r, y_1, \ldots, y_r) \in \{0, \ldots, r\}$ 为满足下列条件的最大元素：
元素 $y_i$ 在 $\Im(\varphi)$ 上整，其中 $1 \leq i \leq s$。定义
$NF(C, \varphi, \tau)$ 为
$r - s = r - s(r, y_1, \ldots, y_r)$ 的最小值，其中遍历所有上述
$r$ 和 $y_1, \ldots, y_r$ 的选择。注意，
$NF(C, \varphi, \tau)$ 为 $0$ 当且仅当 $\varphi$ 有限。

\begin{lemma}
\label{pione-lemma-lower-invariant}
在上述情形中，若 $NF(C, \varphi, \tau) > 0$，则存在 \'etale
$k$-代数映射 $\varphi'$ 和满的 $k$-代数映射 $\tau'$，它们嵌入交换图
$$
\xymatrix{
& B \\
C \ar[r] & C/\varphi'(I)C \ar[u]_{\tau'} \\
k[x_1, \ldots, x_n] \ar[u]^{\varphi'} \ar[r] &
A \ar[u] \ar@/_3em/[uu]_\pi
}
$$
并满足 $NF(C, \varphi', \tau') < NF(C, \varphi, \tau)$。
\end{lemma}

\begin{proof}
选取 $r \geq 0$ 和 $y_1, \ldots, y_r \in C$，它们生成
$C$，作为 $\Im(\varphi)$ 上的代数；再选取 $0 \leq s \leq r$，使得
$y_1, \ldots, y_s$ 在 $\Im(\varphi)$ 上整，并且
$r - s = NF(C, \varphi, \tau) > 0$。由于 $B$ 在 $A$ 上有限，
$y_{s + 1}$ 在 $B$ 中的像满足一个以 $A$ 中元素为系数的首一多项式。
因此可以找到 $d \geq 1$ 和
$f_1, \ldots, f_d \in k[x_1, \ldots, x_n]$，使得
$$
z = y_{s + 1}^d + \varphi(f_1) y_{s + 1}^{d - 1} + \ldots + \varphi(f_d) \in
J = \Ker(C \to C/\varphi(I)C \xrightarrow{\tau} B)
$$
由于 $\varphi : k[x_1, \ldots, x_n] \to C$ 为 \'etale，可以找到非零非常数多项式
$g \in k[T_1, \ldots, T_{n + 1}]$，使得
$$
g(\varphi(x_1), \ldots, \varphi(x_n), z) = 0
\quad\text{于}\quad
C
$$
为说明这一点，例如可用如下事实：
$C \otimes_{\varphi, k[x_1, \ldots, x_n]} k(x_1, \ldots, x_n)$
是 $k(x_1, \ldots, x_n)$ 的有限可分域扩张的有限积
（见《代数》引理 \ref{algebra-lemma-etale-over-field}），因此 $z$ 满足一个
以 $k(x_1, \ldots, x_n)$ 中元素为系数的首一多项式。清除分母即得 $g$。

\medskip\noindent
$g$ 的存在性与《代数》引理 \ref{algebra-lemma-helper-polynomial} 给出整数
$e_1, e_2, \ldots, e_n \geq 1$，使得 $z$ 在子环 $C'$ 上整，而该子环属于 $C$ 并由
$t_1 = \varphi(x_1) + z^{pe_1}, \ldots, t_n = \varphi(x_n) + z^{pe_n}$ 生成。
当然，元素 $\varphi(x_1), \ldots, \varphi(x_n)$ 也在 $C'$ 上整，元素
$y_1, \ldots, y_s$ 亦然。最后，由 $z$ 的选取，元素
$y_{s + 1}$ 也在 $C'$ 上整。

\medskip\noindent
考虑环映射
$$
\varphi' : k[x_1, \ldots, x_n] \longrightarrow C, \quad
x_i \longmapsto t_i
$$
其像为 $C'$。由于
$\text{d}(\varphi(x_i)) = \text{d}(t_i) = \text{d}(\varphi'(x_i))$
在 $\Omega_{C/k}$ 中成立（这里用到 $k$ 的特征为 $p > 0$），所以
$\varphi'$ 为 \'etale，因为 $\varphi$ 为 \'etale；见《代数》引理
\ref{algebra-lemma-characterize-etale-over-polynomial-ring}。注意，
$\varphi'(x_i) - \varphi(x_i) = t_i - \varphi(x_i) = z^{pe_i}$ 属于核
$J$，它是映射 $C \to C/\varphi(I)C \to B$ 的核；这是因为选取的 $z$ 是 $J$ 的元素。
因此，对 $f \in I$，元素
$$
\varphi'(f) =
f(t_1, \ldots, t_n) =
f(\varphi(x_1) + z^{pe_1}, \ldots, \varphi(x_n) + z^{pe_n}) =
\varphi(f) + \text{element of }(z)
$$
也属于 $J$。换言之，$\varphi'(I)C \subset J$，并得到满射
$$
\tau' : C/\varphi'(I)C \longrightarrow C/J \cong B
$$
（在 $A$ 上 \'etale 的代数之间）。最后，代数 $C$ 由元素
$\varphi(x_1), \ldots, \varphi(x_n), y_1, \ldots, y_r$ 生成（在
$C' = \Im(\varphi')$ 上），而
$\varphi(x_1), \ldots, \varphi(x_n), y_1, \ldots, y_{s + 1}$
在 $C' = \Im(\varphi')$ 上整。因此
$NF(C, \varphi', \tau') < r - s = NF(C, \varphi, \tau)$。证明完成。
\end{proof}

\begin{lemma}
\label{pione-lemma-affine-etale-over-affine-space}
设 $k$ 是特征为 $p > 0$ 的域。设 $X \to \mathbf{A}^n_k$ 是 \'etale 态射，
且 $X$ 仿射。则存在有限 \'etale 态射 $X \to \mathbf{A}^n_k$。
\end{lemma}

\begin{proof}
记 $X = \Spec(C)$。令 $A = 0$，并记 $I = k[x_1, \ldots, x_n]$。
由假设，存在某个 \'etale $k$-代数映射
$\varphi : k[x_1, \ldots, x_n] \to C$。以 $\tau : C/\varphi(I)C \to 0$
表示唯一满射。可以选取 $\varphi$ 和 $\tau$，使 $N(C, \varphi, \tau)$ 最小。
由引理 \ref{pione-lemma-lower-invariant}，得到 $N(C, \varphi, \tau) = 0$。
因此 $\varphi$ 有限 \'etale。
\end{proof}

\begin{lemma}
\label{pione-lemma-dominate-affine-space}
设 $k$ 是特征为 $p > 0$ 的域。设 $Z \subset \mathbf{A}^n_k$ 是闭子概形。
设 $Y \to Z$ 有限 \'etale。则存在有限 \'etale 态射
$f : U \to \mathbf{A}^n_k$，使得存在开且闭的浸入
$Y \to f^{-1}(Z)$（在 $Z$ 上）。
\end{lemma}

\begin{proof}
把问题转化为代数。记
$\mathbf{A}^n_k = \Spec(k[x_1, \ldots, x_n])$.
则 $Z = \Spec(A)$，其中 $A = k[x_1, \ldots, x_n]/I$，这里
$I \subset k[x_1, \ldots, x_n]$ 是某个理想。记 $Y = \Spec(B)$，
于是 $Y \to Z$ 对应于有限 \'etale $k$-代数映射 $A \to B$。

\medskip\noindent
由《代数》引理 \ref{algebra-lemma-lift-etale}，存在 \'etale 环映射
$$
\varphi : k[x_1, \ldots, x_n] \to C
$$
以及满的 $A$-代数映射 $\tau : C/\varphi(I)C \to B$。
（甚至可以选取 $C, \varphi, \tau$ 使 $\tau$ 为同构，但这里不用这一点。）
可以选取 $\varphi$ 和 $\tau$，使 $N(C, \varphi, \tau)$ 最小。
由引理 \ref{pione-lemma-lower-invariant}，得到 $N(C, \varphi, \tau) = 0$。
因此 $\varphi$ 有限 \'etale。

\medskip\noindent
令 $f : U = \Spec(C) \to \mathbf{A}^n_k$ 为对应于 $\varphi$ 的有限 \'etale
态射。态射 $Y \to f^{-1}(Z) = \Spec(C/\varphi(I)C)$（由 $\tau$ 诱导）
是闭浸入，因为 $\tau$ 满；
它也是开的，因为由《态射》引理
\ref{morphisms-lemma-etale-permanence}，它是 \'etale 态射。证明完成。
\end{proof}

\noindent
下面是主要结果。

\begin{proposition}
\label{pione-proposition-injective}
设 $p$ 是素数。设 $i : Z \to X$ 是 $\mathbf{F}_p$ 上连通仿射概形的闭浸入。
对任意几何点 $\overline{z}$（取自 $Z$），映射
$$
\pi_1(Z, \overline{z}) \to \pi_1(X, \overline{z})
$$
是单射。
\end{proposition}

\begin{proof}
设 $Y \to Z$ 是有限 \'etale 态射。只须构造有限 \'etale 态射
$f : U \to X$，使 $Y$ 同构于 $f^{-1}(Z)$ 的一个开且闭子概形；见引理
\ref{pione-lemma-functoriality-galois-injective}。记 $Y = \Spec(A)$ 且
$X = \Spec(R)$，于是闭浸入 $Y \to X$ 由满射 $R \to A$ 给出。
可把 $A = \colim A_i$ 写成它的有限型 $\mathbf{F}_p$-子代数的滤过余极限。
由引理 \ref{pione-lemma-limit}，可以找到一个 $i$ 和有限 \'etale 态射
$Y_i \to Z_i = \Spec(A_i)$，使得 $Y = Z \times_{Z_i} Y_i$。

\medskip\noindent
选取满射 $\mathbf{F}_p[x_1, \ldots, x_n] \to A_i$。它确定闭浸入
$$
Z_i = \Spec(A_i)
\longrightarrow
X_i = \mathbf{A}^n_{\mathbf{F}_p} = \Spec(\mathbf{F}_p[x_1, \ldots, x_n])
$$
由多项式代数的泛性质以及 $R \to A$ 为满射这一事实，可以找到交换图
$$
\xymatrix{
\mathbf{F}_p[x_1, \ldots, x_n] \ar[r] \ar[d] &
A_i \ar[d] \\
R \ar[r] &
A
}
$$
（在 $\mathbf{F}_p$-代数范畴中）。因此有交换图
$$
\xymatrix{
Y_i \ar[r] &
Z_i \ar[r] &
X_i \\
Y \ar[u] \ar[r] &
Z \ar[u] \ar[r] &
X \ar[u]
}
$$
其右方形为 Cartesian。显然，若能找到有限 \'etale 的
$f_i : U_i \to X_i$，使 $Y_i$ 同构于 $f_i^{-1}(Z_i)$ 的一个开且闭子概形，
则基变换 $f : U \to X$（即 $f_i$ 沿 $X \to X_i$ 的基变换）就是本问题的解。
因此，应用引理 \ref{pione-lemma-dominate-affine-space} 于
$Y_i \to Z_i \to X_i = \mathbf{A}^n_{\mathbf{F}_p}$ 即得结论。
\end{proof}
